% This LaTeX document needs to be compiled with XeLaTeX.
\documentclass[10pt]{article}
\usepackage[utf8]{inputenc}
\usepackage{amsmath}
\usepackage{amsfonts}
\usepackage{amssymb}
\usepackage[version=4]{mhchem}
\usepackage{stmaryrd}
\usepackage{hyperref}
\hypersetup{colorlinks=true, linkcolor=blue, filecolor=magenta, urlcolor=cyan,}
\urlstyle{same}
\usepackage{graphicx}
\usepackage[export]{adjustbox}
\graphicspath{ {./images/} }
\usepackage{caption}
\usepackage{bbold}
\usepackage{multirow}
\usepackage{fvextra, csquotes}
% \usepackage[fallback]{xeCJK}
\usepackage{polyglossia}
\usepackage{fontspec}
\usepackage{newunicodechar}
\IfFontExistsTF{Noto Serif CJK TC}
{\setCJKmainfont{Noto Serif CJK TC}}
{\IfFontExistsTF{STSong}
  {\setCJKmainfont{STSong}}
  {\IfFontExistsTF{Droid Sans Fallback}
    {\setCJKmainfont{Droid Sans Fallback}}
    {\setCJKmainfont{SimSun}}
}}

\setmainlanguage{english}
\IfFontExistsTF{CMU Serif}
{\setmainfont{CMU Serif}}
{\IfFontExistsTF{DejaVu Sans}
  {\setmainfont{DejaVu Sans}}
  {\setmainfont{Georgia}}
}

 Finite-Sample Properties of OLS }
%New command to display footnote whose markers will always be hidden
\let\svthefootnote\thefootnote
\newcommand\blfootnotetext[1]{%
  \let\thefootnote\relax\footnote{#1}%
  \addtocounter{footnote}{-1}%
  \let\thefootnote\svthefootnote%
}
%Overriding the \footnotetext command to hide the marker if its value is `0`
\let\svfootnotetext\footnotetext
\renewcommand\footnotetext[2][?]{%
  \if\relax#1\relax%
    \ifnum\value{footnote}=0\blfootnotetext{#2}\else\svfootnotetext{#2}\fi%
  \else%
    \if?#1\ifnum\value{footnote}=0\blfootnotetext{#2}\else\svfootnotetext{#2}\fi%
    \else\svfootnotetext[#1]{#2}\fi%
  \fi
}
\newunicodechar{×}{\ifmmode\times\else{$\times$}\fi}
\newunicodechar{⋮}{\ifmmode\vdots\else{$\vdots$}\fi}
\newunicodechar{⇔}{\ifmmode\Leftrightarrow\else{$\Leftrightarrow$}\fi}
\newunicodechar{¥}{\ifmmode\text{\yen}\else\yen\fi}
\title{Chapter 2: Large-Sample Theory}

\author{}
\date{}

\begin{document}
\maketitle

\section*{CHAPTER 2 Large-Sample Theory}
\begin{abstract}
In the previous chapter, we derived the exact- or finite-sample distribution of the OLS estimator and its associated test statistics. However, not very often in economics are the assumptions of the exact distribution satisfied. The finite-sample theory breaks down if one of the following three assumptions is violated: (1) the exogeneity of regressors, (2) the normality of the error term, and (3) the linearity of the regression equation. This chapter develops an alternative approach, retaining only the third assumption. The approach, called asymptotic or large-sample theory, derives an approximation to the distribution of the estimator and its associated statistics assuming that the sample size is sufficiently large.\\
Rather than making assumptions on the sample of a given size, large-sample theory makes assumptions on the stochastic process that generates the sample. The first two sections of this chapter provide the necessary language for describing stochastic processes.\\
The concepts introduced in this chapter are essential for rational expectations econometrics, as illustrated in Fama's classic paper on the Fisher Hypothesis that the real interest rate is constant. The hypothesis has a very strong policy implication: monetary and fiscal policies cannot influence aggregate demand through the real interest rate. Very surprisingly, one cannot reject the hypothesis for the United States (at least if the sample period ends in the early 1970s).
\end{abstract}

\subsection*{2.1 Review of Limit Theorems for Sequences of Random Variables}
The material of this section concerns the limiting behavior of a sequence of random variables, $\left(z_{1}, z_{2}, \ldots\right)$. Since the material may already be familiar to you, we present it rather formally, in a series of definitions and theorems. An authoritative source is Rao (1973, Chapter 2c) which gives proofs of all the theorems included in this section. In this section and the rest of this book, a sequence ( $z_{1}, z_{2}, \ldots$ ) will be denoted by $\left\{z_{n}\right\}$.

\section*{Various Modes of Convergence}
Convergence in Probability\\
A sequence of random scalars $\left\{z_{n}\right\}$ converges in probability to a constant (nonrandom) $\alpha$ if, for any $\varepsilon>0$,


\begin{equation*}
\lim _{n \rightarrow \infty} \operatorname{Prob}\left(\left|z_{n}-\alpha\right|>\varepsilon\right)=0 . \tag{2.1.1}
\end{equation*}


The constant $\alpha$ is called the probability limit of $z_{n}$ and is written as "plim ${ }_{n \rightarrow \infty} z_{n}= \alpha$ " or " $z_{n} \rightarrow_{\mathrm{p}} \alpha$ ". Evidently,

$$
" z_{n} \underset{\mathrm{p}}{\rightarrow} \alpha " \text { is the same as } " z_{n}-\alpha \underset{\mathrm{p}}{\rightarrow} 0 \text {." }
$$

This definition of convergence in probability is extended to a sequence of random vectors or random matrices (by viewing a matrix as a vector whose elements have been rearranged) by requiring element-by-element convergence in probability. That is, a sequence of $K$-dimensional random vectors $\left\{\mathbf{z}_{n}\right\}$ converges in probability to a $K$-dimensional vector of constants $\alpha$ if, for any $\varepsilon>0$,


\begin{equation*}
\lim _{n \rightarrow \infty} \operatorname{Prob}\left(\left|z_{n k}-\alpha_{k}\right|>\varepsilon\right)=0 \text { for all } k(=1,2, \ldots, K), \tag{2.1.2}
\end{equation*}


where $z_{n k}$ is the $k$-th element of $\mathbf{z}_{n}$ and $\alpha_{k}$ the $k$-th element of $\alpha$.

\section*{Almost Sure Convergence}
A sequence of random scalars $\left\{z_{n}\right\}$ converges almost surely to a constant $\alpha$ if


\begin{equation*}
\operatorname{Prob}\left(\lim _{n \rightarrow \infty} z_{n}=\alpha\right)=1 . \tag{2.1.3}
\end{equation*}


We write this as " $z_{n} \rightarrow_{\text {a.s. }} \alpha$." The extension to random vectors is analogous to that for convergence in probability. As will be mentioned below, this concept of convergence is stronger than convergence in probability; that is, if a sequence converges almost surely, then it converges in probability. The concept involved in (2.1.3) is harder to grasp because the probability is about an event concerning an infinite sequence ( $z_{1}, z_{2}, \ldots$ ). For our purposes, however, all that matters is that almost sure convergence is stronger than convergence in probability. If we can show that a sequence converges almost surely, that is one way to prove the sequence converges in probability.

Convergence in Mean Square\\
A sequence of random scalars $\left\{z_{n}\right\}$ converges in mean square (or in quadratic mean) to $\alpha$ (written as " $z_{n} \rightarrow_{\text {m.s. }} \alpha$ ") if


\begin{equation*}
\lim _{n \rightarrow \infty} \mathrm{E}\left[\left(z_{n}-\alpha\right)^{2}\right]=0 \tag{2.1.4}
\end{equation*}


The extension to random vectors is analogous to that for convergence in probability: $\mathbf{z}_{n} \rightarrow_{\text {m.s. }} \boldsymbol{\alpha}$ if each element of $\mathbf{z}_{n}$ converges in mean square to the corresponding component of $\boldsymbol{\alpha}$.

\section*{Convergence to a Random Variable}
In these definitions of convergence, the limit is a constant (i.e., a real number). The limit can be a random variable. We say that a sequence of $K$-dimensional random variables $\left\{\mathbf{z}_{n}\right\}$ converges to a $K$-dimensional random variable $\mathbf{z}$ and write $\mathbf{z}_{n} \rightarrow_{\mathrm{p}} \mathbf{z}$ if $\left\{\mathbf{z}_{n}-\mathbf{z}\right\}$ converges to $\mathbf{0}$ :


\begin{equation*}
" \mathbf{z}_{n} \underset{\mathrm{p}}{\rightarrow} \mathbf{z} " \text { is the same as } " \mathbf{z}_{n}-\mathbf{z} \underset{\mathrm{p}}{\rightarrow} \mathbf{0} \text {." } \tag{2.1.5a}
\end{equation*}


Similarly,


\begin{gather*}
" \mathbf{z}_{n} \underset{\text { a.s. }}{\rightarrow} \mathbf{z} " \text { is the same as } " \mathbf{z}_{n}-\mathbf{z} \underset{\text { a.s. }}{\rightarrow} \mathbf{0}, "  \tag{2.1.5b}\\
" \mathbf{z}_{n} \underset{\text { m.s. }}{\rightarrow} \mathbf{z} " \text { is the same as } " \mathbf{z}_{n}-\mathbf{z} \underset{\text { m.s. }}{\rightarrow} \mathbf{0} . " \tag{2.1.5c}
\end{gather*}


\section*{Convergence in Distribution}
Let $\left\{z_{n}\right\}$ be a sequence of random scalars and $F_{n}$ be the cumulative distribution function (c.d.f.) of $z_{n}$. We say that $\left\{z_{n}\right\}$ converges in distribution to a random scalar $z$ if the c.d.f. $F_{n}$ of $z_{n}$ converges to the c.d.f. $F$ of $z$ at every continuity point of $F^{1}$. We write " $z_{n} \rightarrow_{\mathrm{d}} z$ " or " $z_{n} \rightarrow_{\mathrm{L}} z$ " and call $F$ the asymptotic (or limit or limiting) distribution of $z_{n}$. Sometimes we write " $z_{n} \rightarrow_{\mathrm{d}} F$ ", when the distribution $F$ is well-known. For example, " $z_{n} \rightarrow_{\mathrm{d}} N(0,1)$ " should read " $z_{n} \rightarrow_{\mathrm{d}} z$ and the distribution of $z$ is $N(0,1)$ (normal distribution with mean 0 and variance 1)." It can be shown from the definition that convergence in probability is stronger than convergence in distribution, that is,


\begin{equation*}
" \mathbf{z}_{n} \underset{\mathrm{p}}{\rightarrow} \mathbf{z} " \Rightarrow " \mathbf{z}_{n} \underset{\mathrm{~d}}{\rightarrow} \mathbf{z} . " \tag{2.1.6}
\end{equation*}


\footnotetext{${ }^{1}$ Do not be concerned about the qualifier "at every continuity point of $F$." For the most part, except possibly for the chapters on the discrete or limited dependent variable, the relevant distribution is continuous, and the distribution function is continuous at all points.
}A special case of convergence in distribution is that $z$ is a constant (a trivial random variable).

The extension to a sequence of random vectors is immediate: $\mathbf{z}_{n} \rightarrow_{\mathrm{d}} \mathbf{Z}$ if the joint c.d.f. $F_{n}$ of the random vector $\mathbf{z}_{n}$ converges to the joint c.d.f. $F$ of $\mathbf{z}$ at every continuity point of $F$. Note, however, that, unlike the other concepts of convergence, for convergence in distribution, element-by-element convergence does not necessarily mean convergence for the vector sequence. That is, "each element of $\mathbf{z}_{n} \rightarrow_{\mathrm{d}}$ corresponding element of $\mathbf{z}$ " does not necessarily imply " $\mathbf{z}_{n} \rightarrow_{\mathrm{d}} \mathbf{z}$." A common way to establish the connection between scalar convergence and vector convergence in distribution is

Multivariate Convergence in Distribution Theorem: (stated in Rao, 1973, p. 128) Let $\left\{\mathbf{z}_{n}\right\}$ be a sequence of $K$-dimensional random vectors. Then:\\
" $\mathbf{z}_{n} \xrightarrow[\mathrm{~d}]{ } \mathbf{z}^{\prime \prime} \Leftrightarrow$ " $\lambda^{\prime} \mathbf{z}_{n} \xrightarrow[\mathrm{~d}]{ } \lambda^{\prime} \mathbf{z}$ for any $K$-dimensional vector of real numbers."

\section*{Convergence in Distribution vs. Convergence in Moments}
It is worth emphasizing that the moments of the limit distribution of $z_{n}$ are not necessarily equal to the limits of the moments of $z_{n}$. For example, " $z_{n} \rightarrow_{\mathrm{d}} z$ " does not necessarily imply " $\lim _{n \rightarrow \infty} \mathrm{E}\left(z_{n}\right)=\mathrm{E}(z)$." However,

Lemma 2.1 (convergence in distribution and in moments): Let $\alpha_{s n}$ be the $s$-th moment of $z_{n}$ and $\lim _{n \rightarrow \infty} \alpha_{s n}=\alpha_{s}$ where $\alpha_{s}$ is finite (i.e., a real number). Then:

$$
" z_{n} \underset{\mathrm{~d}}{\rightarrow} z " \Rightarrow " \alpha_{s} \text { is the } s \text {-th moment of } z . "
$$

Thus, for example, if the variance of a sequence of random variables converging in distribution converges to some finite number, then that number is the variance of the limiting distribution.

\section*{Relation among Modes of Convergence}
Some modes of convergence are weaker than others. The following theorem establishes the relationship between the four modes of convergence.

\section*{Lemma 2.2 (relationship among the four modes of convergence):}
(a) " $\mathbf{z}_{n} \rightarrow_{\text {m.s. }} \boldsymbol{\alpha} " \Rightarrow " \mathbf{z}_{n} \rightarrow_{\mathrm{p}} \boldsymbol{\alpha} . "$ So $" \mathbf{z}_{n} \rightarrow_{\text {m.s. }} \mathbf{z} " \Rightarrow " \mathbf{z}_{n} \rightarrow_{\mathrm{p}} \mathbf{z} . "$\\
(b) " $\mathbf{z}_{n} \rightarrow_{\text {a.s. }} \boldsymbol{\alpha} " \Rightarrow " \mathbf{z}_{n} \rightarrow_{\mathrm{p}} \boldsymbol{\alpha} . "$ So $" \mathbf{z}_{n} \rightarrow_{\text {a.s. }} \mathbf{z} " \Rightarrow " \mathbf{z}_{n} \rightarrow_{\mathrm{p}} \mathbf{z} . "$\\
(c) " $\mathbf{z}_{n} \rightarrow_{\mathrm{p}} \boldsymbol{\alpha}$ " $\Leftrightarrow$ " $\mathbf{z}_{n} \rightarrow_{\mathrm{d}} \boldsymbol{\alpha}$." (That is, if the limiting random variable is a constant [a trivial random variable], convergence in distribution is the same as convergence in probability.)

\section*{Three Useful Results}
Having defined the modes of convergence, we can state three results essential for developing large-sample theory.

Lemma 2.3 (preservation of convergence for continuous transformation): Suppose $\mathbf{a}(\cdot)$ is a vector-valued continuous function that does not depend on $n$.\\
(a) " $\mathbf{z}_{n} \rightarrow_{\mathrm{p}} \boldsymbol{\alpha} " \Rightarrow " \mathbf{a}\left(\mathbf{z}_{n}\right) \rightarrow_{\mathrm{p}} \mathbf{a}(\boldsymbol{\alpha})$." Stated differently,

$$
\operatorname{plim}_{n \rightarrow \infty} \mathbf{a}\left(\mathbf{z}_{n}\right)=\mathbf{a}\left(\operatorname{plim}_{n \rightarrow \infty} \mathbf{z}_{n}\right)
$$

provided the plim exists.\\
(b) " $\mathbf{z}_{n} \rightarrow_{\mathrm{d}} \mathbf{z} " \Rightarrow " \mathbf{a}\left(\mathbf{z}_{n}\right) \rightarrow_{\mathrm{d}} \mathbf{a}(\mathbf{z})$."

An immediate implication of Lemma 2.3(a) is that the usual arithmetic operations preserve convergence in probability. For example:

$$
\begin{aligned}
& " x_{n} \underset{\mathrm{p}}{\rightarrow} \beta, y_{n} \underset{\mathrm{p}}{\rightarrow} \gamma " \Rightarrow " x_{n}+y_{n} \rightarrow \beta+\gamma " \\
& " x_{n} \underset{\mathrm{p}}{\rightarrow} \beta, y_{n} \underset{\mathrm{p}}{\rightarrow} \gamma " \Rightarrow " x_{n} y_{n} \underset{\mathrm{p}}{\rightarrow} \beta \gamma . " \\
& " x_{n} \underset{\mathrm{p}}{\rightarrow} \beta, y_{n} \underset{\mathrm{p}}{\rightarrow} \gamma " \Rightarrow " x_{n} / y_{n} \underset{\mathrm{p}}{\rightarrow} \beta / \gamma, \text { " provided that } \gamma \neq 0 . \\
& \quad " \mathbf{Y}_{n} \underset{\mathrm{p}}{\rightarrow} \boldsymbol{\Gamma} " \Rightarrow " \mathbf{Y}_{n}^{-1} \underset{\mathrm{p}}{\rightarrow} \boldsymbol{\Gamma}^{-1} \text { ", provided that } \boldsymbol{\Gamma} \text { is invertible. }
\end{aligned}
$$

The next result about combinations of convergence in probability and in distribution will be used repeatedly to derive the asymptotic distribution of the estimator.

\section*{Lemma 2.4:}
(a) " $\mathbf{x}_{n} \rightarrow_{\mathrm{d}} \mathbf{x}, \mathbf{y}_{n} \rightarrow_{\mathrm{p}} \boldsymbol{\alpha} " \Rightarrow " \mathbf{x}_{n}+\mathbf{y}_{n} \rightarrow_{\mathrm{d}} \mathbf{x}+\boldsymbol{\alpha}$."\\
(b) " $\mathbf{x}_{n} \rightarrow_{\mathrm{d}} \mathbf{x}, \mathbf{y}_{n} \rightarrow_{\mathrm{p}} \mathbf{0} " \Rightarrow " \mathbf{y}_{n}^{\prime} \mathbf{x}_{n} \rightarrow_{\mathrm{p}} \mathbf{0}$."\\
(c) " $\mathbf{x}_{n} \rightarrow_{\mathrm{d}} \mathbf{x}, \mathbf{A}_{n} \rightarrow_{\mathrm{p}} \mathbf{A}$ " $\Rightarrow$ " $\mathbf{A}_{n} \mathbf{x}_{n} \rightarrow_{\mathrm{d}} \mathbf{A x}$," provided that $\mathbf{A}_{n}$ and $\mathbf{x}_{n}$ are conformable. In particular, if $\mathbf{x} \sim N(\mathbf{0}, \mathbf{\Sigma})$, then $\mathbf{A}_{n} \mathbf{x}_{n} \rightarrow_{\mathrm{d}} N\left(\mathbf{0}, \mathbf{A} \boldsymbol{\Sigma} \mathbf{A}^{\prime}\right)$.\\
(d) " $\mathbf{x}_{n} \rightarrow_{\mathrm{d}} \mathbf{x}, \mathbf{A}_{n} \rightarrow_{\mathrm{p}} \mathbf{A}$ " $\Rightarrow$ " $\mathbf{x}_{n}^{\prime} \mathbf{A}_{n}^{-1} \mathbf{x}_{n} \rightarrow_{\mathrm{d}} \mathbf{x}^{\prime} \mathbf{A}^{-1} \mathbf{x}$," provided that $\mathbf{A}_{n}$ and $\mathbf{x}_{n}$ are conformable and $\mathbf{A}$ is nonsingular.

Parts (a) and (c) are sometimes called Slutzky's Theorem. By setting $\boldsymbol{\alpha}=\mathbf{0}$, part (a) implies:


\begin{equation*}
" \mathbf{x}_{n} \underset{\mathrm{~d}}{\rightarrow} \mathbf{x}, \mathbf{y}_{n} \underset{\mathrm{p}}{\rightarrow} \mathbf{0} " \Rightarrow " \mathbf{x}_{n}+\mathbf{y}_{n} \underset{\mathrm{~d}}{\rightarrow} \mathbf{x} . " \tag{2.1.7}
\end{equation*}


That is, if $\mathbf{z}_{n}=\mathbf{x}_{n}+\mathbf{y}_{n}$ and $\mathbf{y}_{n} \rightarrow_{\mathrm{p}} \mathbf{0}$ (i.e., if $\mathbf{z}_{n}-\mathbf{x}_{n} \rightarrow_{\mathrm{p}} \mathbf{0}$ ), then the asymptotic distribution of $\mathbf{z}_{n}$ is the same as that of $\mathbf{x}_{n}$. When $\mathbf{z}_{n}-\mathbf{x}_{n} \rightarrow_{\mathrm{p}} \mathbf{0}$, we sometimes (but not always) say that the two sequences are asymptotically equivalent and write it as

$$
\text { " } \mathbf{z}_{n} \underset{\mathrm{a}}{\sim} \mathbf{x}_{n} \text { " or " } \mathbf{z}_{n}=\mathbf{x}_{n}+o_{p} \text {," }
$$

where $o_{p}$ is some suitable random variable ( $\mathbf{y}_{n}$ here) that converges to zero in probability.

A standard trick in deriving the asymptotic distribution of a sequence of random variables is to find an asymptotically equivalent sequence whose asymptotic distribution is easier to derive. In particular, by replacing $\mathbf{y}_{n}$ by $\mathbf{y}_{n}-\boldsymbol{\alpha}$ in part (b) of the lemma, we obtain


\begin{equation*}
" \mathbf{x}_{n} \underset{\mathrm{~d}}{\rightarrow} \mathbf{x}, \mathbf{y}_{n} \underset{\mathrm{p}}{\rightarrow} \boldsymbol{\alpha}^{\prime \prime} \Rightarrow " \mathbf{y}_{n}^{\prime} \mathbf{x}_{n} \underset{\mathrm{a}}{\sim} \boldsymbol{\alpha}^{\prime} \mathbf{x}_{n} " \text { or " } \mathbf{y}_{n}^{\prime} \mathbf{x}_{n}=\boldsymbol{\alpha}^{\prime} \mathbf{x}_{n}+o_{p} . " \tag{2.1.8}
\end{equation*}


The $o_{p}$ here is $\left(\mathbf{y}_{n}-\boldsymbol{\alpha}\right)^{\prime} \mathbf{x}_{n}$. Therefore, replacing $\mathbf{y}_{n}$ by its probability limit does not change the asymptotic distribution of $\mathbf{y}_{n}^{\prime} \mathbf{x}_{n}$, provided $\mathbf{x}_{n}$ converges in distribution to some random variable.

The third result will allow us to test nonlinear hypotheses given the asymptotic distribution of the estimator.

Lemma 2.5 (the "delta method"): Suppose $\left\{\mathbf{x}_{n}\right\}$ is a sequence of $K$-dimensional random vectors such that $\mathbf{x}_{n} \rightarrow_{\mathrm{p}} \boldsymbol{\beta}$ and

$$
\sqrt{n}\left(\mathbf{x}_{n}-\boldsymbol{\beta}\right) \underset{\mathrm{d}}{\rightarrow} \mathbf{z},
$$

and suppose $\mathbf{a}(\cdot): \mathbb{R}^{K} \rightarrow \mathbb{R}^{r}$ has continuous first derivatives with $\mathbf{A}(\boldsymbol{\beta})$ denoting the $r \times K$ matrix of first derivatives evaluated at $\beta$ :

$$
\underset{(r \times K)}{\mathbf{A}(\boldsymbol{\beta})} \equiv \frac{\partial \mathbf{a}(\boldsymbol{\beta})}{\partial \boldsymbol{\beta}^{\prime}} .
$$

Then

$$
\sqrt{n}\left[\mathbf{a}\left(\mathbf{x}_{n}\right)-\mathbf{a}(\boldsymbol{\beta})\right] \underset{\mathrm{d}}{\rightarrow} \mathbf{A}(\boldsymbol{\beta}) \mathbf{z}
$$

In particular:

$$
" \sqrt{n}\left(\mathbf{x}_{n}-\boldsymbol{\beta}\right) \underset{\mathrm{d}}{\rightarrow} N(\mathbf{0}, \boldsymbol{\Sigma}) " \Rightarrow " \sqrt{n}\left[\mathbf{a}\left(\mathbf{x}_{n}\right)-\mathbf{a}(\boldsymbol{\beta})\right] \underset{\mathrm{d}}{\rightarrow} N\left(\mathbf{0}, \mathbf{A}(\boldsymbol{\beta}) \boldsymbol{\Sigma} \mathbf{A}(\boldsymbol{\beta})^{\prime}\right) . "
$$

Proving this is a good way to learn how to use the results covered so far.

Proof. By the mean-value theorem from calculus (see Section 7.3 for a statement of the theorem), there exists a $K$-dimensional vector $\mathbf{y}_{n}$ between $\mathbf{x}_{n}$ and $\boldsymbol{\beta}$ such that

$$
\mathbf{a}\left(\mathbf{x}_{n}\right)-\mathbf{a}(\boldsymbol{\beta})=\underset{(r \times K)}{\mathbf{A}}\left(\mathbf{y}_{n}\right) \underset{(K \times 1)}{\left(\mathbf{x}_{n}-\boldsymbol{\beta}\right)}
$$

Multiplying both sides by $\sqrt{n}$, we obtain

$$
\sqrt{n}\left[\mathbf{a}\left(\mathbf{x}_{n}\right)-\mathbf{a}(\boldsymbol{\beta})\right]=\mathbf{A}\left(\mathbf{y}_{n}\right) \sqrt{n}\left(\mathbf{x}_{n}-\boldsymbol{\beta}\right)
$$

Since $\mathbf{y}_{n}$ is between $\mathbf{x}_{n}$ and $\boldsymbol{\beta}$ and since $\mathbf{x}_{n} \rightarrow_{\mathrm{p}} \boldsymbol{\beta}$, we know that $\mathbf{y}_{n} \rightarrow_{\mathrm{p}} \boldsymbol{\beta}$. Moreover, the first derivative $\mathbf{A}(\cdot)$ is continuous by assumption. So by Lemma 2.3(a),

$$
\mathbf{A}\left(\mathbf{y}_{n}\right) \underset{\mathrm{p}}{\rightarrow} \mathbf{A}(\boldsymbol{\beta})
$$

By Lemma 2.4(c), this and the hypothesis that $\sqrt{n}\left(\mathbf{x}_{n}-\boldsymbol{\beta}\right) \rightarrow_{\mathrm{d}} \mathbf{z}$ imply that

$$
\mathbf{A}\left(\mathbf{y}_{n}\right) \sqrt{n}\left(\mathbf{x}_{n}-\boldsymbol{\beta}\right)\left(=\sqrt{n}\left[\mathbf{a}\left(\mathbf{x}_{n}\right)-\mathbf{a}(\boldsymbol{\beta})\right]\right) \underset{\mathrm{d}}{\rightarrow} \mathbf{A}(\boldsymbol{\beta}) \mathbf{z}
$$

\section*{Viewing Estimators as Sequences of Random Variables}
Let $\hat{\boldsymbol{\theta}}_{n}$ be an estimator of a parameter vector $\boldsymbol{\theta}$ based on a sample of size $n$. The sequence $\left\{\hat{\boldsymbol{\theta}}_{n}\right\}$ is an example of a sequence of random variables, so the concepts introduced in this section for sequences of random variables are applicable to $\left\{\hat{\boldsymbol{\theta}}_{n}\right\}$. We say that an estimator $\hat{\boldsymbol{\theta}}_{n}$ is consistent for $\boldsymbol{\theta}$ if

$$
\operatorname{plim}_{n \rightarrow \infty} \hat{\boldsymbol{\theta}}_{n}=\boldsymbol{\theta} \quad \text { or } \quad \hat{\boldsymbol{\theta}}_{n} \underset{\mathrm{p}}{\rightarrow} \boldsymbol{\theta}
$$

The asymptotic bias of $\hat{\boldsymbol{\theta}}_{n}$ is defined as $\operatorname{plim}_{n \rightarrow \infty} \hat{\boldsymbol{\theta}}_{n}-\boldsymbol{\theta} .^{2}$ So if the estimator is consistent, its asymptotic bias is zero. A consistent estimator $\hat{\boldsymbol{\theta}}_{n}$ is asymptotically normal if

$$
\sqrt{n}\left(\hat{\boldsymbol{\theta}}_{n}-\boldsymbol{\theta}\right) \underset{\mathrm{d}}{\rightarrow} N(\mathbf{0}, \boldsymbol{\Sigma}) .
$$

Such an estimator is called $\sqrt{n}$-consistent. The acronym sometimes used for "consistent and asymptotically normal" is CAN. The variance matrix $\boldsymbol{\Sigma}$ is called the asymptotic variance and is denoted $\operatorname{Avar}\left(\hat{\boldsymbol{\theta}}_{n}\right)$. Some authors use the notation $\operatorname{Avar}\left(\hat{\boldsymbol{\theta}}_{n}\right)$ to mean $\boldsymbol{\Sigma} / n$ (which is zero in the limit). In this book, $\operatorname{Avar}\left(\hat{\boldsymbol{\theta}}_{n}\right)$ is the variance of the limiting distribution of $\sqrt{n}\left(\hat{\boldsymbol{\theta}}_{n}-\boldsymbol{\theta}\right)$.

\section*{Laws of Large Numbers and Central Limit Theorems}
For a sequence of random scalars $\left\{z_{i}\right\}$, the sample mean $\bar{z}_{n}$ is defined as

$$
\bar{z}_{n} \equiv \frac{1}{n} \sum_{i=1}^{n} z_{i} .
$$

Consider the sequence $\left\{\bar{z}_{n}\right\}$. Laws of large numbers (LLNs) concern conditions under which $\left\{\bar{z}_{n}\right\}$ converges either in probability or almost surely. An LLN is called strong if the convergence is almost surely and weak if the convergence is in probability. We can derive the following weak LLN easily from Part (a) of Lemma 2.2.

\section*{A Version of Chebychev's Weak LLN:}
$$
" \lim _{n \rightarrow \infty} \mathrm{E}\left(\bar{z}_{n}\right)=\mu, \lim _{n \rightarrow \infty} \operatorname{Var}\left(\bar{z}_{n}\right)=0 " \Rightarrow " \bar{z}_{n} \underset{\mathrm{p}}{\rightarrow \mu .} "
$$

This holds because, under the condition specified, it is easy to prove (see an analytical question) that $\bar{z}_{n} \rightarrow_{\text {m.s. }} \mu$. The following strong LLN assumes that $\left\{z_{i}\right\}$ is i.i.d. (independently and identically distributed), but the variance does not need to be finite.

Kolmogorov's Second Strong Law of Large Numbers: Let $\left\{z_{i}\right\}$ be i.i.d. with $\mathrm{E}\left(z_{i}\right)=\mu .^{3}$ Then $\bar{z}_{n} \rightarrow_{\text {a.s. }} \mu$.

\footnotetext{${ }^{2}$ Some authors use the term "asymptotic bias" differently. Amemiya (1985), for example, defines it to mean $\lim _{n \rightarrow \infty} \mathrm{E}\left(\hat{\boldsymbol{\theta}}_{n}\right)-\boldsymbol{\theta}$.\\
${ }^{3}$ So the mean exists and is finite (a real number). When a moment (e.g., the mean) is indicated, as here, then by implication the moment is assumed to exist and is finite.
}These LLNs extend readily to random vectors by requiring element-by-element convergence.

Central Limit Theorems (CLTs) are about the limiting behavior of the difference between $\bar{z}_{n}$ and $\mathrm{E}\left(\bar{z}_{n}\right)$ (which equals $\mathrm{E}\left(z_{i}\right)$ if $\left\{z_{i}\right\}$ is i.i.d.) blown up by $\sqrt{n}$. The only Central Limit Theorem we need for the case of i.i.d. sequences is:

Lindeberg-Levy CLT: Let $\left\{\mathbf{z}_{i}\right\}$ be i.i.d. with $\mathrm{E}\left(\mathbf{z}_{i}\right)=\boldsymbol{\mu}$ and $\operatorname{Var}\left(\mathbf{z}_{i}\right)=\mathbf{\Sigma}$. Then

$$
\sqrt{n}\left(\overline{\mathbf{z}}_{n}-\boldsymbol{\mu}\right)=\frac{1}{\sqrt{n}} \sum_{i=1}^{n}\left(\mathbf{z}_{i}-\boldsymbol{\mu}\right) \underset{\mathrm{d}}{\rightarrow} N(\mathbf{0}, \mathbf{\Sigma})
$$

This reads: a sequence of random vectors $\left\{\sqrt{n}\left(\overline{\mathbf{z}}_{n}-\boldsymbol{\mu}\right)\right\}$ converges in distribution to a random vector whose distribution is $N(\mathbf{0}, \boldsymbol{\Sigma})$. (Usually, the Lindeberg-Levy CLT is for a sequence of scalar random variables. The vector version displayed above is derived from the scalar version as follows. Let $\left\{\mathbf{z}_{i}\right\}$ be i.i.d. with $\mathrm{E}\left(\mathbf{z}_{i}\right)=\boldsymbol{\mu}$ and $\operatorname{Var}\left(\mathbf{z}_{i}\right)=\boldsymbol{\Sigma}$, and let $\boldsymbol{\lambda}$ be any vector of real numbers of the same dimension. Then $\left\{\boldsymbol{\lambda}^{\prime} \mathbf{z}_{n}\right\}$ is a sequence of scalar random variables with $\mathrm{E}\left(\boldsymbol{\lambda}^{\prime} \mathbf{z}_{n}\right)=\boldsymbol{\lambda}^{\prime} \boldsymbol{\mu}$ and $\operatorname{Var}\left(\boldsymbol{\lambda}^{\prime} \mathbf{z}_{n}\right)=\boldsymbol{\lambda}^{\prime} \boldsymbol{\Sigma} \boldsymbol{\lambda}$. The scalar version of Linderberg-Levy then implies that

$$
\sqrt{n}\left(\boldsymbol{\lambda}^{\prime} \overline{\mathbf{z}}_{n}-\boldsymbol{\lambda}^{\prime} \boldsymbol{\mu}\right)=\boldsymbol{\lambda}^{\prime} \sqrt{n}\left(\overline{\mathbf{z}}_{n}-\boldsymbol{\mu}\right) \underset{\mathrm{d}}{\rightarrow} N\left(0, \boldsymbol{\lambda}^{\prime} \boldsymbol{\Sigma} \boldsymbol{\lambda}\right)
$$

But this limit distribution is the distribution of $\boldsymbol{\lambda}^{\prime} \mathbf{x}$ where $\mathbf{x} \sim N(\mathbf{0}, \boldsymbol{\Sigma})$. So by the Multivariate Convergence in Distribution Theorem stated a few pages back, $\left\{\sqrt{n}\left(\overline{\mathbf{z}}_{n}-\boldsymbol{\mu}\right)\right\} \rightarrow_{\mathrm{d}} \mathbf{x}$, which is the claim of the vector version of Lindeberg-Levy.)

\section*{QUESTIONS FOR REVIEW}
\begin{enumerate}
  \item (Usual convergence vs. convergence in probability) A sequence of real numbers is a trivial example of a sequence of random variables. Is it true that " $\lim _{n \rightarrow \infty} z_{n}=\alpha$ " $\Rightarrow$ " $\operatorname{plim}_{n \rightarrow \infty} z_{n}=\alpha$ "? Hint: Look at the definition of plim. Since $\lim _{n \rightarrow \infty} z_{n}=\alpha,\left|z_{n}-\alpha\right|<\varepsilon$ for $n$ sufficiently large.
  \item (Alternative definition of convergence for vector sequences) Verify that the definition in the text of " $\mathbf{z}_{n} \rightarrow_{\text {m.s. }} \mathbf{z}$ " is equivalent to
\end{enumerate}

$$
\lim _{n \rightarrow \infty} \mathrm{E}\left[\left(\mathbf{z}_{n}-\mathbf{z}\right)^{\prime}\left(\mathbf{z}_{n}-\mathbf{z}\right)\right]=0
$$

Hint: $\mathrm{E}\left[\left(\mathbf{z}_{n}-\mathbf{z}\right)^{\prime}\left(\mathbf{z}_{n}-\mathbf{z}\right)\right]=\mathrm{E}\left[\left(z_{n 1}-z_{1}\right)^{2}\right]+\cdots+\mathrm{E}\left[\left(z_{n K}-z_{K}\right)^{2}\right]$, where $K$ is the dimension of $\mathbf{z}$. Similarly, verify that the definition in the text of\\
" $\mathbf{z}_{n} \rightarrow_{\mathrm{p}} \boldsymbol{\alpha}$ " is equivalent to

$$
\lim _{n \rightarrow \infty} \operatorname{Prob}\left(\left(\mathbf{z}_{n}-\boldsymbol{\alpha}\right)^{\prime}\left(\mathbf{z}_{n}-\boldsymbol{\alpha}\right)>\varepsilon\right)=0 \text { for any } \varepsilon>0 .
$$

\begin{enumerate}
  \setcounter{enumi}{2}
  \item Prove Lemma 2.4(c) from Lemma 2.4(a) and (b). Hint: $\mathbf{A}_{n} \mathbf{x}_{n}=\left(\mathbf{A}_{n}-\mathbf{A}\right) \mathbf{x}_{n}+ \mathbf{A} \mathbf{x}_{n}$. By (b), $\left(\mathbf{A}_{n}-\mathbf{A}\right) \mathbf{x}_{n} \rightarrow_{\mathrm{p}} \mathbf{0}$.
  \item Suppose $\sqrt{n}\left(\hat{\theta}_{n}-\theta\right) \rightarrow_{\mathrm{d}} N\left(0, \sigma^{2}\right)$. Does it follow that $\hat{\theta}_{n} \rightarrow_{\mathrm{p}} \theta$ ? Hint:
\end{enumerate}

$$
\hat{\theta}_{n}-\theta=\frac{1}{\sqrt{n}} \cdot \sqrt{n}\left(\hat{\theta}_{n}-\theta\right), \quad \operatorname{plim}_{n \rightarrow \infty} \frac{1}{\sqrt{n}}=0 .
$$

\begin{enumerate}
  \setcounter{enumi}{4}
  \item (Combine Delta method with Lindeberg-Levy) Let $\left\{z_{i}\right\}$ be a sequence of i.i.d. (independently and identically distributed) random variables with $\mathrm{E}\left(z_{i}\right)=\mu \neq$ 0 and $\operatorname{Var}\left(z_{i}\right)=\sigma^{2}$, and let $\bar{z}_{n}$ be the sample mean. Show that
\end{enumerate}

$$
\sqrt{n}\left(\frac{1}{\bar{z}_{n}}-\frac{1}{\mu}\right) \rightarrow N\left(0, \frac{\sigma^{2}}{\mu^{4}}\right) .
$$

Hint: In Lemma 2.5, set $\boldsymbol{\beta}=\mu, \mathbf{a}(\boldsymbol{\beta})=1 / \mu, \mathbf{x}_{n}=\bar{z}_{n}$.

\subsection*{2.2 Fundamental Concepts in Time-Series Analysis}
In this section, we define the very basic concepts in time-series analysis that will form an integral part of our language. The key concept is a stochastic process, which is just a fancy name for a sequence of random variables. If the index for the random variables is interpreted as representing time, the stochastic process is called a time series. If $\left\{z_{i}\right\}(i=1,2, \ldots)$ is a stochastic process, its realization or a sample path is an assignment to each $i$ of a possible value of $z_{i}$. So a realization of $\left\{z_{i}\right\}$ is a sequence of real numbers. We will frequently use the term time series to mean both the realization and the process of which it is a realization.

\section*{Need for Ergodic Stationarity}
The fundamental problem in time-series analysis is that we can observe the realization of the process only once. For example, the sample on the U.S. annual inflation rate for the period from 1946 to 1995 is a string of 50 particular numbers, which is just one possible outcome of the underlying stochastic process for the inflation rate; if history took a different course, we could have obtained a different sample.

If we could observe the history many times over, we could assemble many samples, each containing a possibly different string of 50 numbers. The mean inflation rate for, say, 1995 can then be estimated by taking the average of the 1995 inflation rate (the 50th element of the string) across those samples. The population mean thus estimated is sometimes called the ensemble mean. In the language of general equilibrium theory in economics, the ensemble mean is the average across all the possible different states of nature at any given calendar time.

Of course, it is not feasible to observe many different alternative histories. But if the distribution of the inflation rate remains unchanged (this property will be referred to as stationarity), the particular string of 50 numbers we do observe can be viewed as 50 different values from the same distribution. Furthermore, if the process is not too persistent (what's called ergodicity has this property), each element of the string will contain some information not available from the other elements, and, as shown below, the time average over the elements of the single string will be consistent for the ensemble mean.

\section*{Various Classes of Stochastic Processes ${ }^{\mathbf{4}}$}
\section*{Stationary Processes}
A stochastic process $\left\{\mathbf{z}_{i}\right\}(i=1,2, \ldots)$ is (strictly) stationary if, for any given finite integer $r$ and for any set of subscripts, $i_{1}, i_{2}, \ldots, i_{r}$, the joint distribution of $\left(\mathbf{z}_{i}, \mathbf{z}_{i_{1}}, \mathbf{z}_{i_{2}}, \ldots, \mathbf{z}_{i_{r}}\right)$ depends only on $i_{1}-i, i_{2}-i, i_{3}-i, \ldots, i_{r}-i$ but not on $i$. For example, the joint distribution of ( $\mathbf{z}_{1}, \mathbf{z}_{5}$ ) is the same as that of ( $\mathbf{z}_{12}, \mathbf{z}_{16}$ ). What matters for the distribution is the relative position in the sequence. In particular, the distribution of $\mathbf{z}_{i}$ does not depend on the absolute position, $i$, of $\mathbf{z}_{i}$, so the mean, variance, and other higher moments, if they exist, remain the same across $i$. The definition also implies that any transformation (function) of a stationary process is itself stationary, that is, if $\left\{\mathbf{z}_{i}\right\}$ is stationary, then $\left\{f\left(\mathbf{z}_{i}\right)\right\}$ is. ${ }^{5}$ For example, $\left\{\mathbf{z}_{i} \mathbf{z}_{i}^{\prime}\right\}$ is stationary if $\left\{\mathbf{z}_{i}\right\}$ is.

Example 2.1 (i.i.d. sequences): A sequence of independent and identically distributed random variables is a stationary process that exhibits no serial dependence.

\footnotetext{${ }^{4}$ Many of the concepts collected in this subsection can also be found in Section 4.7 of Davidson and MacKinnon (1993).\\
${ }^{5}$ The function $f(\cdot)$ needs to be "measurable" so that $f\left(\mathbf{z}_{i}\right)$ is a well-defined random variable. Any continuous function is measurable. In what follows, we won't bother to add the qualifier "measurable" when a function $f$ of a random variable is understood to be a random variable.
}Example 2.2 (constant series): Draw $z_{1}$ from some distribution and then set $z_{i}=z_{1}(i=2,3, \ldots)$. So the value of the process is frozen at the initial date. The process $\left\{z_{i}\right\}$ thus created is a stationary process that exhibits maximum serial dependence.

Evidently, if a vector process $\left\{\mathbf{z}_{i}\right\}$ is stationary, then each element of the vector forms a univariate stationary process. The converse, however, is not true.

Example 2.3 (element-wise vs. joint stationarity): Let $\left\{\varepsilon_{i}\right\}(i=1,2, \ldots)$ be a scalar i.i.d. process. Create a two-dimensional process $\left\{\mathbf{z}_{i}\right\}$ from it by defining $z_{i 1}=\varepsilon_{i}$ and $z_{i 2}=\varepsilon_{1}$ The scalar process $\left\{z_{i 1}\right\}$ is stationary (this is the process of Example 2.1). The scalar process $\left\{z_{i 2}\right\}$, too, is stationary (the Example 2.2 process). The vector process $\left\{\mathbf{z}_{i}\right\}$, however, is not (jointly) stationary, because the (joint) distribution of $\mathbf{z}_{1}\left(=\left(\varepsilon_{1}, \varepsilon_{1}\right)^{\prime}\right)$ differs from that of $\mathbf{z}_{2}\left(=\left(\varepsilon_{2}, \varepsilon_{1}\right)^{\prime}\right)$.

Most aggregate time series such as GDP are not stationary because they exhibit time trends. A less obvious example of nonstationarity is international exchange rates, which are alleged to have increasing variance. But many time series with trend can be reduced to stationary processes. A process is called trend stationary if it is stationary after subtracting from it a (usually linear) function of time (which is the index $i$ ). If a process is not stationary but its first difference, $z_{i}-z_{i-1}$, is stationary, $\left\{z_{i}\right\}$ is called difference stationary. Trend-stationary processes and difference-stationary processes will be studied in Chapter 9.

\section*{Covariance Stationary Processes}
A stochastic process $\left\{\mathbf{z}_{i}\right\}$ is weakly (or covariance) stationary if:\\
(i) $\mathrm{E}\left(\mathbf{z}_{i}\right)$ does not depend on $i$, and\\
(ii) $\operatorname{Cov}\left(\mathbf{z}_{i}, \mathbf{z}_{i-j}\right)$ exists, is finite, and depends only on $j$ but not on $i$ (for example, $\operatorname{Cov}\left(\mathbf{z}_{1}, \mathbf{z}_{5}\right)$ equals $\operatorname{Cov}\left(\mathbf{z}_{12}, \mathbf{z}_{16}\right)$ ).

The relative, not absolute, position in the sequence matters for the mean and covariance of a covariance-stationary process. Evidently, if a sequence is (strictly) stationary and if the variance and covariances are finite, then the sequence is weakly stationary (hence the term "strict"). An example of a covariance-stationary but not strictly stationary process will be given in Example 2.4 below.

The $j$-th order autocovariance, denoted $\Gamma_{j}$, is defined as

$$
\Gamma_{j} \equiv \operatorname{Cov}\left(\mathbf{z}_{i}, \mathbf{z}_{i-j}\right) \quad(j=0,1,2, \ldots) .
$$

The term "auto" comes about because the two random variables are taken from the same process. $\boldsymbol{\Gamma}_{j}$ does not depend on $i$ because of covariance-stationarity. Also by covariance stationarity, $\boldsymbol{\Gamma}_{j}$ satisfies


\begin{equation*}
\boldsymbol{\Gamma}_{j}=\boldsymbol{\Gamma}_{-j}^{\prime} . \tag{2.2.1}
\end{equation*}


(Showing this is a review question below.) The 0 -th order autocovariance is the variance $\boldsymbol{\Gamma}_{0}=\operatorname{Var}\left(\mathbf{z}_{i}\right)$. The processes in Examples 2.1 and 2.2 are covariance stationary if the variance exists and is finite. For the process of Example 2.1, $\boldsymbol{\Gamma}_{0}$ is the variance of the distribution and $\boldsymbol{\Gamma}_{j}=\mathbf{0}$ for $j \geq 1$. For the process of Example 2.2, $\Gamma_{j}=\Gamma_{0}$.

For a scalar covariance stationary process $\left\{z_{i}\right\}$, the $j$-th order autocovariance is now a scalar. If $\gamma_{j}$ is this autocovariance, it satisfies


\begin{equation*}
\gamma_{j}=\gamma_{-j} . \tag{2.2.2}
\end{equation*}


Take a string of $n$ successive values, ( $z_{i}, z_{i+1}, \ldots, z_{i+n-1}$ ), from a scalar process. By covariance stationarity, its $n \times n$ variance-covariance matrix is the same as that of $\left(z_{1}, z_{2}, \ldots z_{n}\right)$ and is a band spectrum matrix:

$$
\operatorname{Var}\left(z_{i}, z_{i+1}, \ldots z_{i+n-1}\right)=\left[\begin{array}{ccccc}
\gamma_{0} & \gamma_{1} & \gamma_{2} & \ldots & \gamma_{n-1} \\
\gamma_{1} & \gamma_{0} & \gamma_{1} & \ldots & \gamma_{n-2} \\
\vdots & \ddots & \ddots & \ddots & \vdots \\
\gamma_{n-2} & \cdots & \gamma_{1} & \gamma_{0} & \gamma_{1} \\
\gamma_{n-1} & \cdots & \gamma_{2} & \gamma_{1} & \gamma_{0}
\end{array}\right] .
$$

This is called the autocovariance matrix of the process. The $\boldsymbol{j}$-th order autocorrelation coefficient, $\rho_{j}$, is defined as\\
$j$-th order autocorrelation coefficient $\equiv$


\begin{equation*}
\rho_{j} \equiv \frac{\gamma_{j}}{\gamma_{0}}=\frac{\operatorname{Cov}\left(z_{i}, z_{i-j}\right)}{\operatorname{Var}\left(z_{i}\right)} \quad(j=1,2, \ldots) . \tag{2.2.3}
\end{equation*}


For $j=0, \rho_{j}=1$. The plot of $\left\{\rho_{j}\right\}$ against $j=0,1,2, \ldots$ is called the correlogram.

\section*{White Noise Processes}
A very important class of weakly stationary processes is a white noise process, a process with zero mean and no serial correlation:\\
a covariance-stationary process $\left\{\mathbf{z}_{i}\right\}$ is white noise if

$$
\mathrm{E}\left(\mathbf{z}_{i}\right)=\mathbf{0} \text { and } \operatorname{Cov}\left(\mathbf{z}_{i}, \mathbf{z}_{i-j}\right)=\mathbf{0} \text { for } j \neq 0 .
$$

Clearly, an independently and identically distributed (i.i.d.) sequence with mean zero and finite variance is a special case of a white noise process. For this reason, it is called an independent white noise process.

Example 2.4 (a white noise process that is not strictly stationary ${ }^{6}$ ): Let $w$ be a random variable uniformly distributed in the interval $(0,2 \pi)$, and define

$$
z_{i}=\cos (i w) \quad(i=1,2, \ldots) .
$$

It can be shown that $\mathrm{E}\left(z_{i}\right)=0, \operatorname{Var}\left(z_{i}\right)=1 / 2$, and $\operatorname{Cov}\left(z_{i}, z_{j}\right)=0$ for $i \neq j$. So $\left\{z_{i}\right\}$ is white noise. However, clearly, it is not an independent white noise process. It is not even strictly stationary.

\section*{Ergodicity}
A stationary process $\left\{z_{i}\right\}$ is said to be ergodic if, for any two bounded functions $f: \mathbb{R}^{k} \rightarrow \mathbb{R}$ and $g: \mathbb{R}^{\ell} \rightarrow \mathbb{R}$,

$$
\begin{aligned}
& \lim _{n \rightarrow \infty}\left|\mathrm{E}\left[f\left(z_{i}, \ldots, z_{i+k}\right) g\left(z_{i+n}, \ldots, z_{i+n+\ell}\right)\right]\right| \\
& \quad=\left|\mathrm{E}\left[f\left(z_{i}, \ldots, z_{i+k}\right)\right]\right|\left|\mathrm{E}\left[g\left(z_{i+n}, \ldots, z_{i+n+\ell}\right)\right]\right| .
\end{aligned}
$$

Heuristically, a stationary process is ergodic if it is asymptotically independent, that is, if any two random variables positioned far apart in the sequence are almost independently distributed. A stationary process that is ergodic will be called ergodic stationary. Ergodic stationarity will be an integral ingredient in developing largesample theory because of the following property.

Ergodic Theorem: (See, e.g., Theorem 9.5.5 of Karlin and Taylor (1975).) Let $\left\{\mathbf{z}_{i}\right\}$ be a stationary and ergodic process with $\mathrm{E}\left(\mathbf{z}_{i}\right)=\mu .^{7}$ Then

\footnotetext{${ }^{6}$ Drawn from Example 7.8 of Anderson (1971, p. 379).\\
${ }^{7}$ So the mean is assumed to exist and is finite.
}
$$
\overline{\mathbf{z}}_{n} \equiv \frac{1}{n} \sum_{i=1}^{n} \mathbf{z}_{i} \underset{\text { a.s. }}{\rightarrow} \boldsymbol{\mu} .
$$

The Ergodic Theorem, therefore, is a substantial generalization of Kolmogorov's LLN. Serial dependence, which is ruled out by the i.i.d. assumption in Kolmogorov's LLN, is allowed in the Ergodic Theorem, provided that it disappears in the long run. Since, for any (measurable) function $f(\cdot),\left\{f\left(\mathbf{z}_{i}\right)\right\}$ is ergodic stationary whenever $\left\{\mathbf{z}_{i}\right\}$ is, this theorem implies that any moment of a stationary and ergodic process (if it exists and is finite) is consistently estimated by the sample moment. For example, suppose $\left\{\mathbf{z}_{i}\right\}$ is stationary and ergodic and $\mathrm{E}\left(\mathbf{z}_{i} \mathbf{z}_{i}^{\prime}\right)$ exists and is finite. Then, $\frac{1}{n} \sum_{i} \mathbf{z}_{i} \mathbf{z}_{i}^{\prime}$ is consistent for $\mathrm{E}\left(\mathbf{z}_{i} \mathbf{z}_{i}^{\prime}\right)$.

The simplest example of ergodic stationary processes is independent white noise processes. (White noise processes where independence is weakened to no serial correlation are not necessarily ergodic; Example 2.4 above is an example.) Another important example is the AR(1) process satisfying

$$
z_{i}=c+\rho z_{i-1}+\varepsilon_{i},|\rho|<1
$$

where $\left\{\varepsilon_{i}\right\}$ is independent white noise.

\section*{Martingales}
Let $x_{i}$ be an element of $\mathbf{z}_{i}$. The scalar process $\left\{x_{i}\right\}$ is called a martingale with respect to $\left\{\mathbf{z}_{i}\right\}$ if


\begin{equation*}
\mathrm{E}\left(x_{i} \mid \mathbf{z}_{i-1}, \mathbf{z}_{i-2}, \ldots, \mathbf{z}_{1}\right)=x_{i-1} \text { for } i \geq 2 .^{8} \tag{2.2.4}
\end{equation*}


The conditioning set ( $\mathbf{z}_{i-1}, \mathbf{z}_{i-2}, \ldots, \mathbf{z}_{1}$ ) is often called the information set at point (date) $i-1 .\left\{x_{i}\right\}$ is called simply a martingale if the information set is its own past values $\left(x_{i-1}, x_{i-2}, \ldots, x_{1}\right)$. If $\mathbf{z}_{i}$ includes $x_{i}$, then $\left\{x_{i}\right\}$ is a martingale, because

$$
\begin{aligned}
& \mathrm{E}\left(x_{i} \mid x_{i-1}, x_{i-2}, \ldots, x_{1}\right) \\
& =\mathrm{E}\left[\mathrm{E}\left(x_{i} \mid \mathbf{z}_{i-1}, \mathbf{z}_{i-2}, \ldots, \mathbf{z}_{1}\right) \mid x_{i-1}, x_{i-2}, \ldots, x_{1}\right]
\end{aligned}
$$

(Law of Iterated Expectations)

$$
=\mathrm{E}\left(x_{i-1} \mid x_{i-1}, x_{i-2}, \ldots, x_{1}\right)=x_{i-1} .
$$

\footnotetext{${ }^{8}$ If the process started in the infinite past so that $i$ runs from $-\infty$ to $+\infty$, the definition is $\mathrm{E}\left(x_{i} \mid\right. \left.\mathbf{z}_{i-1}, \mathbf{z}_{i-2}, \ldots\right)=x_{i-1}$, and the qualifier " $i \geq 2$ " is not needed. Whether the process started in the infinite past or in date $i=1$ is not important for large-sample theory to be developed below. What will matter is that the process starts before the sample period.
}A vector process $\left\{\mathbf{z}_{i}\right\}$ is called a martingale if


\begin{equation*}
\mathrm{E}\left(\mathbf{z}_{i} \mid \mathbf{z}_{i-1}, \ldots, \mathbf{z}_{1}\right)=\mathbf{z}_{i-1} \text { for } i \geq 2 . \tag{2.2.5}
\end{equation*}


Example 2.5 (Hall's Martingale Hypothesis): Let $\mathbf{z}_{i}$ be a vector containing a set of macroeconomic variables (such as the money supply or GDP) including aggregate consumption $c_{i}$ for period $i$. Hall's (1978) martingale hypothesis is that consumption is a martingale with respect to $\left\{\mathbf{z}_{i}\right\}$ :

$$
\mathrm{E}\left(c_{i} \mid \mathbf{z}_{i-1}, \mathbf{z}_{i-2}, \ldots, \mathbf{z}_{1}\right)=c_{i-1} .
$$

This formalizes the notion in consumption theory called "consumption smoothing": the consumer, wishing to avoid fluctuations in the standard of living, adjusts consumption in date $i-1$ to the level such that no change in subsequent consumption is anticipated.

\section*{Random Walks}
An important example of martingales is a random walk. Let $\left\{\mathbf{g}_{i}\right\}$ be a vector independent white noise process (so it is i.i.d. with mean $\mathbf{0}$ and finite variance matrix). A random walk, $\left\{\mathbf{z}_{i}\right\}$, is a sequence of cumulative sums:


\begin{equation*}
\mathbf{z}_{1}=\mathbf{g}_{1}, \mathbf{z}_{2}=\mathbf{g}_{1}+\mathbf{g}_{2}, \ldots, \mathbf{z}_{i}=\mathbf{g}_{1}+\mathbf{g}_{2}+\cdots+\mathbf{g}_{i}, \cdots . \tag{2.2.6}
\end{equation*}


Given the sequence $\left\{\mathbf{z}_{i}\right\}$, the underlying independent white noise sequence, $\left\{\mathbf{g}_{i}\right\}$, can be backed out by taking first differences:


\begin{equation*}
\mathbf{g}_{1}=\mathbf{z}_{1}, \mathbf{g}_{2}=\mathbf{z}_{2}-\mathbf{z}_{1}, \ldots, \mathbf{g}_{i}=\mathbf{z}_{i}-\mathbf{z}_{i-1}, \ldots . \tag{2.2.7}
\end{equation*}


So the first difference of a random walk is independent white noise. A random walk is a martingale because\\
$\mathrm{E}\left(\mathbf{z}_{i} \mid \mathbf{z}_{i-1}, \ldots, \mathbf{z}_{1}\right)=\mathrm{E}\left(\mathbf{z}_{i} \mid \mathbf{g}_{i-1}, \ldots, \mathbf{g}_{1}\right)$ (since $\left(\mathbf{z}_{i-1}, \ldots, \mathbf{z}_{1}\right)$ and\\
$\left(\mathbf{g}_{i-1}, \ldots, \mathbf{g}_{1}\right)$ have the same information, as just seen)\\
$=\mathrm{E}\left(\mathbf{g}_{1}+\mathbf{g}_{2}+\cdots+\mathbf{g}_{i} \mid \mathbf{g}_{i-1}, \ldots, \mathbf{g}_{1}\right)$\\
$=\mathrm{E}\left(\mathbf{g}_{i} \mid \mathbf{g}_{i-1}, \ldots, \mathbf{g}_{1}\right)+\left(\mathbf{g}_{1}+\cdots+\mathbf{g}_{i-1}\right)$\\
$=\mathbf{g}_{1}+\cdots+\mathbf{g}_{i-1} \quad\left(\mathrm{E}\left(\mathbf{g}_{i} \mid \mathbf{g}_{i-1}, \ldots, \mathbf{g}_{1}\right)=\mathbf{0}\right.$ as $\left\{\mathbf{g}_{i}\right\}$ is independent white noise $)$\\
$=\mathbf{z}_{i-1}$ (by the definition of $\mathbf{z}_{i-1}$ ).

\section*{Martingale Difference Sequences}
A vector process $\left\{\mathbf{g}_{i}\right\}$ with $\mathrm{E}\left(\mathbf{g}_{i}\right)=\mathbf{0}$ is called a martingale difference sequence (m.d.s.) or martingale differences if the expectation conditional on its past values, too, is zero:


\begin{equation*}
\mathrm{E}\left(\mathbf{g}_{i} \mid \mathbf{g}_{i-1}, \mathbf{g}_{i-2}, \ldots, \mathbf{g}_{1}\right)=\mathbf{0} \text { for } i \geq 2 \tag{2.2.9}
\end{equation*}


The process is so called because the cumulative sum $\left\{\mathbf{z}_{i}\right\}$ created from a martingale difference sequence $\left\{\mathbf{g}_{i}\right\}$ is a martingale; the proof is the same as in (2.2.8). Conversely, if $\left\{\mathbf{z}_{i}\right\}$ is martingale, the first differences created as in (2.2.7) are a martingale difference sequence.

A martingale difference sequence has no serial correlation (i.e., $\operatorname{Cov}\left(\mathbf{g}_{i}, \mathbf{g}_{i-j}\right)=$ 0 for all $i$ and $j \neq 0$ ). A proof of this claim is as follows.

Proof. First note that we can assume, without loss of generality, that $j \geq 1$. Since the mean is zero, it suffices to show that $\mathrm{E}\left(\mathbf{g}_{i} \mathbf{g}_{i-j}^{\prime}\right)=\mathbf{0}$. So consider rewriting it as follows.

$$
\begin{aligned}
& \mathrm{E}\left(\mathbf{g}_{i} \mathbf{g}_{i-j}^{\prime}\right) \\
& =\mathrm{E}\left[\mathrm{E}\left(\mathbf{g}_{i} \mathbf{g}_{i-j}^{\prime} \mid \mathbf{g}_{i-j}\right)\right] \quad \text { (by the Law of Total Expectations) } \\
& =\mathrm{E}\left[\mathrm{E}\left(\mathbf{g}_{i} \mid \mathbf{g}_{i-j}\right) \mathbf{g}_{i-j}^{\prime}\right] \quad \text { (by the linearity of conditional expectations). }
\end{aligned}
$$

Now, since $j \geq 1,\left(\mathbf{g}_{i-1}, \ldots, \mathbf{g}_{i-j}, \ldots, \mathbf{g}_{1}\right)$ includes $\mathbf{g}_{i-j}$. Therefore,

$$
\begin{aligned}
& \mathrm{E}\left(\mathbf{g}_{i} \mid \mathbf{g}_{i-j}\right) \\
& =\mathrm{E}\left[\mathrm{E}\left(\mathbf{g}_{i} \mid \mathbf{g}_{i-1}, \ldots, \mathbf{g}_{i-j}, \ldots, \mathbf{g}_{1}\right) \mid \mathbf{g}_{i-j}\right] \quad \text { (by the Law of Iterated Expectations) } \\
& =\mathbf{0}
\end{aligned}
$$

The last equality holds because $\mathrm{E}\left(\mathbf{g}_{i} \mid \mathbf{g}_{i-1}, \ldots, \mathbf{g}_{i-j}, \ldots, \mathbf{g}_{1}\right)=\mathbf{0}$.

\section*{ARCH Processes}
An example of martingale differences, frequently used in analyzing asset returns, is an autoregressive conditional heteroskedastic (ARCH) process introduced by Engle (1982). A process $\left\{g_{i}\right\}$ is said to be an ARCH process of order 1 (ARCH(1)) if it can be written as


\begin{equation*}
g_{i}=\sqrt{\zeta+\alpha g_{i-1}^{2}} \cdot \varepsilon_{i}, \tag{2.2.10}
\end{equation*}


where $\left\{\varepsilon_{i}\right\}$ is i.i.d. with mean zero and unit variance. If $g_{1}$ is the initial value of the process, we can use (2.2.10) to calculate subsequent values of $g_{i}$. For example,

$$
g_{2}=\sqrt{\zeta+\alpha g_{1}^{2}} \cdot \varepsilon_{2} .
$$

More generally, $g_{i}(i \geq 2)$ is a function of $g_{1}$ and $\left(\varepsilon_{2}, \varepsilon_{3}, \ldots, \varepsilon_{i}\right)$. Therefore, $\varepsilon_{i}$ is independent of $\left(g_{1}, g_{2}, \ldots, g_{i-1}\right)$. It is then easy to show that $\left\{g_{i}\right\}$ is an m.d.s. because


\begin{align*}
& \mathrm{E}\left(g_{i} \mid g_{i-1}, g_{i-2}, \ldots, g_{1}\right)  \tag{2.2.11}\\
& =\mathrm{E}\left(\sqrt{\zeta+\alpha g_{i-1}^{2}} \cdot \varepsilon_{i} \mid g_{i-1}, g_{i-2}, \ldots, g_{1}\right) \\
& =\sqrt{\zeta+\alpha g_{i-1}^{2}} \mathrm{E}\left(\varepsilon_{i} \mid g_{i-1}, g_{i-2}, \ldots, g_{1}\right) \\
& =\sqrt{\zeta+\alpha g_{i-1}^{2}} \mathrm{E}\left(\varepsilon_{i}\right) \quad\left(\text { since } \varepsilon_{i} \text { is independent of }\left(g_{1}, g_{2}, \ldots, g_{i-1}\right)\right) \\
& =0 \quad\left(\text { since } \mathrm{E}\left(\varepsilon_{i}\right)=0\right) \tag{2.2.12}
\end{align*}


By a similar argument, it follows that


\begin{equation*}
\mathrm{E}\left(g_{i}^{2} \mid g_{i-1}, g_{i-2}, \ldots, g_{1}\right)=\zeta+\alpha g_{i-1}^{2} . \tag{2.2.13}
\end{equation*}


So the conditional second moment (which equals the conditional variance since $\mathrm{E}\left(g_{i} \mid g_{1}, g_{2}, \ldots, g_{i-1}\right)=0$ ) is a function of its own history of the process. In this sense the process exhibits own conditional heteroskedasticity. It can be shown (see, e.g., Engle, 1982) that the process is strictly stationary and ergodic if $|\alpha|<1$, provided that $g_{1}$ is a draw from an appropriate distribution or provided that the process started in the infinite past. If $g_{i}$ is stationary, the unconditional second moment is easy to obtain. Taking the unconditional expectation of both sides of (2.2.13) and noting that\\
$\mathrm{E}\left[\mathrm{E}\left(g_{i}^{2} \mid g_{i-1}, g_{i-2}, \ldots, g_{1}\right)\right]=\mathrm{E}\left(g_{i}^{2}\right)$ and $\mathrm{E}\left(g_{i}^{2}\right)=\mathrm{E}\left(g_{i-1}^{2}\right)$ if $g_{i}$ is stationary,\\
we obtain


\begin{equation*}
\mathrm{E}\left(g_{i}^{2}\right)=\zeta+\alpha \mathrm{E}\left(g_{i}^{2}\right) \quad \text { or } \quad \mathrm{E}\left(g_{i}^{2}\right)=\frac{\zeta}{1-\alpha} . \tag{2.2.14}
\end{equation*}


If $\alpha>0$, this model captures the characteristic found for asset returns that large values tend to be followed by large values. (For more details of ARCH processes, see, e.g., Hamilton, 1994, Section 21.1.)

\section*{Different Formulation of Lack of Serial Dependence}
Evidently, an independent white noise process is a stationary martingale difference sequence with finite variance. And, as just seen, a martingale difference sequence has no serial correlation. Thus, we have three formulations of a lack of serial dependence for zero-mean covariance stationary processes. They are, in the order of strength,


\begin{align*}
& (1) \text { " }\left\{\mathbf{g}_{i}\right\} \text { is independent white noise." } \\
\Rightarrow & (2) \text { " }\left\{\mathbf{g}_{i}\right\} \text { is stationary m.d.s. with finite variance." } \\
\Rightarrow & (3) \text { " }\left\{\mathbf{g}_{i}\right\} \text { is white noise." } \tag{2.2.15}
\end{align*}


Condition (1) is stronger than (2) because there are processes satisfying (2) but not (1). An $\mathrm{ARCH}(1)$ process (2.2.10) with $|\alpha|<1$ is an example. Figure 2.1 shows how a realization of a process satisfying (1) typically differs from that satisfying (2). Figure 2.1, Panel (a), plots a realization of a sequence of independent and normally distributed random variables with mean 0 and unit variance. Panel (b) plots an $\operatorname{ARCH}(1)$ process (2.2.10) with $\zeta=0.2$ and $\alpha=0.8$ (so that the unconditional variance, $\zeta /(1-\alpha)$, is unity as in panel (a)), where the value of the i.i.d. sequence $\varepsilon_{i}$ in (2.2.10) is taken from Figure 2.1, Panel (a), so that the sign in both panels is the same at all points. The series in Panel (b) is generally less volatile than in Panel (a), but at some points it is much more volatile. Nevertheless, the series is stationary.

Condition (2) is stronger than (3); the process in Example 2.4 is white noise, but (as you will show in a review question) it does not satisfy (2).

\section*{The CLT for Ergodic Stationary Martingale Differences Sequences}
The following CLT extends the Lindeberg-Levy CLT to stationary and ergodic m.d.s.

Ergodic Stationary Martingale Differences CLT (Billingsley, 1961): Let $\left\{\mathbf{g}_{i}\right\}$ be a vector martingale difference sequence that is stationary and ergodic with $\mathrm{E}\left(\mathbf{g}_{i} \mathbf{g}_{i}^{\prime}\right)=\mathbf{\Sigma},{ }^{9}$ and let $\overline{\mathbf{g}} \equiv \frac{1}{n} \sum_{i=1}^{n} \mathbf{g}_{i}$. Then

$$
\sqrt{n} \overline{\mathbf{g}}=\frac{1}{\sqrt{n}} \sum_{i=1}^{n} \mathbf{g}_{i} \rightarrow N(\mathbf{0}, \mathbf{\Sigma}) .
$$

\footnotetext{${ }^{9}$ Since $\left\{\mathbf{g}_{i}\right\}$ is stationary, this matrix of cross moments does not depend on $i$. Also, since a cross moment matrix is indicated, it is implicitly assumed that all the cross moments exist and are finite.
}\begin{figure}[h]
\begin{center}
  \includegraphics[alt={},max width=\textwidth]{14fb0795-6d81-4e76-9788-792298f1ddae-124_1252_1194_255_264}
\captionsetup{labelformat=empty}
\caption{Figure 2.1: Plots of Serially Uncorrelated Time Series: (a) An i.i.d. $N(0,1)$ sequence. (b) $\mathrm{ARCH}(1)$ with shocks taken from panel (a)}
\end{center}
\end{figure}

Unlike in the Lindeberg-Levy CLT, there is no need to subtract the mean from $\mathbf{g}_{i}$ because the unconditional mean of an m.d.s. is by definition zero. For the same reason, $\boldsymbol{\Sigma}$ also equals $\operatorname{Var}\left(\mathbf{g}_{i}\right)$. This CLT, being applicable not just to i.i.d. sequences but also to stationary martingale differences such as $\operatorname{ARCH}(1)$ processes, is more general than Lindeberg-Levy.

We have presented an LLN for serially correlated processes in the form of the Ergodic Theorem. A central limit theorem for serially correlated processes will be presented in Chapter 6.

\section*{QUESTIONS FOR REVIEW}
\begin{enumerate}
  \item Prove that $\boldsymbol{\Gamma}_{-j}=\boldsymbol{\Gamma}_{j}^{\prime}$. Hint: $\operatorname{Cov}\left(\mathbf{z}_{i}, \mathbf{z}_{i-j}\right)=\mathrm{E}\left[\left(\mathbf{z}_{i}-\boldsymbol{\mu}\right)\left(\mathbf{z}_{i-j}-\boldsymbol{\mu}\right)^{\prime}\right]$ where $\boldsymbol{\mu}=\mathrm{E}\left(\mathbf{z}_{i}\right)$. By covariance-stationarity, $\operatorname{Cov}\left(\mathbf{z}_{i}, \mathbf{z}_{i-j}\right)=\operatorname{Cov}\left(\mathbf{z}_{i+j}, \mathbf{z}_{i}\right)$.
  \item (Forecasting white noise) For the white noise process of Example 2.4, $\mathrm{E}\left(z_{i}\right)=$ 0 . What is $\mathrm{E}\left(z_{i} \mid z_{1}\right)$ for $i \geq 2$ ? Hint: You should be able to forecast the future exactly if you know the value of $z_{1}$. Is the process an m.d.s? [Answer: No.]
  \item (No anticipated changes in martingales) Suppose $\left\{x_{i}\right\}$ is a martingale with respect to $\left\{\mathbf{z}_{i}\right\}$. Show that $\mathrm{E}\left(x_{i+j} \mid \mathbf{z}_{i-1}, \mathbf{z}_{i-2}, \ldots, \mathbf{z}_{1}\right)=x_{i-1}$ and $\mathrm{E}\left(x_{i+j+1}-\right. \left.x_{i+j} \mid \mathbf{z}_{i-1}, \mathbf{z}_{i-2}, \ldots, \mathbf{z}_{1}\right)=0$ for $j=0,1, \ldots$. Hint: Use the Law of Iterated Expectations.
  \item Let $\left\{x_{i}\right\}$ be a sequence of real numbers that change with $i$ and $\left\{\varepsilon_{i}\right\}$ be a sequence of i.i.d. random variables with mean 0 and finite variance. Is $\left\{x_{i} \cdot \varepsilon_{i}\right\}$ i.i.d.? [Answer: No.] Is it serially independent? [Answer: Yes.] An m.d.s? [Answer: Yes.] Stationary? [Answer: No.]
  \item Show that a random walk is nonstationary. Hint: Check the variance.
  \item (The first difference of a martingale is a martingale difference sequence) Let $\left\{\mathbf{z}_{i}\right\}$ be a martingale. Show that the process $\left\{\mathbf{g}_{i}\right\}$ created by (2.2.7) is an m.d.s. Hint: $\left(\mathbf{g}_{1}, \ldots, \mathbf{g}_{i}\right)$ and $\left(\mathbf{z}_{1}, \ldots, \mathbf{z}_{i}\right)$ share the same information.
  \item (An m.d.s. that is not independent white noise) Let $g_{i}=\varepsilon_{i} \cdot \varepsilon_{i-1}$, where $\left\{\varepsilon_{i}\right\}$ is an independent white noise process. Evidently, $\left\{g_{i}\right\}$ is not i.i.d. Verify that $\left\{g_{i}\right\}(i=2,3, \ldots)$ is an m.d.s. Hint: $\mathrm{E}\left(g_{i} \mid g_{i-1}, \ldots, g_{2}\right)=\mathrm{E}\left[\mathrm{E}\left(\varepsilon_{i} \cdot \varepsilon_{i-1} \mid\right.\right. \left.\left.\varepsilon_{i-1}, \ldots, \varepsilon_{1}\right) \mid \varepsilon_{i-1} \cdot \varepsilon_{i-2}, \varepsilon_{i-2} \cdot \varepsilon_{i-3}, \ldots, \varepsilon_{2} \cdot \varepsilon_{1}\right]$.
  \item (Revision of expectations is m.d.s.) Let $\left\{y_{i}\right\}$ be a process such that $\mathrm{E}\left(y_{i} \mid\right. \left.y_{i-1}, y_{i-2}, \ldots, y_{1}\right)$ exists and is finite, and define $r_{i 1} \equiv \mathrm{E}\left(y_{i} \mid y_{i-1}, y_{i-2}, \ldots\right.$, $\left.y_{1}\right)-\mathrm{E}\left(y_{i} \mid y_{i-2}, y_{i-3}, \ldots, y_{1}\right)$. So $r_{i 1}$ is the change in the expectation as one more observation is added to the information set. Show that $\left\{r_{i 1}\right\}(i \geq 2)$ is an m.d.s. with respect to $\left\{y_{i}\right\}$.
  \item (Billingsley is stronger than Lindeberg-Levy) Let $\left\{\mathbf{z}_{i}\right\}$ be an i.i.d. sequence with $\mathrm{E}\left(\mathbf{z}_{i}\right)=\boldsymbol{\mu}$ and $\operatorname{Var}\left(\mathbf{z}_{i}\right)=\boldsymbol{\Sigma}$, as in the Lindeberg-Levy CLT. Use the Martingale Differences CLT to prove the claim of the Lindeberg-Levy CLT, namely, that $\sqrt{n}\left(\overline{\mathbf{z}}_{n}-\boldsymbol{\mu}\right) \rightarrow_{\mathrm{d}} N(\mathbf{0}, \boldsymbol{\Sigma})$. Hint: $\left\{\mathbf{z}_{i}-\boldsymbol{\mu}\right\}$ is an independent white noise process and hence is an ergodic stationary m.d.s.
\end{enumerate}

\subsection*{2.3 Large-Sample Distribution of the OLS Estimator}
The importance in econometrics of the OLS procedure, originally developed for the classical regression model of Chapter 1 , lies in the fact that it has good asymptotic properties for a class of models, different from the classical model, that are useful in economics. Of those models, the model presented in this section has probably the widest range of economic applications. No specific distributional assumption (such as the normality of the error term) is required to derive the asymptotic distribution of the OLS estimator. The requirement in finite-sample theory that the regressors be strictly exogenous or "fixed" is replaced by a much weaker requirement that they be "predetermined." (For the sake of completeness, the appendix develops the parallel asymptotic theory for a model with "fixed" regressors.)

\section*{The Model}
We use the term the data generating process (DGP) for the stochastic process that generated the finite sample $(\mathbf{y}, \mathbf{X})$. Therefore, if we specify the DGP, the joint distribution of the finite sample $(\mathbf{y}, \mathbf{X})$ can be determined. In finite-sample theory, where the sample size is fixed and finite, we defined a model as a set of the joint distributions of $(\mathbf{y}, \mathbf{X})$. In large-sample theory, a model is stated as a set of DGPs. The model we study is the set of DGPs satisfying the following set of assumptions.

Assumption 2.1 (linearity):

$$
y_{i}=\mathbf{x}_{i}^{\prime} \boldsymbol{\beta}+\varepsilon_{i} \quad(i=1,2, \ldots, n)
$$

where $\mathbf{x}_{i}$ is a $K$-dimensional vector of explanatory variables (regressors), $\boldsymbol{\beta}$ is a $K$-dimensional coefficient vector, and $\varepsilon_{i}$ is the unobservable error term.

Assumption 2.2 (ergodic stationarity): The $(K+1)$-dimensional vector stochastic process $\left\{y_{i}, \mathbf{x}_{i}\right\}$ is jointly stationary and ergodic.

Assumption 2.3 (predetermined regressors): All the regressors are predetermined in the sense that they are orthogonal to the contemporaneous error term: $\mathrm{E}\left(x_{i k} \varepsilon_{i}\right)=0$ for all $i$ and $k(=1,2, \ldots, K) .^{10}$ This can be written as

\footnotetext{${ }^{10}$ Our definition of the term predetermined is not universal. Some authors say that the regressors are predetermined if $\mathrm{E}\left(\mathbf{x}_{i-j} \cdot \varepsilon_{i}\right)=\mathbf{0}$ for all $j \geq 0$, not just for $j=0$. That is, the error term is orthogonal not only to the contemporaneous but also the past regressors. In Koopmans and Hood (1953), the regressors are said to be predetermined if $\varepsilon_{i}$ is independent of $\mathbf{x}_{i-j}$ for all $j \geq 0$. Our definition is the same as Hamilton's (1994).
}
$$
\mathrm{E}\left[\mathbf{x}_{i} \cdot\left(y_{i}-\mathbf{x}_{i}^{\prime} \boldsymbol{\beta}\right)\right]=\mathbf{0} \text { or equivalently } \mathrm{E}\left(\mathbf{g}_{i}\right)=\mathbf{0}
$$
where $\mathbf{g}_{i} \equiv \mathbf{x}_{i} \cdot \varepsilon_{i}$.

Assumption 2.4 (rank condition): The $K \times K$ matrix $\mathrm{E}\left(\mathbf{x}_{i} \mathbf{x}_{i}^{\prime}\right)$ is nonsingular (and hence finite). We denote this matrix by $\mathbf{\Sigma}_{\mathbf{x x}}$.

Assumption 2.5 ( $\mathbf{g}_{i}$ is a martingale difference sequence with finite second moments): $\left\{\mathbf{g}_{i}\right\}$ is a martingale difference sequence (so a fortiori $\mathrm{E}\left(\mathbf{g}_{i}\right)=\mathbf{0}$ ). The $K \times K$ matrix of cross moments, $\mathrm{E}\left(\mathbf{g}_{i} \mathbf{g}_{i}^{\prime}\right)$, is nonsingular. We use $\mathbf{S}$ for $\operatorname{Avar}(\overline{\mathbf{g}})$ (the variance of the asymptotic distribution of $\sqrt{n} \overline{\mathbf{g}}$, where $\overline{\mathbf{g}} \equiv \frac{1}{n} \sum_{i} \mathbf{g}_{i}$ ). By Assumption 2.2 and the ergodic stationary Martingale Differences CLT, S $=\mathrm{E}\left(\mathbf{g}_{i} \mathbf{g}_{i}^{\prime}\right)$.

The first assumption is just reproducing Assumption 1.1. The rest of the assumptions require some lengthy comments.

\begin{itemize}
  \item (Ergodic stationarity) A trivial but important special case of ergodic stationarity is that $\left\{y_{i}, \mathbf{x}_{i}\right\}$ is i.i.d., that is, the sample is a random sample. ${ }^{11}$ Most existing microdata on households are random samples, with observations randomly drawn from a population of a nation's households. Thus, we are in no way ruling out models that use cross-section data.
  \item (The model accommodates conditional heteroskedasticity) If $\left\{y_{i}, \mathbf{x}_{i}\right\}$ is stationary, then the error term $\varepsilon_{i}=y_{i}-\mathbf{x}_{i}^{\prime} \boldsymbol{\beta}$ is also stationary. Thus, Assumption 2.2 implies that the unconditional second moment $\mathrm{E}\left(\varepsilon_{i}^{2}\right)$-if it exists and is finite-is constant across $i$. That is, the error term is unconditionally homoskedastic. Yet the error can be conditionally heteroskedastic in that the conditional second moment, $\mathrm{E}\left(\varepsilon_{i}^{2} \mid \mathbf{x}_{i}\right)$, can depend on $\mathbf{x}_{i}$. An example in which the error is homoskedastic unconditionally but not conditionally is included in Section 2.6, where the consequence of superimposing conditional homoskedasticity (that $\mathrm{E}\left(\varepsilon_{i}^{2} \mid \mathbf{x}_{i}\right)=\sigma^{2}$ ) on the model will be explored.
  \item $\left(\mathrm{E}\left(\mathbf{x}_{i} \cdot \varepsilon_{i}\right)=\mathbf{0}\right.$ vs. $\left.\mathrm{E}\left(\varepsilon_{i} \mid \mathbf{x}_{i}\right)=0\right)$ Sometimes, instead of the orthogonality condition $\mathrm{E}\left(\mathbf{x}_{i} \cdot \varepsilon_{i}\right)=\mathbf{0}$, it is assumed that the error is unrelated in the sense that $\mathrm{E}\left(\varepsilon_{i} \mid \mathbf{x}_{i}\right)=0$. This is stronger than the orthogonality condition because it
\end{itemize}

\footnotetext{${ }^{11}$ Actually, once the independence assumption is made, the same large-sample results can be proved for the more general case where $\left(y_{i}, \mathbf{x}_{i}\right\}$ is independently but not identically distributed (i.n.i.d.), provided that some conditions on higher moments of the joint distribution of ( $\varepsilon_{i}, \mathbf{x}_{i}$ ) are satisfied. We will not entertain this generalization because the i.i.d. assumption is satisfied in most microdata sets.
}
implies that, for any (measurable) function $f$ of $\mathbf{x}_{i}, f\left(\mathbf{x}_{i}\right)$ is orthogonal to $\varepsilon_{i}$ :
$$
\mathrm{E}\left[f\left(\mathbf{x}_{i}\right) \varepsilon_{i}\right]=\mathrm{E}\left[\mathrm{E}\left(f\left(\mathbf{x}_{i}\right) \varepsilon_{i} \mid \mathbf{x}_{i}\right)\right]=\mathrm{E}\left[f\left(\mathbf{x}_{i}\right) \mathrm{E}\left(\varepsilon_{i} \mid \mathbf{x}_{i}\right)\right]=0 .
$$

In rational expectations models, this stronger condition is satisfied, but for the purpose of developing asymptotic theory, we need only the weaker assumption of the orthogonality condition.

\begin{itemize}
  \item (Predetermined vs. strictly exogenous regressors) The regressors are not required to be strictly exogenous. As we noted in Section 1.2, the exogeneity assumption (Assumption 1.2) implies that, for any regressor $k, \mathrm{E}\left(x_{j k} \varepsilon_{i}\right)=0$ for all $i$ and $j$, not just for $i=j$, which rules out the possibility that the current error term, $\varepsilon_{i}$, is correlated with future regressors, $\mathbf{x}_{i+j}$ for $j \geq 1$. Assumption 2.3, restricting only the contemporaneous relationship between the error term and the regressors, does not rule out that possibility. For example, the AR(1) process, which does not satisfy the exogeneity assumption of the classical regression model, can be accommodated in the model of this chapter. This weaker assumption of predetermined regressors will be further relaxed in the next chapter.
  \item (Rank condition as no multicollinearity in the limit) Since $\mathrm{E}\left(\mathbf{x}_{i} \mathbf{x}_{i}^{\prime}\right)$ is finite by Assumption 2.4, $\lim _{n \rightarrow \infty} \mathbf{S}_{\mathbf{x x}}=\mathbf{\Sigma}_{\mathbf{x x}}$ (where $\mathbf{S}_{\mathbf{x x}} \equiv \frac{1}{n} \sum_{i=1}^{n} \mathbf{x}_{i} \mathbf{x}_{i}^{\prime}$ ) with probability one by the Ergodic Theorem. So, for $n$ sufficiently large, the sample cross moment of the regressors $\mathbf{S}_{\mathbf{x x}}$, which can be written as $\frac{1}{n} \mathbf{X}^{\prime} \mathbf{X}$, is nonsingular by Assumptions 2.2 and 2.4. Since $\frac{1}{n} \mathbf{X}^{\prime} \mathbf{X}$ is nonsingular if and only if $\operatorname{rank}(\mathbf{X})=K$, Assumption 1.3 (no multicollinearity) is satisfied with probability one for sufficiently large $n$. In the OLS formula $\mathbf{b}=\mathbf{S}_{\mathbf{x x}}^{-1} \mathbf{S}_{\mathbf{x y}}$ (where $\left.\mathbf{s}_{\mathbf{x y}} \equiv \frac{1}{n} \sum_{i=1}^{n} \mathbf{x}_{i} \cdot y_{i}\right), \mathbf{S}_{\mathbf{x x}}$ needs to be inverted. If $\mathbf{S}_{\mathbf{x x}}$ is singular in a finite sample (so it cannot be inverted), we just assign an arbitrary value to $\mathbf{b}$ so that the OLS estimator is well-defined for any sample.
  \item (A sufficient condition for $\left\{\mathbf{g}_{i}\right\}$ to be an m.d.s.) Since an m.d.s. (martingale difference sequence) is zero-mean by definition, Assumption 2.5 is stronger than Assumption 2.3. We will need Assumption 2.5 to prove the asymptotic normality of the OLS estimator. The assumption, about the product of the regressors and the error term, may be hard to interpret. A sufficient condition that is easier to interpret is
\end{itemize}


\begin{equation*}
\mathrm{E}\left(\varepsilon_{i} \mid \varepsilon_{i-1}, \varepsilon_{i-2}, \ldots, \varepsilon_{1}, \mathbf{x}_{i}, \mathbf{x}_{i-1}, \ldots, \mathbf{x}_{1}\right)=0 . \tag{2.3.1}
\end{equation*}


Note that the current as well as lagged regressors is included in the information\\
set. This condition implies that the error term is serially uncorrelated and also is uncorrelated with the current and past regressors (the proof is much like the proof in the previous section that an m.d.s. is serially uncorrelated). That (2.3.1) is sufficient for $\left\{\mathbf{g}_{i}\right\}$ to be an m.d.s. can be seen as follows. We have

$$
\begin{aligned}
& \mathrm{E}\left(\mathbf{g}_{i} \mid \mathbf{g}_{i-1}, \ldots, \mathbf{g}_{1}\right) \\
& =\mathrm{E}\left[\mathrm{E}\left(\mathbf{g}_{i} \mid \varepsilon_{i-1}, \varepsilon_{i-2}, \ldots, \varepsilon_{1}, \mathbf{x}_{i}, \mathbf{x}_{i-1}, \ldots, \mathbf{x}_{1}\right) \mid \mathbf{g}_{i-1}, \ldots, \mathbf{g}_{1}\right] .
\end{aligned}
$$

This holds by the Law of Iterated Expectations because there is more information in the "inside" information set $\left(\varepsilon_{i-1}, \varepsilon_{i-2}, \ldots, \varepsilon_{1}, \mathbf{x}_{i}, \mathbf{x}_{i-1}, \ldots, \mathbf{x}_{1}\right)$ than in the "outside" information set $\left(\mathbf{g}_{i-1}, \ldots, \mathbf{g}_{1}\right)$. Therefore,

$$
\begin{aligned}
& \mathrm{E}\left(\mathbf{g}_{i} \mid \mathbf{g}_{i-1}, \ldots, \mathbf{g}_{1}\right) \\
& =\mathrm{E}\left[\mathbf{x}_{i} \mathrm{E}\left(\varepsilon_{i} \mid \varepsilon_{i-1}, \varepsilon_{i-2}, \ldots, \varepsilon_{1}, \mathbf{x}_{i}, \mathbf{x}_{i-1}, \ldots, \mathbf{x}_{1}\right) \mid \mathbf{g}_{i-1}, \ldots, \mathbf{g}_{1}\right]
\end{aligned}
$$

(by the linearity of conditional expectations)


\begin{equation*}
=\mathbf{0} \quad(\text { by }(2.3 .1)) . \tag{2.3.2}
\end{equation*}


\begin{itemize}
  \item (When the regressors include a constant) In virtually all applications, the regressors include a constant. If the regressors include a constant so that $x_{i 1}=1$ for all $i$, then Assumption 2.3 of predetermined regressors can be stated in more familiar terms: the mean of the error term is zero (which is implied by $\mathrm{E}\left(x_{i k} \varepsilon_{i}\right)=0$ for $k=1$ ), and the contemporaneous correlation between the error term and the regressors is zero (which is implied by $\mathrm{E}\left(x_{i k} \varepsilon_{i}\right)$ for $k \neq 1$ and $\mathrm{E}\left(\varepsilon_{i}\right)=0$ ). Also, since the first element of the $K$-dimensional vector $\mathbf{g}_{i}\left(\equiv \mathbf{x}_{i} \cdot \varepsilon_{i}\right)$ is $\varepsilon_{i}$, Assumption 2.5 implies
\end{itemize}

$$
\mathrm{E}\left(\varepsilon_{i} \mid \mathbf{g}_{i-1}, \mathbf{g}_{i-2}, \ldots, \mathbf{g}_{1}\right)=0 .
$$

Then, by the Law of Iterated Expectations, $\left\{\varepsilon_{i}\right\}$ is a scalar m.d.s:


\begin{equation*}
\mathrm{E}\left(\varepsilon_{i} \mid \varepsilon_{i-1}, \varepsilon_{i-2}, \ldots, \varepsilon_{1}\right)=0 \tag{2.3.3}
\end{equation*}


Therefore, Assumption 2.5 implies that the error term itself is an m.d.s. and hence is serially uncorrelated.

\begin{itemize}
  \item ( $\mathbf{S}$ is a matrix of fourth moments) Since $\mathbf{g}_{i} \equiv \mathbf{x}_{i} \cdot \varepsilon_{i}$, the $\mathbf{S}$ in Assumption 2.5 can be written as $\mathrm{E}\left(\varepsilon_{i}^{2} \mathbf{x}_{i} \mathbf{x}_{i}^{\prime}\right)$. Its $(k, j)$ element is $\mathrm{E}\left(\varepsilon_{i}^{2} x_{i k} x_{i j}\right)$. So $\mathbf{S}$ is a matrix of fourth moments (the expectation of products of four different variables). Consistent estimation of $\mathbf{S}$ will require an additional assumption to be specified in Section 2.5.
  \item ( $\mathbf{S}$ will take a different expression without Assumption 2.5) Thanks to the assumption that $\left\{\mathbf{g}_{i}\right\}$ is an m.d.s., $\mathbf{S}(\equiv \operatorname{Avar}(\overline{\mathbf{g}}))$ is equal to $\mathrm{E}\left(\mathbf{g}_{i} \mathbf{g}_{i}^{\prime}\right)$. Without the assumption, as we will see in Chapter 6, the expression for $\mathbf{S}$ is more complicated and involves autocovariances of $\mathbf{g}_{i}$.
\end{itemize}

\section*{Asymptotic Distribution of the OLS Estimator}
We now prove that the OLS estimator is consistent and asymptotically normal. It should be kept in mind throughout the rest of this chapter that the OLS estimator b depends on the sample size $n$ (although the dependence is not made explicit by our choice of not to subscript $\mathbf{b}$ by $n$ ) and that $K$, the number of regressors, is held fixed when we track the sequence of OLS estimators indexed by $n$. For the time being, we presume that there is available some consistent estimator, denoted $\widehat{\mathbf{S}}$, of $\mathbf{S}\left(\equiv \operatorname{Avar}(\overline{\mathbf{g}})=\mathrm{E}\left(\mathbf{g}_{i} \mathbf{g}_{i}^{\prime}\right)=\mathrm{E}\left(\varepsilon_{i}^{2} \mathbf{x}_{i} \mathbf{x}_{i}^{\prime}\right)\right)$. The issue of estimating $\mathbf{S}$ consistently will be taken up later.

\section*{Proposition 2.1 (asymptotic distribution of the OLS Estimator):}
(a) (Consistency of $\mathbf{b}$ for $\boldsymbol{\beta}$ ) Under Assumptions 2.1-2.4, $\operatorname{plim}_{n \rightarrow \infty} \mathbf{b}=\boldsymbol{\beta}$. (So Assumption 2.5 is not needed for consistency.)\\
(b) (Asymptotic Normality of b) If Assumption 2.3 is strengthened as Assumption 2.5, then

$$
\sqrt{n}(\mathbf{b}-\boldsymbol{\beta}) \underset{\mathrm{d}}{\rightarrow} N(\mathbf{0}, \operatorname{Avar}(\mathbf{b})) \text { as } n \rightarrow \infty
$$

where


\begin{equation*}
\operatorname{Avar}(\mathbf{b})=\boldsymbol{\Sigma}_{\mathbf{x x}}^{-1} \mathbf{S} \boldsymbol{\Sigma}_{\mathbf{x x}}^{-1} \tag{2.3.4}
\end{equation*}


(Recall: $\left.\mathbf{\Sigma}_{\mathbf{x x}} \equiv \mathrm{E}\left(\mathbf{x}_{i} \mathbf{x}_{i}^{\prime}\right), \mathbf{S}=\mathrm{E}\left(\mathbf{g}_{i} \mathbf{g}_{i}^{\prime}\right), \mathbf{g}_{i} \equiv \mathbf{x}_{i} \cdot \varepsilon_{i}.\right)$\\
(c) (Consistent Estimate of Avar(b)) Suppose there is available a consistent estimator, $\widehat{\mathbf{S}}$, of $\mathbf{S}(K \times K)$. Then, under Assumption 2.2, Avar(b) is consistently estimated by


\begin{equation*}
\widehat{\operatorname{Avar}(\mathbf{b})}=\mathbf{S}_{\mathbf{k x}}^{-1} \widehat{\mathbf{S}} \mathbf{S}_{\mathbf{x x}}^{-1} \tag{2.3.5}
\end{equation*}


where $\mathbf{S}_{\mathbf{x x}}$ is the sample mean of $\mathbf{x}_{i} \mathbf{x}_{i}^{\prime}$ :


\begin{equation*}
\underset{(K \times K)}{\mathbf{S}_{\mathbf{x x}}} \equiv \frac{1}{n} \sum_{i=1}^{n} \mathbf{x}_{i} \mathbf{x}_{i}^{\prime}=\frac{1}{n} \mathbf{X}^{\prime} \mathbf{X} \tag{2.3.6}
\end{equation*}


The proof is a showcase of all the standard tricks in asymptotics. For proving (a) and (b), three tricks will be employed: (1) write the object in question in terms of sample means, (2) apply the relevant LLN (the Ergodic Theorem in the present context) and CLT (the ergodic stationary Martingale Differences CLT) to sample means, and (3) use Lemma 2.4(c) to derive the asymptotic distribution. Proof of (c) will not be given because it is an immediate implication of ergodic stationarity.

Proof (Parts (a) and (b)).\\
(1) We first write the sampling error $\mathbf{b}-\boldsymbol{\beta}$ in terms of sample means.


\begin{align*}
\mathbf{b}-\boldsymbol{\beta} & =\left(\mathbf{X}^{\prime} \mathbf{X}\right)^{-1} \mathbf{X}^{\prime} \boldsymbol{\varepsilon} \\
& =\left(\frac{1}{n} \mathbf{X}^{\prime} \mathbf{X}\right)^{-1}\left(\frac{1}{n} \mathbf{X}^{\prime} \boldsymbol{\varepsilon}\right) \\
& =\left(\frac{1}{n} \sum_{i=1}^{n} \mathbf{x}_{i} \mathbf{x}_{i}^{\prime}\right)^{-1}\left(\frac{1}{n} \sum_{i=1}^{n} \mathbf{x}_{i} \cdot \varepsilon_{i}\right) \\
& \equiv \mathbf{S}_{\mathbf{x x}}^{-1} \overline{\mathbf{g}} \tag{2.3.7}
\end{align*}


where

$$
\overline{\mathbf{g}} \equiv \frac{1}{n} \sum_{i=1}^{n} \mathbf{g}_{i}, \mathbf{g}_{i} \equiv \mathbf{x}_{i} \cdot \varepsilon_{i}
$$

The sample means $\mathbf{S}_{\mathbf{x x}}$ and $\overline{\mathbf{g}}$ depend on the sample size $n$, although the notation does not make it explicit.\\
(2) (Consistency) Since by Assumption $2.2\left\{\mathbf{x}_{i} \mathbf{x}_{i}^{\prime}\right\}$ is ergodic stationary, $\mathbf{S}_{\mathbf{x x}} \rightarrow_{\mathrm{p}} \boldsymbol{\Sigma}_{\mathbf{x} \mathbf{x}}$. (The convergence is actually almost surely, but almost sure convergence implies convergence in probability.) Since $\mathbf{\Sigma}_{\mathbf{x x}}$ is invertible by Assumption 2.4, $\mathbf{S}_{\mathbf{x x}}^{-1} \rightarrow_{\mathrm{p}} \boldsymbol{\Sigma}_{\mathbf{x x}}^{-1}$ by Lemma 2.3(a). Similarly, $\overline{\mathbf{g}} \rightarrow_{\mathrm{p}} \mathrm{E}\left(\mathbf{g}_{i}\right)$ which by Assumption 2.3 is $\mathbf{0}$. So by Lemma 2.3(a), $\mathbf{S}_{\mathbf{x x}}^{-1} \overline{\mathbf{g}} \rightarrow_{\mathbf{p}} \boldsymbol{\Sigma}_{\mathbf{x x}}^{-1} \mathbf{0}=\mathbf{0}$. Therefore, $\operatorname{plim}_{n \rightarrow \infty}(\mathbf{b}- \boldsymbol{\beta})=\mathbf{0}$, which implies plim $_{n \rightarrow \infty} \mathbf{b}=\boldsymbol{\beta}$.\\
(3) (Asymptotic normality) Rewrite (2.3.7) as


\begin{equation*}
\sqrt{n}(\mathbf{b}-\boldsymbol{\beta})=\mathbf{S}_{\mathbf{x x}}^{-1}(\sqrt{n} \overline{\mathbf{g}}) \tag{2.3.8}
\end{equation*}


As mentioned in the statement of Assumption 2.5, $\sqrt{n} \overline{\mathbf{g}} \rightarrow_{\mathrm{d}} N(\mathbf{0}, \mathbf{S})$. So, by Lemma 2.4(c), $\sqrt{n}(\mathbf{b}-\boldsymbol{\beta})$ converges to a normal distribution with mean $\mathbf{0}$ and variance $\boldsymbol{\Sigma}_{\mathbf{x x}}^{-1} \mathbf{S}\left(\boldsymbol{\Sigma}_{\mathbf{x x}}^{-1}\right)^{\prime}$. But since $\boldsymbol{\Sigma}_{\mathbf{x x}}$ is symmetric, this expression equals (2.3.4).

This result says that the distribution of $\sqrt{n}$ times the sampling error is approximated arbitrarily well by a normal distribution when the sample size is sufficiently large. The natural question is how large is "large": how large must the sample size be for the asymptotic approximation to be valid? The asymptotic result we just derived holds for all the DGPs satisfying the model assumptions. However, the sample size needed to achieve a given measure of proximity to the asymptotic distribution depends on the DGP. We will partially address this issue in the Monte Carlo experiment of this chapter.

\section*{$\boldsymbol{s}^{\mathbf{2}}$ Is Consistent}
We now turn to the OLS estimator, $s^{2}$, of the error variance.

Proposition 2.2 (consistent estimation of error variance): Let $e_{i} \equiv y_{i}-\mathbf{x}_{i}^{\prime} \mathbf{b}$ be the OLS residual for observation $i$. Under Assumptions 2.1-2.4,

$$
s^{2} \equiv \frac{1}{n-K} \sum_{i=1}^{n} e_{i}^{2} \underset{\mathrm{p}}{\rightarrow} \mathrm{E}\left(\varepsilon_{i}^{2}\right)
$$

provided $\mathrm{E}\left(\varepsilon_{i}^{2}\right)$ exists and is finite.

If we could observe the error term $\varepsilon_{i}$, then the obvious estimator would be the sample mean of $\varepsilon_{i}^{2}$. It is consistent by ergodic stationarity. The message of Proposition 2.2 is that the substitution of the OLS residual $e_{i}$ for the true error term $\varepsilon_{i}$ does not impair consistency. Let us go through a sketch of the proof, because knowing how to handle the discrepancy between $\varepsilon_{i}$ and its estimate $e_{i}$ will be useful in other contexts as well. Since

$$
s^{2}=\frac{n}{n-K}\left(\frac{1}{n} \sum_{i=1}^{n} e_{i}^{2}\right)
$$

it suffices to prove that the sample mean of $e_{i}^{2}, \frac{1}{n} \sum_{i} e_{i}^{2}$, converges in probability to $\mathrm{E}\left(\varepsilon_{i}^{2}\right)$. The relationship between $e_{i}$ and $\varepsilon_{i}$ is given by


\begin{align*}
e_{i} & \equiv y_{i}-\mathbf{x}_{i}^{\prime} \mathbf{b} \\
& =y_{i}-\mathbf{x}_{i}^{\prime} \boldsymbol{\beta}-\mathbf{x}_{i}^{\prime}(\mathbf{b}-\boldsymbol{\beta}) \quad\left(\text { by adding and subtracting } \mathbf{x}_{i}^{\prime} \boldsymbol{\beta}\right) \\
& =\varepsilon_{i}-\mathbf{x}_{i}^{\prime}(\mathbf{b}-\boldsymbol{\beta}), \tag{2.3.9}
\end{align*}


so that


\begin{equation*}
e_{i}^{2}=\varepsilon_{i}^{2}-2(\mathbf{b}-\boldsymbol{\beta})^{\prime} \mathbf{x}_{i} \cdot \varepsilon_{i}+(\mathbf{b}-\boldsymbol{\beta})^{\prime} \mathbf{x}_{i} \mathbf{x}_{i}^{\prime}(\mathbf{b}-\boldsymbol{\beta}) . \tag{2.3.10}
\end{equation*}


Summing over $i$, we obtain


\begin{align*}
\frac{1}{n} \sum_{i=1}^{n} e_{i}^{2} & =\frac{1}{n} \sum_{i=1}^{n} \varepsilon_{i}^{2}-2(\mathbf{b}-\boldsymbol{\beta})^{\prime} \frac{1}{n} \sum_{i=1}^{n} \mathbf{x}_{i} \cdot \varepsilon_{i}+(\mathbf{b}-\boldsymbol{\beta})^{\prime}\left(\frac{1}{n} \sum_{i=1}^{n} \mathbf{x}_{i} \mathbf{x}_{i}^{\prime}\right)(\mathbf{b}-\boldsymbol{\beta}) \\
& =\frac{1}{n} \sum_{i=1}^{n} \varepsilon_{i}^{2}-2(\mathbf{b}-\boldsymbol{\beta})^{\prime} \mathbf{g}+(\mathbf{b}-\boldsymbol{\beta})^{\prime} \mathbf{S}_{\mathbf{x x}}(\mathbf{b}-\boldsymbol{\beta}) \tag{2.3.11}
\end{align*}


The rest of the proof, which is to show that the plims of the last two terms are zero so that $\operatorname{plim} \frac{1}{n} \sum_{i} \varepsilon_{i}^{2}=\operatorname{plim} \frac{1}{n} \sum_{i} e_{i}^{2}$, is left as a review question. If you do the review question, it should be clear to you that all that is required for the coefficient estimator is consistency; if some consistent estimator, rather than the OLS estimator $\mathbf{b}$, is used to form the residuals, the error variance estimator is still consistent for $\mathrm{E}\left(\varepsilon_{i}^{2}\right)$.

\section*{QUESTIONS FOR REVIEW}
\begin{enumerate}
  \item Suppose $\mathrm{E}\left(y_{i} \mid \mathbf{x}_{i}\right)=\mathbf{x}_{i}^{\prime} \boldsymbol{\beta}$, that is, suppose the regression of $y_{i}$ on $\mathbf{x}_{i}$ is a linear function of $\mathbf{x}_{i}$. Define $\varepsilon_{i} \equiv y_{i}-\mathbf{x}_{i}^{\prime} \boldsymbol{\beta}$. Show that $\mathbf{x}_{i}$ is orthogonal to $\varepsilon_{i}$. Hint: First show that $\mathrm{E}\left(\varepsilon_{i} \mid \mathbf{x}_{i}\right)=0$.
  \item (Is $\mathrm{E}\left(\varepsilon_{i}^{2}\right)$ assumed to be finite?)\\
(a) Do Assumptions $2.1-2.5$ imply that $\mathrm{E}\left(\varepsilon_{i}^{2}\right)$ exists and is finite? Hint: A strictly stationary process may not have finite second moments.\\
(b) If one of the regressors is a constant in our model, then the variance of the error term is finite. Prove this. Hint: If $x_{i 1}=1$, the $(1,1)$ element of $\mathrm{E}\left(\mathbf{g}_{i} \mathbf{g}_{i}^{\prime}\right)$ is $\varepsilon_{i}^{2}$.
  \item (Alternative expression for $\mathbf{S}$ ) Let $f\left(\mathbf{x}_{i}\right) \equiv \mathrm{E}\left(\varepsilon_{i}^{2} \mid \mathbf{x}_{i}\right)$. Show that $\mathbf{S}(\equiv \left.\mathrm{E}\left(\varepsilon_{i}^{2} \mathbf{x}_{i} \mathbf{x}_{i}^{\prime}\right)\right)$ can be written as
\end{enumerate}

$$
\mathbf{S}=\mathrm{E}\left[f\left(\mathbf{x}_{i}\right) \mathbf{x}_{i} \mathbf{x}_{i}^{\prime}\right]
$$

Hint: Law of Total Expectations.\\
4. Complete the proof of Proposition 2.2. Hint: We have already proved for Proposition 2.1 that plim $\overline{\mathbf{g}}=\mathbf{0}$, $\operatorname{plim} \mathbf{S}_{\mathbf{x x}}=\boldsymbol{\Sigma}_{\mathbf{x x}}$, and $\operatorname{plim}(\mathbf{b}-\boldsymbol{\beta})=\mathbf{0}$ under Assumptions 2.1-2.4. Use Lemma 2.3(a) to show that plim( $\mathbf{b}-\boldsymbol{\beta})^{\prime} \mathbf{g}=0$ and $\operatorname{plim}(\mathbf{b}-\boldsymbol{\beta})^{\prime} \mathbf{S}_{\mathbf{x x}}(\mathbf{b}-\boldsymbol{\beta})=0$.\\
5. (Proposition 2.2 with consistent $\widehat{\boldsymbol{\beta}}$ ) Prove the following generalization of Proposition 2.2:

Let $\hat{\varepsilon}_{i} \equiv y_{i}-\mathbf{x}_{i}^{\prime} \widehat{\boldsymbol{\beta}}$ where $\widehat{\boldsymbol{\beta}}$ is any consistent estimator of $\boldsymbol{\beta}$. Under Assumptions 2.1,2.2, and the assumption that $\mathrm{E}\left(\mathbf{x}_{i} \cdot \varepsilon_{i}\right)$ and $\mathrm{E}\left(\mathbf{x}_{i} \mathbf{x}_{i}^{\prime}\right)$ are finite, $\frac{1}{n} \sum_{i} \hat{\varepsilon}_{i}^{2} \rightarrow_{\mathrm{p}} \mathrm{E}\left(\varepsilon_{i}^{2}\right)$.

So the regressors do not have to be orthogonal to the error term.

\subsection*{2.4 Hypothesis Testing}
Statistical inference in large-sample theory is based on test statistics whose asymptotic distributions are known under the truth of the null hypothesis. Derivation of the distribution of test statistics is easier than in finite-sample theory because we are only concerned about the large-sample approximation to the exact distribution. In this section we derive test statistics, assuming throughout that a consistent estimator, $\widehat{\mathbf{S}}$, of $\mathbf{S}\left(\equiv \mathrm{E}\left(\mathbf{g}_{i} \mathbf{g}_{i}^{\prime}\right)\right)$ is available. The issue of consistent estimation of $\mathbf{S}$ will be taken up in the next section.

\section*{Testing Linear Hypotheses}
Consider testing a hypothesis about the $k$-th coefficient $\beta_{k}$. Proposition 2.1 implies that under the $\mathrm{H}_{0}: \beta_{k}=\bar{\beta}_{k}$,

$$
\left.\sqrt{n}\left(b_{k}-\bar{\beta}_{k}\right) \underset{\mathrm{d}}{\rightarrow} N\left(0, A \operatorname{var}\left(b_{k}\right)\right) \quad \text { and } \quad \widehat{\operatorname{Avar}\left(b_{k}\right.}\right) \underset{\mathrm{p}}{\rightarrow} \operatorname{Avar}\left(b_{k}\right),
$$

where $b_{k}$ is the $k$-th element of $\mathbf{b}$ and $\operatorname{Avar}\left(b_{k}\right)$ is the $(k, k)$ element of the $K \times K$ matrix Avar(b). So Lemma 2.4(c) guarantees that


\begin{equation*}
t_{k} \equiv \frac{\sqrt{n}\left(b_{k}-\bar{\beta}_{k}\right)}{\sqrt{\widehat{\operatorname{Avar}\left(b_{k}\right)}}}=\frac{b_{k}-\bar{\beta}_{k}}{S E^{*}\left(b_{k}\right)} \rightarrow N(0,1) \tag{2.4.1}
\end{equation*}


where

$$
S E^{*}\left(b_{k}\right) \equiv \sqrt{\frac{1}{n} \cdot \widehat{\operatorname{Avar}\left(b_{k}\right)}} \equiv \sqrt{\frac{1}{n} \cdot\left(\mathbf{S}_{\mathbf{x x}}^{-1} \widehat{\mathbf{S}} \mathbf{S}_{\mathbf{x x}}^{-1}\right)_{k k}}
$$

The denominator in this $t$-ratio, $S E^{*}\left(b_{k}\right)$, is called the heteroskedasticityconsistent standard error, (heteroskedasticity-)robust standard error, or White's standard error. The reason for this terminology is that the error term can be conditionally heteroskedastic; recall that we have not assumed conditional homoskedasticity (that $\mathrm{E}\left(\varepsilon_{i}^{2} \mid \mathbf{x}_{i}\right)$ does not depend on $\mathbf{x}_{i}$ ) to derive the asymptotic distribution of $t_{k}$. This $t$-ratio is called the robust $t$-ratio, to distinguish it from the\\
$t$-ratio of Chapter 1. The relationship between these two sorts of $t$-ratio will be discussed in Section 2.6.

Given this $t$-ratio, testing the null hypothesis $\mathrm{H}_{0}: \beta_{k}=\bar{\beta}_{k}$ at a significance level of $\alpha$ proceeds as follows:

Step 1: Calculate $t_{k}$ by the formula (2.4.1).\\
Step 2: Look up the table of $N(0,1)$ to find the critical value $t_{\alpha / 2}$ which leaves $\alpha / 2$ to the upper tail of the standard normal distribution. (example: if $\alpha=5 \%$, $t_{\alpha / 2}=1.96$.)\\
Step 3: Accept the hypothesis if $\left|t_{k}\right|<t_{\alpha / 2}$; otherwise reject.\\
The differences from the finite-sample $t$-test are: (1) the way the standard error is calculated is different, (2) we use the table of $N(0,1)$ rather than that of $t(n-K)$, and (3) the actual size or exact size of the test (the probability of Type I error given the sample size) equals the nominal size (i.e., the desired significance level $\alpha$ ) only approximately, although the approximation becomes arbitrarily good as the sample size increases. The difference between the exact size and the nominal size of a test is called the size distortion. Since $t_{k}$ is asymptotically standard normal, the size distortion of the $t$-test converges to zero as the sample size $n$ goes to infinity.

Thus, we have proved the first half of

Proposition 2.3 (robust $t$-ratio and Wald statistic): Suppose Assumptions 2.12.5 hold, and suppose there is available a consistent estimate $\widehat{\mathbf{S}}$ of $\mathbf{S}\left(=\mathrm{E}\left(\mathbf{g}_{i} \mathbf{g}_{i}^{\prime}\right)\right)$. As before, let

$$
\widehat{\operatorname{Avar}(\mathbf{b})} \equiv \mathbf{S}_{\mathbf{x x}}^{-1} \widehat{\mathbf{S}} \mathbf{S}_{\mathbf{x x}}^{-1}
$$

Then\\
(a) Under the null hypothesis $\mathrm{H}_{0}: \beta_{k}=\bar{\beta}_{k}, t_{k}$ defined in (2.4.1) $\rightarrow_{\mathrm{d}} N(0,1)$.\\
(b) Under the null hypothesis $\mathrm{H}_{0}: \mathbf{R} \boldsymbol{\beta}=\mathbf{r}$, where $\mathbf{R}$ is an \# $\mathbf{r} \times K$ matrix (where \#r, the dimension of $\mathbf{r}$, is the number of restrictions) of full row rank,


\begin{equation*}
\left.W \equiv n \cdot(\mathbf{R} \mathbf{b}-\mathbf{r})^{\prime}\{\mathbf{R}[\widehat{\operatorname{Avar}(\mathbf{b}})] \mathbf{R}^{\prime}\right\}^{-1}(\mathbf{R} \mathbf{b}-\mathbf{r}) \underset{\mathrm{d}}{\rightarrow} \chi^{2}(\# \mathbf{r}) . \tag{2.4.2}
\end{equation*}


What remains to be shown is that $W \rightarrow_{\mathrm{d}} \chi^{2}(\# \mathbf{r})$, which is a straightforward application of Lemma 2.4(d).

Proof (continued). Write $W$ as

$$
W=\mathbf{c}_{n}^{\prime} \mathbf{Q}_{n}^{-1} \mathbf{c}_{n} \quad \text { where } \quad \mathbf{c}_{n} \equiv \sqrt{n}(\mathbf{R} \mathbf{b}-\mathbf{r}) \text { and } \mathbf{Q}_{n} \equiv \mathbf{R} \widehat{\operatorname{Avar}(\mathbf{b})} \mathbf{R}^{\prime} .
$$

Under $\mathrm{H}_{0}, \mathbf{c}_{n}=\mathbf{R} \sqrt{n}(\mathbf{b}-\boldsymbol{\beta})$. So by Proposition 2.1,

$$
\mathbf{c}_{n} \rightarrow \mathbf{d} \quad \text { where } \quad \mathbf{c} \sim N\left(\mathbf{0}, \mathbf{R} \text { Avar(b) } \mathbf{R}^{\prime}\right) .
$$

Also by Proposition 2.1,

$$
\mathbf{Q}_{n} \underset{\mathbf{p}}{\rightarrow} \mathbf{Q} \text { where } \mathbf{Q} \equiv \mathbf{R} \operatorname{Avar}(\mathbf{b}) \mathbf{R}^{\prime} .
$$

Because $\mathbf{R}$ is of full row rank and $\operatorname{Avar}(\mathbf{b})$ is positive definite, $\mathbf{Q}$ is invertible. Therefore, by Lemma 2.4(d),

$$
W \rightarrow \underset{\mathrm{~d}}{\rightarrow} \mathbf{c}^{\prime} \mathbf{Q}^{-1} \mathbf{c} .
$$

Since the \#r-dimensional random vector $\mathbf{c}$ is normally distributed and since $\mathbf{Q}$ equals $\operatorname{Var}(\mathbf{c}), \mathbf{c}^{\prime} \mathbf{Q}^{-1} \mathbf{c} \sim \chi^{2}(\# \mathbf{r})$. \(\square\)

This chi-square statistic $W$ is a Wald statistic because it is based on unrestricted estimates (b and $\widehat{\operatorname{Avar}(\mathbf{b})}$ here) not constrained by the null hypothesis $\mathrm{H}_{0}$. Testing $\mathrm{H}_{0}$ at a significance level of $\alpha$ proceeds as follows.

Step 1: Calculate the $W$ statistic by the formula (2.4.2).\\
Step 2: Look up the table of $\chi^{2}$ (\#r) distribution to find the critical value $\chi_{\alpha}^{2}(\# \mathrm{r})$ that gives $\alpha$ to the upper tail of the $\chi^{2}(\# \mathbf{r})$ distribution.\\
Step 3: If $W<\chi_{\alpha}^{2}(\# \mathrm{r})$, then accept $\mathrm{H}_{0}$; otherwise reject.\\
The probability of Type I error approaches $\alpha$ as the sample becomes larger. As will be made clear in Section 2.6, this Wald statistic is closely related to the familiar $F$-test under conditional homoskedasticity.

\section*{The Test Is Consistent}
Recall from basic statistics that the (finite-sample) power of a test is the probability of rejecting the null hypothesis when it is false given a sample of finite size (that is, the power is 1 minus the probability of Type II error). Power will obviously depend on the DGP (i.e., how the data were actually generated) considered as the alternative as well as on the size (significance level) of the test. For example, consider any DGP $\left\{y_{i}, \mathbf{x}_{i}\right\}$ satisfying Assumptions 2.1-2.5 but not the null hypothesis\\
$\mathrm{H}_{0}: \beta_{k}=\bar{\beta}_{k}$. The power of the $t$-test of size $\alpha$ against this alternative is

$$
\text { power }=\operatorname{Prob}\left(\left|t_{k}\right|>t_{\alpha / 2}\right)
$$

which depends on the DGP in question because the DGP controls the distribution of $t_{k}$. We say that a test is consistent against a set of DGPs, none of which satisfies the null, if the power against any particular member of the set approaches unity as $n \rightarrow \infty$ for any assumed significance level.

That the $t$-test is consistent against the set of alternatives (DGPs) satisfying Assumptions 2.1-2.5 can be seen as follows. Look at the expression (2.4.1) for the $t$-ratio, reproduced here:

$$
t_{k} \equiv \frac{\sqrt{n}\left(b_{k}-\bar{\beta}_{k}\right)}{\sqrt{\widehat{\operatorname{Avar}\left(b_{k}\right)}}}
$$

The denominator converges to $\sqrt{\operatorname{Avar}\left(b_{k}\right)}$ despite the fact that the DGP does not satisfy the null (recall that all parts of Proposition 2.1 hold regardless of the truth of the null, provided Assumptions 2.1-2.5 are satisfied). On the other hand, the numerator tends to $+\infty$ or $-\infty$ because $b_{k}$ converges in probability to the DGP's $\beta_{k}$, which is different from $\bar{\beta}_{k}$. So the power tends to unity as the sample size $n$ tends to infinity, implying that the $t$-test of Proposition 2.3 is consistent against those alternatives, the DGPs that do not satisfy the null. The same is true for the Wald test.

\section*{Asymptotic Power}
For later use in the next chapter, we define here the asymptotic power of a consistent test. As noted above, the power of the $t$-test approaches to unity as the sample size increases while the DGP taken as the alternative is held fixed. But if the DGP gets closer and closer to the null as the sample size increases, the power may not converge to unity. A sequence of such DGPs is called a sequence of local alternatives. For the regression model and for the null of $\mathrm{H}_{0}: \beta_{k}=\bar{\beta}_{k}$, it is a sequence of DGPs such that (i) the $n$-th DGP, $\left\{y_{i}^{(n)}, \mathbf{x}_{i}^{(n)}\right\}(i=1,2, \ldots)$, satisfies Assumptions 2.1-2.5 and converges in a certain sense to a fixed DGP $\left\{y_{i}, \mathbf{x}_{i}\right\},{ }^{12}$ and (ii) the value of $\beta_{k}$ of the $n$-th DGP, $\beta_{k}^{(n)}$, converges to $\bar{\beta}_{k}$. Suppose, further, that $\beta_{k}^{(n)}$ satisfies


\begin{equation*}
\beta_{k}^{(n)}=\bar{\beta}_{k}+\frac{\gamma}{\sqrt{n}} \tag{2.4.3}
\end{equation*}


\footnotetext{${ }^{12}$ See, e.g., Assumption 1 of Newey (1985) for a precise statement.
}
for some given $\gamma \neq 0$. So $\beta_{k}^{(n)}$ approaches to $\bar{\beta}_{k}$ at a rate proportional to $1 / \sqrt{n}$. This special sequence of local alternatives is called a Pitman drift or a Pitman sequence. Substituting (2.4.3) into (2.4.1), the $t$-ratio above can be rewritten as

\begin{equation*}
t_{k}=\frac{\sqrt{n}\left(b_{k}-\beta_{k}^{(n)}\right)}{\sqrt{\widehat{\operatorname{Avar}\left(b_{k}\right)}}}+\frac{\gamma}{\sqrt{\widehat{\operatorname{Avar}\left(b_{k}\right)}}} \tag{2.4.4}
\end{equation*}


If the sample of size $n$ is generated by the $n$-th DGP of a Pitman drift, does $t_{k}$ converge to a nontrivial distribution? Since the $n$-th DGP satisfies Assumptions 2.1-2.5, the first term on the right hand side of (2.4.4) converges in distribution to $N(0,1)$ by parts (b) and (c) of Proposition 2.1. By part (c) of Proposition 2.1 and the fact that $\left\{y_{i}^{(n)}, \mathbf{x}_{i}^{(n)}\right\}$ "converges" to a fixed DGP, the second term converges in probability to


\begin{equation*}
\mu \equiv \frac{\gamma}{\sqrt{\operatorname{Avar}\left(b_{k}\right)}} \tag{2.4.5}
\end{equation*}


where $\operatorname{Avar}\left(b_{k}\right)$ is evaluated at the fixed DGP. Therefore, $t_{k} \rightarrow_{\mathrm{d}} N(\mu, 1)$ along this sequence of local alternatives. If the significance level is $\alpha$, the power converges to


\begin{equation*}
\operatorname{Prob}\left(|x|>t_{\alpha / 2}\right) \tag{2.4.6}
\end{equation*}


where $x \sim N(\mu, 1)$ and $t_{\alpha / 2}$ is the level- $\alpha$ critical value. This probability is called the asymptotic power. It is a measure of the ability of the test to detect small deviations of the model from the null hypothesis. Evidently, the larger is $|\mu|$, the higher is the asymptotic power for any given size $\alpha$. By a similar argument, it is easy to show that the Wald statistic converges to a distribution called noncentral chi-squared.

\section*{Testing Nonlinear Hypotheses}
The Wald statistic can be generalized to a test of a set of nonlinear restrictions on $\boldsymbol{\beta}$. Consider a null hypothesis of the form

$$
\mathrm{H}_{0}: \mathbf{a}(\boldsymbol{\beta})=\mathbf{0} .
$$

Here, $\mathbf{a}$ is a vector-valued function with continuous first derivatives. Let \#a be the dimension of $\mathbf{a}(\boldsymbol{\beta})$ (so the null hypothesis has \#a restrictions), and let $\mathbf{A}(\boldsymbol{\beta})$ be the $\# \mathbf{a} \times K$ matrix of first derivatives evaluated at $\boldsymbol{\beta}: \mathbf{A}(\boldsymbol{\beta})=\partial \mathbf{a}(\boldsymbol{\beta}) / \partial \boldsymbol{\beta}^{\prime}$. For the hypothesis to be well-defined, we assume that $\mathbf{A}(\boldsymbol{\beta})$ is of full row rank (this is the generalization of the requirement for linear hypothesis $\mathbf{R} \boldsymbol{\beta}=\mathbf{r}$ that $\mathbf{R}$ is of full\\
row rank). Lemma 2.5 of Section 2.1 and Proposition 2.1(b) imply that


\begin{equation*}
\sqrt{n}[\mathbf{a}(\mathbf{b})-\mathbf{a}(\boldsymbol{\beta})] \underset{\mathrm{d}}{\rightarrow} \mathbf{c}, \mathbf{c} \sim N\left(\mathbf{0}, \mathbf{A}(\boldsymbol{\beta}) \operatorname{Avar}(\mathbf{b}) \mathbf{A}(\boldsymbol{\beta})^{\prime}\right) . \tag{2.4.7}
\end{equation*}


Since $\mathbf{a}(\boldsymbol{\beta})=\mathbf{0}$ under $\mathrm{H}_{0}$, (2.4.7) becomes


\begin{equation*}
\sqrt{n} \mathbf{a}(\mathbf{b}) \underset{\mathrm{d}}{\rightarrow} \mathbf{c}, \mathbf{c} \sim N\left(\mathbf{0}, \mathbf{A}(\boldsymbol{\beta}) \operatorname{Avar}(\mathbf{b}) \mathbf{A}(\boldsymbol{\beta})^{\prime}\right) . \tag{2.4.8}
\end{equation*}


Since $\mathbf{b} \rightarrow_{\mathrm{p}} \boldsymbol{\beta}$ by Proposition 2.1(a), Lemma 2.3(a) implies that $\mathbf{A}(\mathbf{b}) \rightarrow_{\mathrm{p}} \mathbf{A}(\boldsymbol{\beta})$. By Proposition 2.1(c), $\widehat{\operatorname{Avar}(\mathbf{b})} \rightarrow_{\mathrm{p}} \operatorname{Avar}(\mathbf{b})$. So by Lemma 2.3(a),


\begin{equation*}
\mathbf{A}(\mathbf{b}) \widehat{\operatorname{Avar}(\mathbf{b})} \mathbf{A}(\mathbf{b})^{\prime} \underset{\mathrm{p}}{\rightarrow} \mathbf{A}(\boldsymbol{\beta}) \operatorname{Avar}(\mathbf{b}) \mathbf{A}(\boldsymbol{\beta})^{\prime}=\operatorname{Var}(\mathbf{c}) \tag{2.4.9}
\end{equation*}


Because $\mathbf{A}(\boldsymbol{\beta})$ is of full row rank and $\operatorname{Avar}(\mathbf{b})$ is positive definite, $\operatorname{Var}(\mathbf{c})$ is invertible. Then Lemma 2.4(d), (2.4.8), and (2.4.9) imply


\begin{equation*}
\sqrt{n} \mathbf{a}(\mathbf{b})^{\prime}\left\{\mathbf{A}(\mathbf{b}) \widehat{\operatorname{Avar}(\mathbf{b})} \mathbf{A}(\mathbf{b})^{\prime}\right\}^{-1} \sqrt{n} \mathbf{a}(\mathbf{b}) \underset{\mathrm{d}}{\rightarrow} \mathbf{c}^{\prime} \operatorname{Var}(\mathbf{c})^{-1} \mathbf{c} \sim \chi^{2}(\# \mathbf{a}) . \tag{2.4.10}
\end{equation*}


Combining two $\sqrt{n}$ 's in (2.4.10) into one $n$, we have proved

\section*{Proposition 2.3 (continued):}
(c) Under the null hypothesis with \#a restrictions $\mathrm{H}_{0}: \mathbf{a}(\boldsymbol{\beta})=\mathbf{0}$ such that $\mathbf{A}(\boldsymbol{\beta})$, the \#a × K matrix of continuous first derivatives of $\mathbf{a}(\boldsymbol{\beta})$, is of full row rank, we have


\begin{equation*}
W \equiv n \cdot \mathbf{a}(\mathbf{b})^{\prime}\left\{\mathbf{A}(\mathbf{b}) \widehat{\operatorname{Avar}(\mathbf{b})} \mathbf{A}(\mathbf{b})^{\prime}\right\}^{-1} \mathbf{a}(\mathbf{b}) \underset{\mathrm{d}}{\rightarrow} \chi^{2}(\# \mathbf{a}) . \tag{2.4.11}
\end{equation*}


Part (c) is a generalization of (b); by setting $\mathbf{a}(\boldsymbol{\beta})=\mathbf{R} \boldsymbol{\beta}-\mathbf{r}$, (2.4.11) reduces to (2.4.2), the Wald statistic for linear restrictions.

The choice of $\mathbf{a}(\cdot)$ for representing a given set of restrictions is not unique. For example, $\beta_{1} \beta_{2}=1$ can be written as $a(\boldsymbol{\beta})=0$ with $a(\boldsymbol{\beta})=\beta_{1} \beta_{2}-1$ or with $a(\boldsymbol{\beta})=\beta_{1}-1 / \beta_{2}$. While part (c) of the proposition guarantees that in large samples the outcome of the Wald test is the same regardless of the choice of the function a, the numerical value of the Wald statistic $W$ does depend on the representation, and the test outcome can be different in finite samples. In the above example, the second representation, $a(\boldsymbol{\beta})=\beta_{1}-1 / \beta_{2}$, does not satisfy the requirement of continuous derivatives at $\beta_{2}=0$. Indeed, a Monte Carlo study by Gregory and Veall (1985) reports that, when $\beta_{2}$ is close to zero, the Wald test based on the second representation rejects the null too often in small samples.

\section*{QUESTIONS FOR REVIEW}
\begin{enumerate}
  \item Does $S E^{*}\left(b_{k}\right) \rightarrow_{\mathrm{p}} 0$ as $n \rightarrow \infty$ ?
  \item (Standard error of a nonlinear function) For simplicity let $K=1$ and let $b$ be the OLS estimate of $\beta$. The standard error of $b$ is $\sqrt{\widehat{\operatorname{Avar}(b)} / n}$. Suppose $\lambda= -\log (\beta)$. The estimate of $\lambda$ implied by the OLS estimate of $\beta$ is $\hat{\lambda}=-\log (b)$. Verify that the standard error of $\hat{\lambda}$ is $(1 / b) \cdot \sqrt{\widehat{\operatorname{Avar}(b)} / n}$.
  \item (Invariance [or lack thereof] of the Wald statistic) There is no unique way to write the linear hypothesis $\mathbf{R} \boldsymbol{\beta}=\mathbf{r}$, because for any $\# \mathbf{r} \times \# \mathbf{r}$ nonsingular matrix $\mathbf{F}$, the same set of restrictions can be represented as $\widetilde{\mathbf{R}} \boldsymbol{\beta}=\widetilde{\mathbf{r}}$ with $\widetilde{\mathbf{R}} \equiv \mathbf{F R}$ and $\widetilde{\mathbf{r}} \equiv \mathbf{F r}$. Does a different choice of $\mathbf{R}$ and $\mathbf{r}$ affect the asymptotic distribution of $W$ ? The finite-sample distribution? The numerical value?
\end{enumerate}

\subsection*{2.5 Estimating $\mathrm{E}\left(\varepsilon_{i}^{2} \mathrm{x}_{i} \mathrm{x}_{i}^{\prime}\right)$ Consistently}
The theory developed so far presumes that there is available a consistent estimator, $\widehat{\mathbf{S}}$, of $\mathbf{S}\left(=\mathrm{E}\left(\mathbf{g}_{i} \mathbf{g}_{i}^{\prime}\right) \equiv \mathrm{E}\left(\varepsilon_{i}^{2} \mathbf{x}_{i} \mathbf{x}_{i}^{\prime}\right)\right)$ to be used to calculate the estimated asymptotic variance, $\widehat{\operatorname{Avar}(\mathbf{b})}$. This section explains how to obtain $\widehat{\mathbf{S}}$ from the sample $(\mathbf{y}, \mathbf{X})$.

\section*{Using Residuals for the Errors}
If the error were observable, then the sample mean of $\varepsilon_{i}^{2} \mathbf{x}_{i} \mathbf{x}_{i}^{\prime}$ is obviously consistent by ergodic stationarity. But we do not observe the error term, and the substitution of some consistent estimate of it results in


\begin{equation*}
\widehat{\mathbf{S}} \equiv \frac{1}{n} \sum_{i=1}^{n} \hat{\varepsilon}_{i}^{2} \mathbf{x}_{i} \mathbf{x}_{i}^{\prime}, \tag{2.5.1}
\end{equation*}


where $\hat{\varepsilon}_{i} \equiv y_{i}-\mathbf{x}_{i}^{\prime} \widehat{\boldsymbol{\beta}}$, and $\widehat{\boldsymbol{\beta}}$ is some consistent estimator of $\boldsymbol{\beta}$. (Although the obvious candidate for the consistent estimator $\widehat{\boldsymbol{\beta}}$ is the OLS estimator $\mathbf{b}$, we use $\widehat{\boldsymbol{\beta}}$ rather than $\mathbf{b}$ here, in order to make the point that the results of this section hold for any consistent estimator.) For this estimator to be consistent for $\mathbf{S}$, we need to make a fourth-moment assumption about the regressors.

Assumption 2.6 (finite fourth moments for regressors): $\mathrm{E}\left[\left(x_{i k} x_{i j}\right)^{2}\right]$ exists and is finite for all $k, j(=1,2, \ldots, K)$.

Proposition 2.4 (consistent estimation of S): Suppose the coefficient estimate $\widehat{\boldsymbol{\beta}}$ used for calculating the residual $\hat{\varepsilon}_{i}$ for $\widehat{\mathbf{S}}$ in (2.5.1) is consistent, and suppose $\mathbf{S}=\mathrm{E}\left(\mathbf{g}_{i} \mathbf{g}_{i}^{\prime}\right)$ exists and is finite. Then, under Assumptions 2.1, 2.2, and 2.6, $\widehat{\mathbf{S}}$ given in (2.5.1) is consistent for $\mathbf{S}$.

To indicate why the fourth-moment assumption is needed for the regressors, we provide a sketch of the proof for the special case of $K=1$ (only one regressor). So $\mathbf{x}_{i}$ is now a scalar $x_{i}, \mathbf{g}_{i}$ is a scalar $g_{i}=x_{i} \varepsilon_{i}$, and (2.3.10) (with $\mathbf{b}=\widehat{\boldsymbol{\beta}}$ and $\left.e_{i}=\hat{\varepsilon}_{i}\right)$ simplifies to


\begin{equation*}
\hat{\varepsilon}_{i}^{2}=\varepsilon_{i}^{2}-2(\widehat{\beta}-\beta) x_{i} \varepsilon_{i}+(\widehat{\beta}-\beta)^{2} x_{i}^{2} . \tag{2.5.2}
\end{equation*}


By multiplying both sides by $x_{i}^{2}$ and summing over $i$,


\begin{equation*}
\frac{1}{n} \sum_{i=1}^{n} \hat{\varepsilon}_{i}^{2} x_{i}^{2}-\frac{1}{n} \sum_{i=1}^{n} \varepsilon_{i}^{2} x_{i}^{2}=-2(\widehat{\beta}-\beta) \frac{1}{n} \sum_{i=1}^{n} \varepsilon_{i} x_{i}^{3}+(\widehat{\beta}-\beta)^{2} \frac{1}{n} \sum_{i=1}^{n} x_{i}^{4} . \tag{2.5.3}
\end{equation*}


Now we can see why the finite fourth-moment assumption on $x_{i}$ is required: if the fourth moment $\mathrm{E}\left(x_{i}^{4}\right)$ is finite, then by ergodic stationarity the sample average of $x_{i}^{4}$ converges in probability to some finite number, so that the last term in (2.5.3) vanishes (converges to 0 in probability) if $\widehat{\beta}$ is consistent for $\beta$. It can also be shown (see Analytical Exercise 4 for proof) that, by combining the same fourthmoment assumption about the regressors and the fourth-moment assumption that $\mathrm{E}\left(\mathbf{g}_{i} \mathbf{g}_{i}^{\prime}\right)\left(=\mathrm{E}\left(\varepsilon_{i}^{2} \mathbf{x}_{i} \mathbf{x}_{i}^{\prime}\right)\right)$ is finite, the sample average of $x_{i}^{3} \varepsilon_{i}$ converges in probability to some finite number, so that the other term on the RHS of (2.5.3), too, vanishes.

According to Proposition 2.1(a), the assumptions made in Proposition 2.3 are sufficient to guarantee that $\mathbf{b}$ is consistent, so we can set $\mathbf{b}=\widehat{\boldsymbol{\beta}}$ in (2.5.1) and use the OLS residual to calculate $\widehat{\mathbf{S}}$. Also, the assumption made in Proposition 2.4 that $\mathrm{E}\left(\mathbf{g}_{i} \mathbf{g}_{i}^{\prime}\right)$ is finite is part of Assumption 2.5, which is assumed in Proposition 2.3. Therefore, the import of Proposition 2.4 is:

If Assumption 2.6 is added to the hypothesis in Proposition 2.3, then the $\widehat{\mathbf{S}}$ given in (2.5.1) with $\mathbf{b}=\widehat{\boldsymbol{\beta}}$ (so $\hat{\varepsilon}_{i}$ is the OLS residual $e_{i}$ ) can be used in (2.3.5) to calculate the estimated asymptotic variance:


\begin{equation*}
\widehat{\operatorname{Avar}(\mathbf{b})}=\mathbf{S}_{\mathbf{x x}}^{-1}\left(\frac{1}{n} \sum_{i=1}^{n} e_{i}^{2} \mathbf{x}_{i} \mathbf{x}_{i}^{\prime}\right) \mathbf{S}_{\mathbf{x x}}^{-1} \tag{2.5.4}
\end{equation*}


which is consistent for Avar(b).

\section*{Data Matrix Representation of $\mathbf{S}$}
If $\mathbf{B}$ is the $n \times n$ diagonal matrix whose $i$-th diagonal element is $\hat{\varepsilon}_{i}^{2}$, then the $\widehat{\mathbf{S}}$ in (2.5.1) can be represented in terms of data matrices as


\begin{equation*}
\widehat{\mathbf{S}}=\frac{\mathbf{X}^{\prime} \mathbf{B X}}{n} \tag{$\prime$}
\end{equation*}


with

$$
\mathbf{B}=\left[\begin{array}{ccc}
\hat{\varepsilon}_{1}^{2} & & \\
& \ddots & \\
& & \hat{\varepsilon}_{n}^{2}
\end{array}\right] .
$$

So (2.5.4) can be rewritten (with $\hat{\varepsilon}_{i}$ in $\mathbf{B}$ set to $e_{i}$ ) as


\begin{equation*}
\widehat{\operatorname{Avar}(\mathbf{b})}=n \cdot\left(\mathbf{X}^{\prime} \mathbf{X}\right)^{-1}\left(\mathbf{X}^{\prime} \mathbf{B X}\right)\left(\mathbf{X}^{\prime} \mathbf{X}\right)^{-1} . \tag{$\prime$}
\end{equation*}


These expressions, although useful for some purposes, should not be used for computation purposes, because the $n \times n$ matrix $\mathbf{B}$ will take up too much of the computer's working memory, particularly when the sample size is large. To compute $\widehat{\mathbf{S}}$ from the sample, the formula (2.5.1) is more useful than (2.5.1').

\section*{Finite-Sample Considerations}
Of course, in finite samples, the power may well be far below one against certain alternatives. Also, the probability of rejecting the null when the DGP does satisfy the null (the Type I error) may be very different from the assumed significance level. Davidson and MacKinnon (1993, Section 16.3) report that, at least for the Monte Carlo simulations they have seen, the robust $t$-ratio based on (2.5.1) rejects the null too often and that simply replacing the denominator $n$ in (2.5.1) by the degrees of freedom $n-K$ or equivalently multiplying (2.5.1) by $n /(n-K)$ (this is a degrees of freedom correction) mitigates the problem of overrejecting. They also report that the robust $t$-ratios based on the following adjustments on $\widehat{\mathbf{S}}$ perform even better:


\begin{equation*}
\widehat{\mathbf{S}}=\frac{1}{n} \sum_{i=1}^{n} \frac{e_{i}^{2}}{\left(1-p_{i}\right)^{d}} \mathbf{x}_{i} \mathbf{x}_{i}^{\prime}, \quad d=1 \text { or } 2, \tag{2.5.5}
\end{equation*}


where $p_{i}$ is the $p_{i}$ defined in the context of the influential analysis in Chapter 1: it is the $i$-th diagonal element of the projection matrix $\mathbf{P}$, that is,

$$
p_{i} \equiv \mathbf{x}_{i}^{\prime}\left(\mathbf{X}^{\prime} \mathbf{X}\right)^{-1} \mathbf{x}_{i}=\frac{\mathbf{x}_{i}^{\prime} \mathbf{S}_{\mathbf{x x}}^{-1} \mathbf{x}_{i}}{n}
$$

\section*{QUESTIONS FOR REVIEW}
\begin{enumerate}
  \item (Computation of robust standard errors) In Review Question 9 of Section 1.2, we observed that the standard errors of the OLS coefficient estimates can be calculated from $\mathbf{S}_{\mathbf{x x}}, \mathbf{S}_{\mathbf{x y}}$ (the sample mean of $\left.\mathbf{x}_{i} \cdot y_{i}\right), \mathbf{y}^{\prime} \mathbf{y} / n$, and $\bar{y}$, so the sample moments need to be computed just once. Is the same true for the robust standard errors where the $\widehat{\mathbf{S}}$ is calculated according to the formula (2.5.1) with $\hat{\varepsilon}_{i}=e_{i}$ ?
  \item The finite-sample variance of the OLS estimator in the generalized regression model of Chapter 1 is $\operatorname{Var}(\mathbf{b} \mid \mathbf{X})=\left(\mathbf{X}^{\prime} \mathbf{X}\right)^{-1} \mathbf{X}^{\prime}\left(\sigma^{2} \mathbf{V}\right) \mathbf{X}\left(\mathbf{X}^{\prime} \mathbf{X}\right)^{-1}$. Compare this to (2.5.4'). What are the differences?
\end{enumerate}

\subsection*{2.6 Implications of Conditional Homoskedasticity}
The test statistics developed in Sections 2.4 and 2.5 are different from the finitesample counterparts of Section 1.4 designed for testing the same null hypothesis. How are they related? What is the asymptotic distribution of the $t$ and $F$ statistics of Chapter 1? This section answers these questions.

It turns out that the robust $t$-ratio is numerically equal to the $t$-ratio of Section 1.4, for a particular choice of $\widehat{\mathbf{S}}$. Therefore, the asymptotic distribution of the $t$ ratio of Section 1.4 is the same as that of the robust $t$-ratio, if that particular choice is consistent for $\mathbf{S}$. The same relationship holds between the $F$-ratio of Section 1.4 and the Wald statistic $W$ of this chapter. Under the conditional homoskedasticity assumption stated below, that particular choice is indeed consistent.

\section*{Conditional versus Unconditional Homoskedasticity}
The conditional homoskedasticity assumption is:

\section*{Assumption 2.7 (conditional homoskedasticity):}

\begin{equation*}
\mathrm{E}\left(\varepsilon_{i}^{2} \mid \mathbf{x}_{i}\right)=\sigma^{2}>0 . \tag{2.6.1}
\end{equation*}


This assumption implies that the unconditional second moment $\mathrm{E}\left(\varepsilon_{i}^{2}\right)$ equals $\sigma^{2}$ by the Law of Total Expectations. To be clear about the distinction between unconditional and conditional homoskedasticity, consider the following example.

Example 2.6 (unconditionally homoskedastic but conditionally heteroskedastic errors): As already observed, if $\left\{y_{i}, \mathbf{x}_{i}\right\}$ is stationary, so is $\left\{\varepsilon_{i}\right\}$, and the error is unconditionally homoskedastic in that $\mathrm{E}\left(\varepsilon_{i}^{2}\right)$ does not depend on $i$. To illustrate that the error can nevertheless be conditionally heteroskedastic, suppose that $\varepsilon_{i}$ is written as $\varepsilon_{i} \equiv \eta_{i} f\left(\mathbf{x}_{i}\right)$, where $\left\{\eta_{i}\right\}$ is zero-mean $\mathrm{E}\left(\eta_{i}\right)=0$ and is independent of $\mathbf{x}_{i}$. The conditional second moment of $\varepsilon_{i}$ depends on $\mathbf{x}_{i}$ because

$$
\begin{aligned}
& \left.\mathrm{E}\left(\varepsilon_{i}^{2} \mid \mathbf{x}_{i}\right)=\mathrm{E}\left(\eta_{i}^{2} f\left(\mathbf{x}_{i}\right)^{2} \mid \mathbf{x}_{i}\right) \quad \text { (since } \varepsilon_{i} \equiv \eta_{i} f\left(\mathbf{x}_{i}\right)\right) \\
& =f\left(\mathbf{x}_{i}\right)^{2} \mathrm{E}\left(\eta_{i}^{2} \mid \mathbf{x}_{i}\right) \quad \text { (by the linearity of conditional expectations) } \\
& =f\left(\mathbf{x}_{i}\right)^{2} \mathrm{E}\left(\eta_{i}^{2}\right) \quad \text { (since } \eta_{i} \text { is independent of } \mathbf{x}_{i} \text { by assumption) }
\end{aligned}
$$

which varies across $i$ because of the variation in $f\left(\mathbf{x}_{i}\right)$ across $i$.

\section*{Reduction to Finite-Sample Formulas}
To examine the large-sample distribution of the $t$ and $F$ statistics of Chapter 1 under the additional Assumption 2.7, we first examine the algebraic relationship to their robust counterparts. Consider the following choice for the estimate of $\mathbf{S}$ :

$$
\widehat{\mathbf{S}}=s^{2} \mathbf{S}_{\mathbf{x x}},
$$

where $s^{2}$ is the OLS estimate of $\sigma^{2}$. (We will show in a moment that this estimator is consistent under conditional homoskedasticity.) Then the expression (2.3.5) for $\widehat{\text { Avar(b) }}$ becomes


\begin{equation*}
\widehat{\operatorname{Avar}(\mathbf{b})}=s^{2} \mathbf{S}_{\mathbf{x x}}^{-1}=n \cdot s^{2} \cdot\left(\mathbf{X}^{\prime} \mathbf{X}\right)^{-1} . \tag{2.6.2}
\end{equation*}


Substituting this expression into (2.4.1), we see that the robust standard error becomes


\begin{equation*}
\sqrt{s^{2} \text { times }(k, k) \text { element of }\left(\mathbf{X}^{\prime} \mathbf{X}\right)^{-1}} \tag{2.6.3}
\end{equation*}


which is the usual standard error in finite-sample theory. So the robust $t$-ratio is numerically identical to the usual finite-sample $t$-ratio when we set $\widehat{\mathbf{S}}=s^{2} \mathbf{S}_{\mathbf{x x}}$. Similarly, substituting (2.6.2) into the expression for the Wald statistic (2.4.2), we obtain

$$
\begin{aligned}
W & =n \cdot(\mathbf{R} \mathbf{b}-\mathbf{r})^{\prime}\left\{\mathbf{R}\left[n \cdot s^{2} \cdot\left(\mathbf{X}^{\prime} \mathbf{X}\right)^{-1}\right] \mathbf{R}^{\prime}\right\}^{-1}(\mathbf{R} \mathbf{b}-\mathbf{r}) \\
& =(\mathbf{R} \mathbf{b}-\mathbf{r})^{\prime}\left\{\mathbf{R}\left[s^{2} \cdot\left(\mathbf{X}^{\prime} \mathbf{X}\right)^{-1}\right] \mathbf{R}^{\prime}\right\}^{-1}(\mathbf{R} \mathbf{b}-\mathbf{r}) \quad \text { (the two } n \text { 's cancel) } \\
& =(\mathbf{R} \mathbf{b}-\mathbf{r})^{\prime}\left\{\mathbf{R}\left[\left(\mathbf{X}^{\prime} \mathbf{X}\right)^{-1}\right] \mathbf{R}^{\prime}\right\}^{-1}(\mathbf{R} \mathbf{b}-\mathbf{r}) / s^{2} \\
& =r \cdot F \quad(\text { by the definition of (1.4.9) of the } F \text {-ratio) } \\
& =\left(S S R_{R}-S S R_{U}\right) / s^{2} \quad(\text { by }(1.4 .11))
\end{aligned}
$$

Thus, when we set $\widehat{\mathbf{S}}=s^{2} \mathbf{S}_{\mathbf{x x}}$, the Wald statistic $W$ is numerically identical to $\# \mathbf{r} \cdot F$ (where \#r is the number of restrictions in the null hypothesis).

\section*{Large-Sample Distribution of $\boldsymbol{t}$ and $\boldsymbol{F}$ Statistics}
It then follows from Proposition 2.3 that the $t$-ratio (2.4.1) is asymptotically $N(0,1)$ and \# $\mathbf{r} \cdot F$ asymptotically $\chi^{2}(\# \mathbf{r})$, if $s^{2} \mathbf{S}_{\mathbf{x x}}$ is consistent for $\mathbf{S}$. That $s^{2} \mathbf{S}_{\mathbf{x x}}$ is consistent for $\mathbf{S}$ can be seen as follows. Under conditional homoskedasticity, the matrix of fourth moments $\mathbf{S}$ can be expressed as a product of second moments:


\begin{align*}
\mathbf{S} & =\mathrm{E}\left(\mathbf{g}_{i} \mathbf{g}_{i}^{\prime}\right)=\mathrm{E}\left(\mathbf{x}_{i} \mathbf{x}_{i}^{\prime} \varepsilon_{i}^{2}\right) \quad\left(\text { since } \mathbf{g}_{i} \equiv \mathbf{x}_{i} \cdot \varepsilon_{i}\right) \\
& =\mathrm{E}\left[\mathrm{E}\left(\mathbf{x}_{i} \mathbf{x}_{i}^{\prime} \varepsilon_{i}^{2} \mid \mathbf{x}_{i}\right)\right] \quad \text { (by the Law of Total Expectations) } \\
& =\mathrm{E}\left[\mathbf{x}_{i} \mathbf{x}_{i}^{\prime} \mathrm{E}\left(\varepsilon_{i}^{2} \mid \mathbf{x}_{i}\right)\right] \quad \text { (by the linearity of conditional expectations) } \\
& =\mathrm{E}\left(\mathbf{x}_{i} \mathbf{x}^{\prime} \sigma^{2}\right) \quad \text { (by Assumption 2.7) } \\
& =\sigma^{2} \mathrm{E}\left(\mathbf{x}_{i} \mathbf{x}_{i}^{\prime}\right)=\sigma^{2} \mathbf{\Sigma}_{\mathbf{x x}} \tag{2.6.4}
\end{align*}


This decomposition has several implications.

\begin{itemize}
  \item ( $\boldsymbol{\Sigma}_{\mathbf{x x}}$ is nonsingular) Since by Assumption $2.5 \mathbf{S}$ is nonsingular, this decomposition of $\mathbf{S}$ implies that $\sigma^{2}>0$ and $\boldsymbol{\Sigma}_{\mathbf{x x}}$ is nonsingular. Hence, Assumption 2.4 (rank condition) is implied.
  \item (No need for fourth-moment assumption) By ergodic stationarity $\mathbf{S}_{\mathbf{x x}} \rightarrow_{\mathrm{p}} \boldsymbol{\Sigma}_{\mathbf{x x}}$. By Proposition 2.2, $s^{2}$ is consistent for $\sigma^{2}$ under Assumptions 2.1-2.4. Thus, $s^{2} \mathbf{S}_{\mathbf{x x}} \rightarrow_{\mathrm{p}} \sigma^{2} \boldsymbol{\Sigma}_{\mathbf{x x}}=\mathbf{S}$. We do not need the fourth-moment assumption (Assumption 2.6) for consistency.
\end{itemize}

As another implication of (2.6.4), the expression for Avar(b) can be simplified: inserting (2.6.4) into (2.3.4) of Proposition 2.1, the expression for Avar(b) becomes


\begin{equation*}
\operatorname{Avar}(\mathbf{b})=\sigma^{2} \boldsymbol{\Sigma}_{\mathbf{x x}}^{-1} \tag{2.6.5}
\end{equation*}


Thus, we have proved

\section*{Proposition 2.5 (large-sample properties of $b, t$, and $F$ under conditional homoskedasticity): Suppose Assumptions 2.1-2.5 and 2.7 are satisfied. Then}
(a) (Asymptotic distribution of $\mathbf{b}$ ) The OLS estimator $\mathbf{b}$ of $\boldsymbol{\beta}$ is consistent and asymptotically normal with $\operatorname{Avar}(\mathbf{b})=\sigma^{2} \boldsymbol{\Sigma}_{\mathbf{x x}}^{-1}$.\\
(b) (Consistent estimation of asymptotic variance) Under the same set of assumptions, $\operatorname{Avar}(\mathbf{b})$ is consistently estimated by $\widehat{\operatorname{Avar}(\mathbf{b})}=s^{2} \mathbf{S}_{\mathbf{x x}}^{-1}=n \cdot s^{2} \cdot\left(\mathbf{X}^{\prime} \mathbf{X}\right)^{-1}$.\\
(c) (Asymptotic distribution of the $t$ and $F$ statistics of the finite-sample theory) Under $\mathrm{H}_{0}: \beta_{k}=\bar{\beta}_{k}$, the usual $t$-ratio (1.4.5) is asymptotically distributed as $N(0,1)$. Under $\mathrm{H}_{0}: \mathbf{R} \boldsymbol{\beta}=\mathbf{r}, \# \mathbf{r} \cdot F$ is asymptotically $\chi^{2}(\# \mathbf{r})$, where $F$ is the $F$ statistic from (1.4.9) and \# $\mathbf{r}$ is the number of restrictions in $\mathrm{H}_{0}$.

\section*{Variations of Asymptotic Tests under Conditional Homoskedasticity}
According to this result, you should look up the $N(0,1)$ table to find the critical value to be compared with the $t$-ratio (1.4.5) and the $\chi^{2}$ table for the statistic \# $\mathbf{r} \cdot F$ derived from (1.4.9). Some researchers replace the $s^{2}$ in (1.4.5) and (1.4.9) by $\frac{1}{n} \sum_{i} e_{i}^{2}$. That is, the degrees of freedom $n-K$ is replaced by $n$, or the degrees of freedom adjustment implicit in $s^{2}$ is removed. The difference this substitution makes vanishes in large samples, because $\lim _{n \rightarrow \infty}[n /(n-K)]=1$. Therefore, regardless of which test to use, the outcome of the test will be the same if the sample size is sufficiently large.

Another variation is to retain the degrees of freedom $n-K$ but use the $t(n- K$ ) table for the $t$ ratio and the $F(\# \mathbf{r}, n-K)$ table for $F$, which is exactly the prescription of finite-sample theory. This, too, is asymptotically valid because, as $n-K$ tends to infinity (which is what happens when $n \rightarrow \infty$ with $K$ fixed), the $t(n-K)$ distribution converges to $N(0,1)$ (just compare the $t$ table for large degrees of freedom with the standard normal table) and $F(\# \mathbf{r}, n-K)$ to $\chi^{2}(\# \mathbf{r}) / \# \mathbf{r}$. Put differently, even if the error is not normally distributed and the regressors are merely predetermined (orthogonal to the error term) and not strictly exogenous, the distribution of the $t$-ratio (1.4.5) is well approximated by $t(n-K)$, and the distribution of the $F$ ratio by $F(\# \mathrm{r}, n-K)$.

These variations are all asymptotically equivalent in that the differences in the values vanishes in large samples and hence (by Lemma 2.4(a)) their asymptotic distributions are the same. However, when the sample size is only moderately large, the approximation to the finite-sample or exact distribution of test statistics may be better with $t(n-K)$ and $F(\# \mathbf{r}, n-K)$, rather than with $N(0,1)$ and $\chi^{2}(\# \mathbf{r})$. Because the exact distribution depends on the DGP, there is no simple guide as to which variation works better in finite samples. This issue of which table- $N(0,1)$\\
or $t(n-K)$ - should be used for moderate sample sizes will be taken up in the Monte Carlo exercise of this chapter.

\section*{QUESTIONS FOR REVIEW}
\begin{enumerate}
  \item (Inconsistency of finite sample-formulas without conditional homoskedasticity) Without Assumption 2.7, Avar(b) is given by (2.3.4) in Proposition 2.1. Is it consistently estimated by (2.6.2) without Assumption 2.7? [The answer is no. Why?] Is the $t$-ratio (1.4.5) asymptotically standard normal without the assumption? [Answer: No.]
  \item (Advantage of finite-sample formulas under conditional homoskedasticity) Conversely, under Assumption 2.7, Avar(b) is given by (2.6.5). Is it consistently estimated by (2.3.5) under Assumption 2.7? If Assumption 2.7 holds, what do you think is the advantage of using (2.6.2) over (2.3.5) to estimate the asymptotic variance? [Note: The finite-sample properties of an estimator are generally better, the fewer the number of population parameters estimated to form the estimator. How many population parameters need to be estimated to form (2.3.5)? (2.6.2)?]
  \item (Relation of $F$ to $\chi^{2}$ ) Find the 5 percent critical value of $F(10, \infty)$ and compare it to the 5 percent critical value of $\chi^{2}(10)$. What is the relationship between the two?
  \item Without conditional homoskedasticity, is $\left(S S R_{R}-S S R_{U}\right) / s^{2}$ asymptotically $\chi^{2}(\# \mathbf{r})$ ? [Answer: No.]
  \item ( $n R^{2}$ test) For a regression with a constant, consider the null hypothesis that the coefficients of the $K-1$ nonconstant regressors are all zero. Show that $n R^{2} \rightarrow_{\mathrm{d}} \chi^{2}(K-1)$ under the hypothesis of Proposition 2.5. Hint: You have proved for a Chapter 1 analytical exercise that the algebraic relationship between the $F$-ratio for the null and $R^{2}$ is
\end{enumerate}

$$
F=\frac{R^{2} /(K-1)}{\left(1-R^{2}\right) /(n-K)} .
$$

Can you use the $n R^{2}$ statistic when the error is not conditionally homoskedastic? [Answer: No.]

\subsection*{2.7 Testing Conditional Homoskedasticity}
With the advent of robust standard errors allowing us to do inference without specifying the conditional second moment $\mathrm{E}\left(\varepsilon_{i}^{2} \mid \mathbf{x}_{i}\right)$, testing conditional homoskedasticity is not as important as it used to be. This section presents only the most popular test due to White (1980) for the case of random samples. ${ }^{13}$

Recall that $\widehat{\mathbf{S}}$ given in (2.5.1) with $\hat{\varepsilon}_{i}=e_{i}$ ) is consistent for $\mathbf{S}$ (Proposition 2.4), and $s^{2} \mathbf{S}_{\mathbf{x x}}$ is consistent for $\sigma^{2} \boldsymbol{\Sigma}_{\mathbf{x x}}$ (an implication of Proposition 2.2). But under conditional homoskedasticity, $\mathbf{S}=\sigma^{2} \boldsymbol{\Sigma}_{\mathbf{x x}}$ (see (2.6.4)), so the difference between the two consistent estimators should vanish:


\begin{align*}
\widehat{\mathbf{S}}-s^{2} \mathbf{S}_{\mathbf{x x}} & =\frac{1}{n} \sum_{i=1}^{n} e_{i}^{2} \mathbf{x}_{i} \mathbf{x}_{i}^{\prime}-s^{2} \frac{1}{n} \sum_{i=1}^{n} \mathbf{x}_{i} \mathbf{x}_{i}^{\prime} \\
& =\frac{1}{n} \sum_{i=1}^{n}\left(e_{i}^{2}-s^{2}\right) \mathbf{x}_{i} \mathbf{x}_{i}^{\prime} \underset{\mathrm{p}}{\rightarrow} \mathbf{0} . \tag{2.7.1}
\end{align*}


Let $\psi_{i}$ be a vector collecting unique and nonconstant elements of the $K \times K$ symmetric matrix $\mathbf{x}_{i} \mathbf{x}_{i}^{\prime}$. (Construction of $\psi_{i}$ from $\mathbf{x}_{i} \mathbf{x}_{i}^{\prime}$ will be illustrated in Example 2.7 below.) Then (2.7.1) implies


\begin{equation*}
\mathbf{c}_{n} \equiv \frac{1}{n} \sum_{i=1}^{n}\left(e_{i}^{2}-s^{2}\right) \boldsymbol{\psi}_{i} \underset{\mathrm{p}}{\rightarrow} \mathbf{0} . \tag{2.7.2}
\end{equation*}


This $\mathbf{c}_{n}$ is a sample mean converging to zero. Under some conditions appropriate for a Central Limit Theorem to be applicable, we would expect $\sqrt{n} \mathbf{c}_{n}$ to converge in probability to a normal distribution with mean zero and some asymptotic variance $\mathbf{B}$, so for any consistent estimator $\widehat{\mathbf{B}}$ of $\mathbf{B}$,


\begin{equation*}
n \cdot \mathbf{c}_{n}^{\prime} \widehat{\mathbf{B}}^{-1} \mathbf{c}_{n} \underset{\mathrm{~d}}{\rightarrow} \chi^{2}(m) \tag{2.7.3}
\end{equation*}


where $m$ is the dimension of $\mathbf{c}_{n}$. For a certain choice of $\widehat{\mathbf{B}}$, this statistic can be computed as $n R^{2}$ from the following auxiliary regression:


\begin{equation*}
\text { regress } e_{i}^{2} \text { on a constant and } \boldsymbol{\psi}_{i} \text {. } \tag{2.7.4}
\end{equation*}


${ }^{13}$ See, e.g., Judge et al. (1985, Section 11.3) or Greene (1997, Section 12.3) for other tests.

White (1980) rigorously developed this argument ${ }^{14}$ to prove

Proposition 2.6 (White's Test for Conditional Heteroskedasticity): In addition to Assumptions 2.1 and 2.4, suppose that (a) $\left\{y_{i}, \mathbf{x}_{i}\right\}$ is i.i.d. with finite $\mathrm{E}\left(\varepsilon_{i}^{2} \mathbf{x}_{i} \mathbf{x}_{i}^{\prime}\right)$ (thus strengthening Assumptions 2.2 and 2.5), (b) $\varepsilon_{i}$ is independent of $\mathbf{x}_{i}$ (thus strengthening Assumption 2.3 and conditional homoskedasticity), and (c) a certain condition holds on the moments of $\varepsilon_{i}$ and $\mathbf{x}_{i}$. Then,

$$
n R^{2} \underset{\mathrm{~d}}{\rightarrow} \chi^{2}(m),
$$

where $R^{2}$ is the $R^{2}$ from the auxiliary regression (2.7.4), and $m$ is the dimension of $\boldsymbol{\psi}_{i}$.

Example 2.7 (regressors in White's $\boldsymbol{n} \boldsymbol{R}^{\mathbf{2}}$ test): Consider the Cobb-Douglas cost function of Section 1.7:

$$
\log \left(\frac{T C_{i}}{p_{i 3}}\right)=\beta_{1}+\beta_{2} \log \left(Q_{i}\right)+\beta_{3} \log \left(\frac{p_{i 1}}{p_{i 3}}\right)+\beta_{4} \log \left(\frac{p_{i 2}}{p_{i 3}}\right)+\varepsilon_{i} .
$$

Here, $\mathbf{x}_{i}^{\prime}=\left(1, \log \left(Q_{i}\right), \log \left(p_{i 1} / p_{i 3}\right), \log \left(p_{i 2} / p_{i 3}\right)\right)$, a four-dimensional vector. There are $10(=4 \cdot 5 / 2)$ unique elements in $\mathbf{x}_{i} \mathbf{x}_{i}^{\prime}$ :

$$
\begin{gathered}
1, \log \left(Q_{i}\right), \log \left(\frac{p_{i 1}}{p_{i 3}}\right), \log \left(\frac{p_{i 2}}{p_{i 3}}\right), \\
{\left[\log \left(Q_{i}\right)\right]^{2}, \log \left(Q_{i}\right) \cdot \log \left(\frac{p_{i 1}}{p_{i 3}}\right), \log \left(Q_{i}\right) \cdot \log \left(\frac{p_{i 2}}{p_{i 3}}\right),} \\
{\left[\log \left(\frac{p_{i 1}}{p_{i 3}}\right)\right]^{2}, \log \left(\frac{p_{i 1}}{p_{i 3}}\right) \cdot \log \left(\frac{p_{i 2}}{p_{i 3}}\right),\left[\log \left(\frac{p_{i 2}}{p_{i 3}}\right)\right]^{2}}
\end{gathered}
$$

$\boldsymbol{\psi}_{i}$ is a nine-dimensional vector excluding a constant from this list. So the $m$ in Proposition 2.6 is 9.

If White's test accepts the null of conditional homoskedasticity, then the results of Section 2.6 apply, and statistical inference can be based on the $t$ - and $F$-ratios from Chapter 1. Otherwise inference should be based on the robust $t$ and Wald statistics of Proposition 2.3.

Because the regressors $\boldsymbol{\psi}_{i}$ in the auxiliary regression have many elements consisting of squares and cross-products of the elements of $\mathbf{x}_{i}$, the test should be consistent (i.e., the power approaches unity as $n \rightarrow \infty$ ) against most heteroskedas-

\footnotetext{${ }^{14}$ The original statement of White's theorem is more general as it covers the case where $\left\{y_{i}, \mathbf{x}_{i}\right\}$ is independent, but not identically distributed (i.n.i.d.).
}
tic alternatives but may require a fairly large sample to have power close to unity. If the researcher knows that some of the elements of $\boldsymbol{\psi}_{i}$ do not affect the conditional second moment, then they can be dropped from the auxiliary regression and the power might be increased in finite samples. The downside of it, of course, is that if such knowledge is false, the test will have no power against heteroskedastic alternatives that relate the conditional second moment to those elements that are excluded from the auxiliary regression.

\section*{QUESTION FOR REVIEW}
\begin{enumerate}
  \item (Dimension of $\boldsymbol{\psi}_{i}$ ) Suppose $\mathbf{x}_{i}=\left(1, q_{i}, q_{i}^{2}, p_{i}\right)^{\prime}$, a four-dimensional vector. How many nonconstant and unique elements are there in $\mathbf{x}_{i} \mathbf{x}_{i}^{\prime}$ ? [Answer: 8.]
\end{enumerate}

\subsection*{2.8 Estimation with Parameterized Conditional Heteroskedasticity (optional)}
Even when the error is found to be conditionally heteroskedastic, the OLS estimator is still consistent and asymptotically normal, and valid statistical inference can be conducted with robust standard errors and robust Wald statistics. However, in the (somewhat unlikely) case of a priori knowledge of the functional form of the conditional second moment $\mathrm{E}\left(\varepsilon_{i}^{2} \mid \mathbf{x}_{i}\right)$, it should be possible to obtain sharper estimates with smaller asymptotic variance. Indeed, in finite-sample theory, the WLS (weighted least squares) can incorporate such knowledge for increased efficiency in the sense of smaller finite-sample variance. Does this finite-sample result carry over to large-sample theory? This section deals with the large-sample properties of the WLS estimator. To simplify the discussion, throughout this section we strengthen Assumptions 2.2 and 2.5 by assuming that $\left\{y_{i}, \mathbf{x}_{i}\right\}$ is i.i.d. This is a natural assumption to make, because it is usually in cross-section contexts where WLS is invoked.

\section*{The Functional Form}
The parametric functional form for the conditional second moment we consider is


\begin{equation*}
\mathrm{E}\left(\varepsilon_{i}^{2} \mid \mathbf{x}_{i}\right)=\mathbf{z}_{i}^{\prime} \boldsymbol{\alpha}, \tag{2.8.1}
\end{equation*}


where $\mathbf{z}_{i}$ is a function of $\mathbf{x}_{i}$.

Example 2.8 (parametric form of the conditional second moment): The functional form used in the WLS estimation in the empirical exercise to Chapter 1 was

$$
\mathrm{E}\left(\varepsilon_{i}^{2} \mid \log \left(Q_{i}\right), \log \left(Q_{i}\right)^{2}, \log \left(\frac{p_{i 1}}{p_{i 3}}\right), \log \left(\frac{p_{i 2}}{p_{i 3}}\right)\right)=\alpha_{1}+\alpha_{2} \cdot\left(\frac{1}{Q_{i}}\right) .
$$

Since the elements of $\mathbf{z}_{i}$ can be nonlinear functions of $\mathbf{x}_{i}$, this specification is more flexible than it might first look, but still it rules out some nonlinearities. For example, the functional form


\begin{equation*}
\mathrm{E}\left(\varepsilon_{i}^{2} \mid \mathbf{x}_{i}\right)=\exp \left(\mathbf{z}_{i}^{\prime} \boldsymbol{\alpha}\right) \tag{2.8.2}
\end{equation*}


might be more attractive because its value is guaranteed to be positive. We consider the linear specification (2.8.1) only because estimating parameters in nonlinear specifications such as (2.8.2) requires the use of nonlinear least squares.

\section*{WLS with Known a}
To isolate the complication arising from the fact that the unknown parameter vector $\boldsymbol{\alpha}$ must be estimated, we first examine the large-sample distribution of the WLS estimator with known $\boldsymbol{\alpha}$. If $\boldsymbol{\alpha}$ is known, the conditional second moment can be calculated from data as $\mathbf{z}_{i}^{\prime} \boldsymbol{\alpha}$, and WLS proceeds exactly as in Section 1.6: dividing both sides of the estimation equation $y_{i}=\mathbf{x}_{i}^{\prime} \boldsymbol{\beta}+\varepsilon_{i}$ by the square root of $\mathbf{z}_{i}^{\prime} \boldsymbol{\alpha}$, to obtain


\begin{equation*}
\tilde{y}_{i}=\tilde{\mathbf{x}}_{i}^{\prime} \boldsymbol{\beta}+\tilde{\varepsilon}_{i}, \tag{2.8.3}
\end{equation*}


where

$$
\tilde{y}_{i} \equiv \frac{y_{i}}{\sqrt{\mathbf{z}_{i}^{\prime} \boldsymbol{\alpha}}}, \tilde{\mathbf{x}}_{i} \equiv \frac{\mathbf{x}_{i}}{\sqrt{\mathbf{z}_{i}^{\prime} \boldsymbol{\alpha}}}, \tilde{\varepsilon}_{i} \equiv \frac{\varepsilon_{i}}{\sqrt{\mathbf{z}_{i}^{\prime} \boldsymbol{\alpha}}}
$$

and then apply OLS. For later reference, write the resulting WLS estimator as $\widehat{\boldsymbol{\beta}}(\mathbf{V})$. It is given by


\begin{align*}
\widehat{\boldsymbol{\beta}}(\mathbf{V}) & \equiv\left(\frac{1}{n} \sum_{i=1}^{n} \tilde{\mathbf{x}}_{i} \tilde{\mathbf{x}}_{i}^{\prime}\right)^{-1} \frac{1}{n} \sum_{i=1}^{n} \tilde{\mathbf{x}}_{i} \cdot \tilde{y}_{i} \\
& =\left(\frac{1}{n} \sum_{i=1}^{n} \frac{1}{\mathbf{z}_{i}^{\prime} \boldsymbol{\alpha}} \mathbf{x}_{i} \mathbf{x}_{i}^{\prime}\right)^{-1} \frac{1}{n} \sum_{i=1}^{n} \frac{1}{\mathbf{z}_{i}^{\prime} \boldsymbol{\alpha}} \mathbf{x}_{i} \cdot y_{i} \\
& =\left(\mathbf{X}^{\prime} \mathbf{V}^{-1} \mathbf{X}\right)^{-1} \mathbf{X}^{\prime} \mathbf{V}^{-1} \mathbf{y} \tag{2.8.4}
\end{align*}


with

$$
\mathbf{V} \equiv\left[\begin{array}{ccc}
\mathbf{z}_{1}^{\prime} \alpha & & \\
& \ddots & \\
& & \mathbf{z}_{n}^{\prime} \alpha
\end{array}\right]
$$

If Assumption 2.3 is strengthened by the condition that


\begin{equation*}
\mathrm{E}\left(\varepsilon_{i} \mid \mathbf{x}_{i}\right)=0, \tag{2.8.5}
\end{equation*}


then $\mathrm{E}\left(\tilde{\varepsilon}_{i} \mid \tilde{\mathbf{x}}_{i}\right)=0$. To see this, note first that, because $\mathbf{z}_{i}$ is a function of $\mathbf{x}_{i}$,


\begin{equation*}
\mathrm{E}\left(\tilde{\varepsilon}_{i} \mid \mathbf{x}_{i}\right)=\mathrm{E}\left(\left.\frac{\varepsilon_{i}}{\sqrt{\mathbf{z}_{i}^{\prime} \boldsymbol{\alpha}}} \right\rvert\, \mathbf{x}_{i}\right)=\frac{1}{\sqrt{\mathbf{z}_{i}^{\prime} \boldsymbol{\alpha}}} \mathrm{E}\left(\varepsilon_{i} \mid \mathbf{x}_{i}\right)=0 . \tag{2.8.6}
\end{equation*}


Second, because $\tilde{\mathbf{x}}_{i}$ is a function of $\mathbf{x}_{i}$, there is no more information in $\tilde{\mathbf{x}}_{i}$ than in $\mathbf{x}_{i}$. So by the Law of Iterated Expectations we have

$$
\mathrm{E}\left(\tilde{\varepsilon}_{i} \mid \tilde{\mathbf{x}}_{i}\right)=\mathrm{E}\left[\mathrm{E}\left(\tilde{\varepsilon}_{i} \mid \mathbf{x}_{i}\right) \mid \tilde{\mathbf{x}}_{i}\right]=0 .
$$

Therefore, provided that $\mathrm{E}\left(\tilde{\mathbf{x}}_{i} \tilde{\mathbf{x}}_{i}^{\prime}\right)$ is nonsingular, Assumptions 2.1-2.5 are satisfied for equation (2.8.3). Furthermore, by construction, the error $\tilde{\varepsilon}_{i}$ is conditionally homoskedastic: $\mathrm{E}\left(\tilde{\varepsilon}_{i}^{2} \mid \tilde{\mathbf{x}}_{i}\right)=1$. So Proposition 2.5 applies with $\sigma^{2}=1$ : the WLS estimator is consistent and asymptotically normal, and the asymptotic variance is


\begin{align*}
\operatorname{Avar}(\widehat{\boldsymbol{\beta}}(\mathbf{V})) & =\mathrm{E}\left(\tilde{\mathbf{x}}_{i} \tilde{\mathbf{x}}_{i}^{\prime}\right)^{-1} \quad(\text { since the error variance is } 1) \\
& =\operatorname{plim}\left(\frac{1}{n} \sum_{i=1}^{n} \tilde{\mathbf{x}}_{i} \tilde{\mathbf{x}}_{i}^{\prime}\right)^{-1} \quad(\text { by ergodic stationarity }) \\
& =\operatorname{plim}\left(\frac{1}{n} \sum_{i=1}^{n} \frac{1}{\mathbf{z}_{i}^{\prime} \boldsymbol{\alpha}} \mathbf{x}_{i} \mathbf{x}_{i}^{\prime}\right)^{-1} \quad\left(\text { since } \tilde{\mathbf{x}}_{i}=\mathbf{x}_{i} / \sqrt{\mathbf{z}_{i} \boldsymbol{\alpha}}\right) \\
& =\operatorname{plim}\left(\frac{1}{n} \mathbf{X}^{\prime} \mathbf{V}^{-1} \mathbf{X}\right)^{-1} \tag{2.8.7}
\end{align*}


So $\left(\frac{1}{n} \mathbf{X}^{\prime} \mathbf{V}^{-1} \mathbf{X}\right)^{-1}$ is a consistent estimator of $\operatorname{Avar}(\widehat{\boldsymbol{\beta}}(\mathbf{V}))$.

\section*{Regression of $\boldsymbol{e}_{i}^{2}$ on $\mathbf{z}_{i}$ Provides a Consistent Estimate of $\boldsymbol{\alpha}$}
If $\boldsymbol{\alpha}$ is unknown, it can be estimated by running a separate regression. (2.8.1) says that $\mathbf{z}_{i}^{\prime} \boldsymbol{\alpha}$ is a regression of $\varepsilon_{i}^{2}$. If we define $\eta_{i} \equiv \varepsilon_{i}^{2}-\mathrm{E}\left(\varepsilon_{i}^{2} \mid \mathbf{x}_{i}\right)$, then (2.8.1) can be\\
written as a regression equation:


\begin{equation*}
\varepsilon_{i}^{2}=\mathbf{z}_{i}^{\prime} \boldsymbol{\alpha}+\eta_{i} . \tag{2.8.8}
\end{equation*}


By construction, $\mathrm{E}\left(\eta_{i} \mid \mathbf{x}_{i}\right)=0$, which, together with the fact that $\mathbf{z}_{i}$ is a function of $\mathbf{x}_{i}$, implies that the regressors $\mathbf{z}_{i}$ are orthogonal to the error term $\eta_{i}$. Hence, provided that $\mathrm{E}\left(\mathbf{z}_{i} \mathbf{z}_{i}^{\prime}\right)$ is nonsingular, Proposition 2.1 is applicable to this auxiliary regression (2.8.8): the OLS estimator of $\boldsymbol{\alpha}$ is consistent and asymptotically normal.

Of course, this we cannot do because we don't observe the error term. However, since the OLS estimator $\mathbf{b}$ for the original regression $y_{i}=\mathbf{x}_{i}^{\prime} \boldsymbol{\beta}+\varepsilon_{i}$ is consistent despite the presence of conditional heteroskedasticity, the OLS residual $e_{i}$ provides a consistent estimate of $\varepsilon_{i}$. It is left to you as an analytical exercise to show that, when $\varepsilon_{i}$ is replaced by $e_{i}$ in the regression (2.8.8), the OLS estimator, call it $\hat{\boldsymbol{\alpha}}$, is consistent for $\boldsymbol{\alpha}$.

\section*{WLS with Estimated $\boldsymbol{\alpha}$}
The WLS estimation of the original equation $y_{i}=\mathbf{x}_{i}^{\prime} \boldsymbol{\beta}+\varepsilon_{i}$ with estimated $\boldsymbol{\alpha}$ is $\widehat{\boldsymbol{\beta}}(\widehat{\mathbf{V}})$, where $\widehat{\mathbf{V}}$ is the $n \times n$ diagonal matrix whose $i$-th diagonal is $\mathbf{z}_{i}^{\prime} \hat{\boldsymbol{\alpha}}$. Under suitable additional conditions (see, e.g., Amemiya (1977) for an explicit statement of such conditions), it can be shown that\\
(a) $\sqrt{n}(\widehat{\boldsymbol{\beta}}(\mathbf{V})-\boldsymbol{\beta})$ and $\sqrt{n}(\widehat{\boldsymbol{\beta}}(\widehat{\mathbf{V}})-\boldsymbol{\beta})$ are asymptotically equivalent in that the difference converges to zero in probability as $n \rightarrow \infty$. Therefore, by Lemma 2.4(a), the asymptotic distribution of $\sqrt{n}(\widehat{\boldsymbol{\beta}}(\widehat{\mathbf{V}})-\boldsymbol{\beta})$ is the same as that of $\sqrt{n}(\widehat{\boldsymbol{\beta}}(\mathbf{V})-\boldsymbol{\beta})$. So $\operatorname{Avar}(\widehat{\boldsymbol{\beta}}(\widehat{\mathbf{V}}))$ equals $\operatorname{Avar}(\widehat{\boldsymbol{\beta}}(\mathbf{V}))$, which in turn is given by (2.8.7);\\
(b) $p \lim \frac{1}{n} \mathbf{X}^{\prime} \widehat{\mathbf{V}}^{-1} \mathbf{X}=p \operatorname{plim} \frac{1}{n} \mathbf{X}^{\prime} \mathbf{V}^{-1} \mathbf{X}$.

So $\left(\frac{1}{n} \mathbf{X}^{\prime} \widehat{\mathbf{V}}^{-1} \mathbf{X}\right)^{-1}$ is consistent for $\operatorname{Avar}(\widehat{\boldsymbol{\beta}}(\widehat{\mathbf{V}}))$.\\
All this may sound complicated, but the operational implication for the WLS estimation of the equation $y_{i}=\mathbf{x}_{i}^{\prime} \boldsymbol{\beta}+\varepsilon_{i}$ is very clear:

Step 1: Estimate the equation $y_{i}=\mathbf{x}_{i}^{\prime} \boldsymbol{\beta}+\varepsilon_{i}$ by OLS and compute the OLS residuals $e_{i}$.\\
Step 2: Regress $e_{i}^{2}$ on $\mathbf{z}_{i}$, to obtain the OLS coefficient estimate $\hat{\boldsymbol{\alpha}}$.\\
Step 3: Re-estimate the equation $y_{i}=\mathbf{x}_{i}^{\prime} \boldsymbol{\beta}+\varepsilon_{i}$ by WLS, using $1 / \sqrt{\mathbf{z}_{i}^{\prime} \hat{\boldsymbol{\alpha}}}$ as the weight for observation $i$.

Since the correct estimate of the asymptotic variance, $\left(\frac{1}{n} \mathbf{X}^{\prime} \widehat{\mathbf{V}}^{-1} \mathbf{X}\right)^{-1}$ in (b), is ( $n$ times) the estimated variance matrix routinely printed out by standard regression\\
packages for the Step 3 regression, calculating appropriate test statistics is quite straightforward: the standard $t$ - and $F$-ratios from Step 3 regression can be used to do statistical inference.

\section*{OLS versus WLS}
Thus we have two consistent and asymptotically normal estimators, the OLS and WLS estimators. We say that a consistently and asymptotically normal estimator is asymptotically more efficient than another consistent and asymptotically normal estimator of the same parameter if the asymptotic variance of the former is no larger than that of the latter. It is left to you as an analytical exercise to show that the WLS estimator is asymptotically more efficient than the OLS estimator.

The superiority of WLS over OLS, however, rests on the premise that the sample size is sufficiently large and the functional form of the conditional second moment is correctly specified. If the functional form is misspecified, the WLS estimator would still be consistent, but its asymptotic variance may or may not be smaller than Avar(b). In finite samples, even if the functional form is correctly specified, the large-sample approximation will probably work less well for the WLS estimator than for OLS because of the estimation of extra parameters ( $\boldsymbol{\alpha}$ ) involved in the WLS procedure.

\section*{QUESTIONS FOR REVIEW}
\begin{enumerate}
  \item Prove: " $\mathrm{E}\left(\eta_{i} \mid \mathbf{x}_{i}\right)=0, \mathbf{z}_{i}$ is a function of $\mathbf{x}_{i}$ " $\Rightarrow$ " $\mathrm{E}\left(\mathbf{z}_{i} \cdot \eta_{i}\right)=\mathbf{0}$." Hint: Law of Total Expectations.
  \item Is the error conditionally homoskedastic in the auxiliary regression (2.8.8)? If so, does it matter for the asymptotic distribution of the WLS estimator?
\end{enumerate}

\subsection*{2.9 Least Squares Projection}
What if the assumptions justifying the large-sample properties of the OLS estimator (except for ergodic stationarity) are not satisfied but we nevertheless go ahead and apply OLS to the sample? What is it that we estimate? The answer is in this section. OLS provides an estimate of the best way linearly to combine the explanatory variables to predict the dependent variable. The linear combination is called the least squares projection.

\section*{Optimally Predicting the Value of the Dependent Variable}
We have been concerned about estimating unknown parameters from a sample. Let us temporarily suspend the role of econometrician and put ourselves in the following situation. There is a random scalar $y$ and a random vector $\mathbf{x}$. We know the joint distribution of $(y, \mathbf{x})$ and the value of $\mathbf{x}$. On the basis of this knowledge we wish to predict $y$. So a predictor is a function $f(\mathbf{x})$ of $\mathbf{x}$ with the functional form $f(\cdot)$ determined by the joint distribution of $(y, \mathbf{x})$. Naturally, we choose the function $f(\cdot)$ so as to minimize some index that is a function of the forecast error $y-f(\mathbf{x})$. We take the loss function to be the mean squared error $\mathrm{E}\left[(y-f(\mathbf{x}))^{2}\right]$ because it seems as reasonable a loss function as any other and, more importantly, because it produces the following convenient result:

Proposition 2.7: $\mathrm{E}(y \mid \mathbf{x})$ is the best.predictor of $y$ in that it minimizes the mean squared error

We have seen in Chapter 1 that the add-and-subtract strategy is effective in showing that the candidate solution minimizes a quadratic function. Let us apply the strategy to the squared error here. Let $f(\mathbf{x})$ be any forecast. Add $\mathrm{E}(y \mid \mathbf{x})$ to the forecast error $y-f(\mathbf{x})$ and then subtract it to obtain the decomposition


\begin{equation*}
y-f(\mathbf{x})=(y-\mathrm{E}(y \mid \mathbf{x}))+(\mathrm{E}(y \mid \mathbf{x})-f(\mathbf{x})) . \tag{2.9.1}
\end{equation*}


So the squared forecast error is


\begin{gather*}
(y-f(\mathbf{x}))^{2}=(y-\mathrm{E}(y \mid \mathbf{x}))^{2}+2(y-\mathrm{E}(y \mid \mathbf{x}))(\mathrm{E}(y \mid \mathbf{x})-f(\mathbf{x})) \\
+(\mathrm{E}(y \mid \mathbf{x})-f(\mathbf{x}))^{2} \tag{2.9.2}
\end{gather*}


Take the expectation of both sides to obtain


\begin{align*}
& \text { mean squared error } \equiv \mathrm{E}\left[(y-f(\mathbf{x}))^{2}\right] \\
& =\mathrm{E}\left[(y-\mathrm{E}(y \mid \mathbf{x}))^{2}\right]+2 \mathrm{E}[(y-\mathrm{E}(y \mid \mathbf{x}))(\mathrm{E}(y \mid \mathbf{x})-f(\mathbf{x}))] \\
& \quad+\mathrm{E}\left[(\mathrm{E}(y \mid \mathbf{x})-f(\mathbf{x}))^{2}\right] \tag{2.9.3}
\end{align*}


It is a straightforward application of the Law of Total Expectations to show that the middle term, which is the covariance between the optimal forecast error and the difference in forecasts, is zero (a review question). Therefore,


\begin{align*}
& \text { mean squared error }=\mathrm{E}\left[(y-\mathrm{E}(y \mid \mathbf{x}))^{2}\right]+\mathrm{E}\left[(\mathrm{E}(y \mid \mathbf{x})-f(\mathbf{x}))^{2}\right] \\
& \geq \mathrm{E}\left[(y-\mathrm{E}(y \mid \mathbf{x}))^{2}\right] \tag{2.9.4}
\end{align*}


which shows that the mean squared error is bounded from below by $\mathrm{E}[(y-\mathrm{E}(y \mid \left.\mathbf{x} \boldsymbol{)})^{2}\right]$, and this lower bound is achieved by the conditional expectation.

\section*{Best Linear Predictor}
It requires the knowledge of the joint distribution of $(y, \mathbf{x})$ to calculate $\mathrm{E}(y \mid \mathbf{x})$, which may be highly nonlinear. We now restrict the predictor to being a linear function of $\mathbf{x}$ and ask: what is the best (in the sense of minimizing the mean squared error) linear predictor of $y$ based on $\mathbf{x}$ ? For this purpose, consider $\boldsymbol{\beta}^{*}$ that satisfies the orthogonality condition


\begin{equation*}
\mathrm{E}\left[\mathbf{x} \cdot\left(y-\mathbf{x}^{\prime} \boldsymbol{\beta}^{*}\right)\right]=\mathbf{0} \text { or } \mathrm{E}\left(\mathbf{x x}^{\prime}\right) \boldsymbol{\beta}^{*}=\mathrm{E}(\mathbf{x} \cdot y) . \tag{2.9.5}
\end{equation*}


The idea is to choose $\boldsymbol{\beta}^{*}$ so that the forecast error $y-\mathbf{x}^{\prime} \boldsymbol{\beta}^{*}$ is orthogonal to $\mathbf{x}$. If $\mathrm{E}\left(\mathbf{x x}^{\prime}\right)$ is nonsingular, the orthogonality condition can be solved for $\boldsymbol{\beta}^{*}$ :


\begin{equation*}
\boldsymbol{\beta}^{*}=\left[\mathrm{E}\left(\mathbf{x} \mathbf{x}^{\prime}\right)\right]^{-1} \mathrm{E}(\mathbf{x} \cdot y) . \tag{2.9.6}
\end{equation*}


The least squares (or linear) projection of $y$ on $\mathbf{x}$, denoted $\widehat{\mathrm{E}}^{*}(y \mid \mathbf{x})$, is defined as $\mathbf{x}^{\prime} \boldsymbol{\beta}^{*}$, where $\boldsymbol{\beta}^{*}$ satisfies (2.9.5) and is called the least squares projection coefficients.

Proposition 2.8: The least squares projection $\widehat{\mathrm{E}}^{*}(y \mid \mathbf{x})$ is the best linear predictor of $y$ in that it minimizes the mean squared error.

The add-and-subtract strategy also works here.

Proof. For any linear predictor $\mathbf{x}^{\prime} \boldsymbol{\mathcal { B }}$,

$$
\begin{aligned}
& \text { mean squared error } \equiv \mathrm{E}\left[\left(y-\mathbf{x}^{\prime} \tilde{\boldsymbol{\beta}}\right)^{2}\right] \\
& \quad=\mathrm{E}\left\{\left[\left(y-\mathbf{x}^{\prime} \boldsymbol{\beta}^{*}\right)+\mathbf{x}^{\prime}\left(\boldsymbol{\beta}^{*}-\widetilde{\boldsymbol{\beta}}\right)\right]^{2}\right\} \quad \text { (by the add-and-subtract strategy) } \\
& \quad=\mathrm{E}\left[\left(y-\mathbf{x}^{\prime} \boldsymbol{\beta}^{*}\right)^{2}\right]+2\left(\boldsymbol{\beta}^{*}-\widetilde{\boldsymbol{\beta}}\right)^{\prime} \mathrm{E}\left[\mathbf{x} \cdot\left(y-\mathbf{x}^{\prime} \boldsymbol{\beta}^{*}\right)\right]+\mathrm{E}\left[\left(\mathbf{x}^{\prime}\left(\boldsymbol{\beta}^{*}-\widetilde{\boldsymbol{\beta}}\right)\right)^{2}\right] \\
& \quad=\mathrm{E}\left[\left(y-\mathbf{x}^{\prime} \boldsymbol{\beta}^{*}\right)^{2}\right]+\mathrm{E}\left[\left(\mathbf{x}^{\prime}\left(\widetilde{\boldsymbol{\beta}}-\boldsymbol{\beta}^{*}\right)\right)^{2}\right]
\end{aligned}
$$

(by the orthogonality condition (2.9.5))

$$
\geq \mathrm{E}\left[\left(y-\mathbf{x}^{\prime} \boldsymbol{\beta}^{*}\right)^{2}\right] .
$$ \(\square\)

In contrast to the best predictor, which is the conditional expectation, the best linear predictor requires only the knowledge of the second moments of the joint distribution of ( $y, \mathbf{x}$ ) to calculate (see (2.9.6)).

If one of the regressors $\mathbf{x}$ is a constant, the least squares projection coefficients can be written in terms of variances and covariances. Let $\tilde{\mathbf{x}}$ be the vector of nonconstant regressors so that

$$
\mathbf{x}=\left[\begin{array}{c}
1 \\
\tilde{\mathbf{x}}
\end{array}\right]
$$

An analytical exercise to this chapter asks you to prove that


\begin{equation*}
\widehat{\mathrm{E}}^{*}(y \mid \mathbf{x}) \equiv \widehat{\mathrm{E}}^{*}(y \mid 1, \tilde{\mathbf{x}})=\mu+\gamma^{\prime} \tilde{\mathbf{x}} \tag{2.9.7}
\end{equation*}


where

$$
\boldsymbol{\gamma}=\operatorname{Var}(\tilde{\mathbf{x}})^{-1} \operatorname{Cov}(\tilde{\mathbf{x}}, y), \mu=\mathrm{E}(y)-\boldsymbol{\gamma}^{\prime} \mathrm{E}(\tilde{\mathbf{x}}) .
$$

This formula is the population analogue of the "deviations-from-the-mean regression" formula of Chapter 1 (Analytical Exercise 3).

\section*{OLS Consistently Estimates the Projection Coefficients}
Now let us put the econometrician's hat back on and consider estimating $\boldsymbol{\beta}^{*}$. Suppose we have a sample of size $n$ drawn from an ergodic stationary stochastic process $\left\{y_{i}, \mathbf{x}_{i}\right\}$ with the joint distribution $\left(y_{i}, \mathbf{x}_{i}\right)$ (which does not depend on $i$ because of stationarity) identical to that of ( $y, \mathbf{x}$ ) above. So, for example, $\mathrm{E}\left(\mathbf{x}_{i} \mathbf{x}_{i}^{\prime}\right)=\mathrm{E}\left(\mathbf{x x}^{\prime}\right)$. By the Ergodic Theorem the second moments in (2.9.6) can be consistently estimated by the corresponding sample second moments. Thus a consistent estimator of the projection coefficients $\boldsymbol{\beta}^{*}$ is

$$
\left(\frac{1}{n} \sum_{i=1}^{n} \mathbf{x}_{i} \mathbf{x}_{i}^{\prime}\right)^{-1}\left(\frac{1}{n} \sum_{i=1}^{n} \mathbf{x}_{i} \cdot y_{i}\right)=\left(\mathbf{X}^{\prime} \mathbf{X}\right)^{-1} \mathbf{X}^{\prime} \mathbf{y}
$$

which is none other than the OLS estimator $\mathbf{b}$. That is, under Assumption 2.2 (ergodic stationarity) and Assumption 2.4 guaranteeing the nonsingularity of $\mathrm{E}\left(\mathbf{x x}^{\prime}\right)$, the OLS estimator is always consistent for the projection coefficient vector, the $\boldsymbol{\beta}^{*}$ that satisfies the orthogonality condition (2.9.5).

\section*{QUESTIONS FOR REVIEW}
\begin{enumerate}
  \item (Unforecastability of forecast error) For the minimum mean square error forecast $\mathrm{E}(y \mid \mathbf{x})$, the forecast error is orthogonal to any function $\phi(\mathbf{x})$ of $\mathbf{x}$. That is, $\mathrm{E}[\eta \phi(\mathbf{x})]=0$ where $\eta \equiv y-\mathrm{E}(y \mid \mathbf{x})$. Prove this. Hint: The Law of Total Expectations. Show that the middle term on the RHS of (2.9.3) is zero by setting $\phi(\mathbf{x})=\mathrm{E}(y \mid \mathbf{x})-f(\mathbf{x})$.
  \item (Forecasting white noise) Suppose $\left\{\varepsilon_{i}\right\}$ is white noise. What is $\widehat{\mathrm{E}}^{*}\left(\varepsilon_{i} \mid \varepsilon_{i-1}\right.$, $\left.\varepsilon_{i-2}, \ldots, \varepsilon_{i-m}\right)$ ? What is $\widehat{\mathrm{E}}^{*}\left(\varepsilon_{i} \mid 1, \varepsilon_{i-1}, \varepsilon_{i-2}, \ldots, \varepsilon_{i-m}\right)$ ? Is it true that $\mathrm{E}\left(\varepsilon_{i} \mid\right. \left.\varepsilon_{i-1}, \varepsilon_{i-2}, \ldots, \varepsilon_{i-m}\right)=0$ ? Hint: Is it zero for the process of Example 2.4?
  \item (Conditional expectations that are linear) Suppose $\mathrm{E}(y \mid \tilde{\mathbf{x}})=\mu+\boldsymbol{\gamma}^{\prime} \tilde{\mathbf{x}}$. Show: $\widehat{\mathrm{E}}^{*}(y \mid 1, \tilde{\mathbf{x}})=\mathrm{E}(y \mid \tilde{\mathbf{x}})$.
  \item (Partitioned projection) Consider the model $y_{i}=\mathbf{x}_{i}^{\prime} \boldsymbol{\beta}+\mathbf{z}_{i}^{\prime} \boldsymbol{\delta}+\varepsilon_{i}$ with $\mathrm{E}\left(\mathbf{x}_{i} \cdot\right. \left.\varepsilon_{i}\right)=\mathbf{0}, \mathrm{E}\left(\mathbf{z}_{i} \cdot \varepsilon_{i}\right) \neq \mathbf{0}$, and $\mathrm{E}\left(\mathbf{z}_{i} \mathbf{x}_{i}^{\prime}\right)=\mathbf{0}$. Thus, $\mathbf{z}_{i}$ is not predetermined (i.e., not orthogonal to the error term), but it is unrelated to the predetermined regressor $\mathbf{x}_{i}$ in that the cross moments are zero.\\
(a) Show that the least squares projection coefficient of $\mathbf{x}_{i}$ in the projection of $y_{i}$ on $\mathbf{x}_{i}$ and $\mathbf{z}_{i}$ is $\boldsymbol{\beta}$. Hint: Calculate $\widehat{\mathrm{E}}^{*}\left(\varepsilon_{i} \mid \mathbf{x}_{i}, \mathbf{z}_{i}\right)$.\\
(b) What is the least squares projection coefficient of $\mathbf{x}_{i}$ in the least squares projection of $y_{i}$ on $\mathbf{x}_{i}$ ? Hint: Treat $\mathbf{z}_{i}^{\prime} \boldsymbol{\delta}+\varepsilon_{i}$ as the error term.\\
(c) Which projection would you use for estimating $\boldsymbol{\beta}$ ? Hint: You want the error variance to be smaller.
\end{enumerate}

\subsection*{2.10 Testing for Serial Correlation}
As remarked in Section 2.3 (see (2.3.3)), when the regressors include a constant (true in virtually all known applications), Assumption 2.5 implies that the error term is a scalar martingale difference sequence (m.d.s.), so if the error is found to be serially correlated, that is an indication of a failure of Assumption 2.5. Serial correlation has traditionally been an important subject in econometrics, and there are available a number of tests for serial correlation (i.e., tests of the null of no serial correlation in the error term). Some of them, however, require that the regressors be strictly exogenous. The test to be presented in this section does not require strict exogeneity. Because the issue of serial correlation arises only in timeseries models, we use the subscript " $t$ " instead of " $i$ " in this section (and the next). Throughout this section we assume that the regressors include a constant.

It would be nice to have those tests extended to cover serial correlation in $\mathbf{g}_{t}$ ( $\equiv \mathbf{x}_{t} \cdot \varepsilon_{t}$ ), but no such tests have been proposed to gain acceptance. This is a gap in the literature, but not a serious one, because nowadays researchers know how to live with serial correlation in $\mathbf{g}_{\text {. }}$ That is, as will be shown in Chapter 6, there is available a method to do inference in the presence of serial correlation in $\mathbf{g}_{t}$.

Earlier in this chapter we learned how to calculate standard errors that are robust to conditional heteroskedasticity. In Chapter 6, those standard errors will be made robust to serial correlation as well.

\section*{Box-Pierce and Ljung-Box}
Before turning to the tests for serial correlation in the error term, we temporarily step outside the regression framework and consider serial correlation in a univariate time series. Suppose we have a sample of size $n,\left\{z_{1}, \ldots, z_{n}\right\}$, drawn from a scalar covariance-stationary process. In Section 2.2 we defined the (population) $j$-th order autocovariance $\gamma_{j}$. The sample $j$-th order autocovariance is


\begin{equation*}
\hat{\gamma}_{j} \equiv \frac{1}{n} \sum_{t=j+1}^{n}\left(z_{t}-\bar{z}_{n}\right)\left(z_{t-j}-\bar{z}_{n}\right) \quad(j=0,1, \ldots), \tag{2.10.1}
\end{equation*}


where

$$
\bar{z}_{n} \equiv \frac{1}{n} \sum_{t=1}^{n} z_{t} .
$$

(If the population mean $\mathrm{E}\left(z_{t}\right)$ is known, it can replace the sample mean $\bar{z}_{n}$; doing so would improve the small sample property.) Here, even though only $n-j$ terms are in the sum, the denominator is $n$ rather than $n-j$. Whether the sum of $n-j$ terms is divided by $n$ or by $n-j$ does not affect large-sample results. For moderate sample sizes, however, the numerical difference can be substantial, and you should always be explicit about which is used as the denominator. The sample $\boldsymbol{j}$-th order autocorrelation coefficient, $\hat{\rho}_{j}$, is defined as


\begin{equation*}
\hat{\rho}_{j} \equiv \frac{\hat{\gamma}_{j}}{\hat{\gamma}_{0}} \quad(j=1,2, \ldots) . \tag{2.10.2}
\end{equation*}


If $\left\{z_{t}\right\}$ is ergodic stationary, then it is easy to show (see a review question) that $\hat{\gamma}_{j}$ is consistent for $\gamma_{j}(j=0,1,2, \ldots)$. Hence, by Lemma 2.3(a), $\hat{\rho}_{j}$ is consistent for $\rho_{j}(j=1,2, \ldots)$. In particular, if $\left\{z_{t}\right\}$ is serially uncorrelated, then all the sample autocorrelation coefficients converge to 0 in probability. To test for serial correlation, however, we need to know the asymptotic distribution of $\sqrt{n} \hat{\rho}_{j}$. It is provided by

Proposition 2.9 (special case of Theorem 6.7 of Hall and Heyde (1980)): Suppose $\left\{z_{t}\right\}$ can be written as $\mu+\varepsilon_{t}$, where $\varepsilon_{t}$ is a stationary martingale difference sequence with "own" conditional homoskedasticity:\\
(own conditional homoskedasticity) $\mathrm{E}\left(\varepsilon_{t}^{2} \mid \varepsilon_{t-1}, \varepsilon_{t-2}, \ldots\right)=\sigma^{2}, \sigma^{2}>0$.

Let the sample autocorrelation $\hat{\rho}_{j}$ be defined as in (2.10.1) and (2.10.2). Then

$$
\sqrt{n} \hat{\boldsymbol{\gamma}} \underset{\mathrm{~d}}{\rightarrow} N\left(\mathbf{0}, \sigma^{4} \mathbf{I}_{p}\right) \text { and } \sqrt{n} \hat{\boldsymbol{\rho}} \underset{\mathrm{~d}}{\rightarrow} N\left(\mathbf{0}, \mathbf{I}_{p}\right),
$$

where $\hat{\boldsymbol{\gamma}}=\left(\hat{\gamma}_{1}, \hat{\gamma}_{2}, \ldots, \hat{\gamma}_{p}\right)^{\prime}$ and $\hat{\boldsymbol{\rho}}=\left(\hat{\rho}_{1}, \hat{\rho}_{2}, \ldots, \hat{\rho}_{p}\right)^{\prime}$.

Here, the process $\varepsilon_{t}$ is not required to be ergodic. Proving this under the additional condition of ergodicity is left as an analytical exercise. Thus, asymptotically, ( $\sqrt{n}$ times) the autocorrelations are i.i.d. and the distribution is $N(0,1)$. The process $\left\{\varepsilon_{t}\right\}$ assumed here is more general than independent white noise processes, but the conditional second moment has to be constant, so this result does not cover ARCH processes, for example.

One way to test for serial correlation in the series is to check whether the firstorder autocorrelation, $\rho_{1}$, is 0 . Proposition 2.9 implies that


\begin{equation*}
\sqrt{n} \hat{\rho}_{1}=\frac{\hat{\rho}_{1}}{1 / \sqrt{n}} \underset{\mathrm{~d}}{\rightarrow} N(0,1) \tag{2.10.3}
\end{equation*}


So the $t$-statistic formed as the ratio of $\hat{\rho}_{1}$ to a "standard error" of $1 / \sqrt{n}$ is asymptotically standard normal.

We can also test whether a group of autocorrelations are simultaneously zero. Let $\hat{\boldsymbol{\rho}}=\left(\hat{\rho}_{1}, \ldots, \hat{\rho}_{p}\right)^{\prime}$ be the $p$-dimensional vector collecting the first $p$ sample autocorrelations. Since the elements of $\sqrt{n} \hat{\rho}$ are asymptotically independent and individually distributed as standard normal, their squared sum, called the BoxPierce $Q$ because it was first considered by Box and Pierce (1970), is asymptotically chi-squared:


\begin{equation*}
\text { Box-Pierce } Q \text { statistic } \equiv n \sum_{j=1}^{p} \hat{\rho}_{j}^{2}=\sum_{j=1}^{p}\left(\sqrt{n} \hat{\rho}_{j}\right)^{2} \rightarrow \chi^{2}(p) \tag{2.10.4}
\end{equation*}


It is easy to show (see a review question) that the following modification, called the Ljung-Box $\boldsymbol{Q}$, is asymptotically equivalent in that its difference from the BoxPierce $Q$ vanishes in large samples. So by Lemma 2.4(a) it, too, is asymptotically chi-squared:


\begin{equation*}
\text { Ljung-Box } Q \text { statistic } \equiv n \cdot(n+2) \sum_{j=1}^{p} \frac{\hat{\rho}_{j}^{2}}{n-j}=\sum_{j=1}^{p} \frac{n+2}{n-j}\left(\sqrt{n} \hat{\rho}_{j}\right)^{2} \underset{\mathrm{~d}}{\rightarrow} \chi^{2}(p) . \tag{2.10.5}
\end{equation*}


This modification often provides a better approximation to the chi-square distribution for moderate sample sizes (you will be asked to verify this in a Monte Carlo exercise). For either statistic, there is no clear guide to the choice of $p$. If $p$ is too small, there is a danger of missing the existence of higher-order autocorrelations, but if $p$ is too large relative to the sample size, its finite-sample distribution is likely to deteriorate, diverging greatly from the chi-square distribution.

\section*{Sample Autocorrelations Calculated from Residuals}
Now go back to the regression model described by Assumptions 2.1-2.5. If the error term $\varepsilon_{t}$ were observable, we would calculate the sample autocorrelations as


\begin{equation*}
\tilde{\rho}_{j} \equiv \frac{\tilde{\gamma}_{j}}{\tilde{\gamma}_{0}}(j=1,2, \ldots), \tag{2.10.6}
\end{equation*}


where


\begin{equation*}
\tilde{\gamma}_{j} \equiv \frac{1}{n} \sum_{t=j+1}^{n} \varepsilon_{t} \varepsilon_{t-j} \quad(j=0,1,2, \ldots) \tag{2.10.7}
\end{equation*}


(There is no need to subtract the sample mean because the population mean is zero, an implication of the inclusion of a constant in the regressors.) Since $\left\{\varepsilon_{t} \varepsilon_{t-j}\right\}$ is ergodic stationary by Assumption 2.2, $\tilde{\gamma}_{j}$ converges in probability to the corresponding population mean, $\mathrm{E}\left(\varepsilon_{t} \varepsilon_{t-j}\right)$, for all $j$, and $\tilde{\rho}_{j}$ is consistent for the population $j$-th autocorrelation coefficient of $\varepsilon_{t}$.

Next consider the more realistic case where we do not observe the error term. We can replace $\varepsilon_{t}$ in the above formula by the OLS estimate $e_{t}$ and calculate the sample autocorrelations as


\begin{equation*}
\hat{\rho}_{j} \equiv \frac{\hat{\gamma}_{j}}{\hat{\gamma}_{0}} \quad(j=1,2, \ldots) \tag{2.10.8}
\end{equation*}


where


\begin{equation*}
\hat{\gamma}_{j} \equiv \frac{1}{n} \sum_{t=j+1}^{n} e_{t} e_{t-j} \quad(j=0,1,2, \ldots) \tag{2.10.9}
\end{equation*}


(Because the regressors include a constant, the normal equation corresponding to the constant ensures that the sample mean of $e_{t}$ is zero. So there is no need to subtract the sample mean.) Is it all right to use $\hat{\rho}_{j}$ (calculated from the residuals) instead of $\tilde{\rho}_{j}$ and the residual-based $Q$ statistics derived from $\left\{\hat{\rho}_{j}\right\}$ for testing for serial correlation? The answer is yes, but only if the regressors are strictly exogenous.

Recall that the expression linking the residual $e_{t}$ to the true error term $\varepsilon_{t}$ was given in (2.3.9). Using this, the difference between $\tilde{\gamma}_{j}$ and $\hat{\gamma}_{j}$ can be written as


\begin{align*}
\hat{\gamma}_{j} \equiv & \frac{1}{n} \sum_{t=j+1}^{n} e_{t} e_{t-j} \\
= & \frac{1}{n} \sum_{t=j+1}^{n}\left[\varepsilon_{t}-\mathbf{x}_{t}^{\prime}(\mathbf{b}-\boldsymbol{\beta})\right]\left[\varepsilon_{t-j}-\mathbf{x}_{t-j}^{\prime}(\mathbf{b}-\boldsymbol{\beta})\right] \\
= & \tilde{\gamma}_{j}-\frac{1}{n} \sum_{t=j+1}^{n}\left(\mathbf{x}_{t-j} \cdot \varepsilon_{t}+\mathbf{x}_{t} \cdot \varepsilon_{t-j}\right)^{\prime}(\mathbf{b}-\boldsymbol{\beta}) \\
& +(\mathbf{b}-\boldsymbol{\beta})^{\prime}\left(\frac{1}{n} \sum_{t=j+1}^{n} \mathbf{x}_{t} \mathbf{x}_{t-j}^{\prime}\right)(\mathbf{b}-\boldsymbol{\beta}) \tag{2.10.10}
\end{align*}


If $\mathrm{E}\left(\mathbf{x}_{t} \cdot \varepsilon_{t-j}\right), \mathrm{E}\left(\mathbf{x}_{t-j} \cdot \varepsilon_{t}\right)$, and $\mathrm{E}\left(\mathbf{x}_{t} \mathbf{x}_{t-j}^{\prime}\right)$ are all finite, then the second and the third terms vanish (converges to zero in probability) because $\mathbf{b}-\boldsymbol{\beta} \rightarrow_{\mathrm{p}} \mathbf{0}$. Therefore,

$$
\hat{\gamma}_{j}-\tilde{\gamma}_{j} \underset{\mathrm{p}}{\rightarrow} 0 \quad(j=0,1,2, \ldots),
$$

and thus the difference between $\tilde{\rho}_{j}$ and $\hat{\rho}_{j}$, too, vanishes in large samples.\\
However, $\sqrt{n}$ times the difference does not. $\sqrt{n} \tilde{\rho}_{j}$ and $\sqrt{n} \hat{\rho}_{j}$ can be written as


\begin{equation*}
\sqrt{n} \tilde{\rho}_{j}=\frac{\sqrt{n} \tilde{\gamma}_{j}}{\tilde{\gamma}_{0}} \text { and } \sqrt{n} \hat{\rho}_{j}=\frac{\sqrt{n} \hat{\gamma}_{j}}{\hat{\gamma}_{0}} . \tag{2.10.11}
\end{equation*}


We have just seen that, for both $\sqrt{n} \tilde{\rho}_{j}$ and $\sqrt{n} \hat{\rho}_{j}$, the denominator $\rightarrow_{\mathrm{p}} \sigma^{2}$. So the difference between $\sqrt{n} \tilde{\rho}_{j}$ and $\sqrt{n} \hat{\rho}_{j}$ will vanish if the difference between $\sqrt{n} \tilde{\gamma}_{j}$ and $\sqrt{n} \hat{\gamma}_{j}$ does (if you are not convinced, see Review Question 3 below). Now, by multiplying both sides of (2.10.10) by $\sqrt{n}$, we obtain


\begin{align*}
\sqrt{n} \hat{\gamma}_{j}= & \sqrt{n} \tilde{\gamma}_{j}-\frac{1}{n} \sum_{t=j+1}^{n}\left(\mathbf{x}_{t-j} \cdot \varepsilon_{t}+\mathbf{x}_{t} \cdot \varepsilon_{t-j}\right)^{\prime} \sqrt{n}(\mathbf{b}-\boldsymbol{\beta}) \\
& +\sqrt{n}(\mathbf{b}-\boldsymbol{\beta})^{\prime}\left(\frac{1}{n} \sum_{t=j+1}^{n} \mathbf{x}_{t} \mathbf{x}_{t-j}^{\prime}\right)(\mathbf{b}-\boldsymbol{\beta}) \tag{2.10.12}
\end{align*}


Because $\sqrt{n}(\mathbf{b}-\boldsymbol{\beta})$ converges to a random variable (whose distribution is normal), the third term on the RHS vanishes by Lemma 2.4(b). Regarding the second term, we have


\begin{equation*}
\frac{1}{n} \sum_{t=j+1}^{n}\left(\mathbf{x}_{t-j} \cdot \varepsilon_{t}+\mathbf{x}_{t} \cdot \varepsilon_{t-j}\right) \underset{\mathrm{p}}{\rightarrow} \mathrm{E}\left(\mathbf{x}_{t-j} \cdot \varepsilon_{t}\right)+\mathrm{E}\left(\mathbf{x}_{t} \cdot \varepsilon_{t-j}\right) . \tag{2.10.13}
\end{equation*}


If the regressors are strictly exogenous in the sense that $\mathrm{E}\left(\mathbf{x}_{t} \cdot \varepsilon_{s}\right)=\mathbf{0}$ for all $t, s$, then


\begin{equation*}
\mathrm{E}\left(\mathbf{x}_{t-j} \cdot \varepsilon_{t}\right)+\mathrm{E}\left(\mathbf{x}_{t} \cdot \varepsilon_{t-j}\right)=\mathbf{0} \tag{2.10.14}
\end{equation*}


and by Lemma 2.4(b) the second term converges to zero in probability. So the difference between $\sqrt{n} \tilde{\rho}_{j}$ and $\sqrt{n} \hat{\rho}_{j}$ vanishes, which means that the $Q$ statistic calculated from the regression residuals $\left\{e_{t}\right\}$, too, is asymptotically chi-squared, and we can use this residual-based $Q$ to test for serial correlation. If, on the other hand, the regressors are not strictly exogenous, then there is no guarantee that (2.10.14) holds. Consequently, the residual-based $Q$ statistic may not be asymptotically chisquared.

\section*{Testing with Predetermined, but Not Strictly Exogenous, Regressors}
Therefore, when the regressors are not strictly exogenous, we need to modify the $Q$ statistic to restore its asymptotic distribution. For this purpose, consider two restrictions:\\
(stronger form of predeterminedness) $\mathrm{E}\left(\varepsilon_{t} \mid \varepsilon_{t-1}, \varepsilon_{t-2}, \ldots, \mathbf{x}_{t}, \mathbf{x}_{t-1}, \ldots\right)=0$,\\
(stronger form of conditional homoskedasticity)


\begin{equation*}
\mathrm{E}\left(\varepsilon_{t}^{2} \mid \varepsilon_{t-1}, \varepsilon_{t-2}, \ldots, \mathbf{x}_{t}, \mathbf{x}_{t-1}, \ldots\right)=\sigma^{2}>0 . \tag{2.10.16}
\end{equation*}


The first condition is just reproducing (2.3.1) (with " $t$ " now being used as the subscript). As was shown in Section 2.3, this is stronger than Assumption 2.3 and implies that $\mathbf{g}_{t}\left(\equiv \mathbf{x}_{t} \cdot \varepsilon_{t}\right)$ is an m.d.s. Condition (2.10.16), with the conditioning set including $\mathbf{x}_{t}$, is obviously stronger than Assumption 2.7 (conditional homoskedasticity). It also is stronger than the own conditional homoskedasticity assumption in Proposition 2.9 because the conditioning set includes current and past $\mathbf{x}$ as well as past $\varepsilon$. The next result shows that under these additional conditions there is an appropriate modification of the $Q$ statistic.

Proposition 2.10 (testing for serial correlation with predetermined regressors): Suppose that Assumptions 2.1, 2.2, 2.4, (2.10.15), and (2.10.16) are satisfied. Let the sample autocorrelation of the OLS residuals, $\hat{\rho}_{j}$, be defined as in (2.10.8). Then,


\begin{equation*}
\sqrt{n} \hat{\boldsymbol{\gamma}} \underset{\mathrm{~d}}{\rightarrow} N\left(\mathbf{0}, \sigma^{4} \cdot\left(\mathbf{I}_{p}-\boldsymbol{\Phi}\right)\right) \text { and } \sqrt{n} \hat{\boldsymbol{\rho}} \underset{\mathrm{~d}}{\rightarrow} N\left(\mathbf{0}, \mathbf{I}_{p}-\boldsymbol{\Phi}\right), \tag{2.10.17}
\end{equation*}


where $\hat{\boldsymbol{\gamma}}=\left(\hat{\gamma}_{1}, \hat{\gamma}_{2}, \ldots, \hat{\gamma}_{p}\right)^{\prime}, \hat{\boldsymbol{\rho}}=\left(\hat{\rho}_{1}, \hat{\rho}_{2}, \ldots, \hat{\rho}_{p}\right)^{\prime}$, and $\phi_{j k}(\equiv(j, k)$ element of the $p \times p$ matrix $\boldsymbol{\Phi}$ ) is given by


\begin{equation*}
\phi_{j k}=\mathrm{E}\left(\mathbf{x}_{t} \cdot \varepsilon_{t-j}\right)^{\prime} \mathrm{E}\left(\mathbf{x}_{t} \mathbf{x}_{t}^{\prime}\right)^{-1} \mathrm{E}\left(\mathbf{x}_{t} \cdot \varepsilon_{t-k}\right) / \sigma^{2} . \tag{2.10.18}
\end{equation*}


The proof, though not very difficult, is relegated to Appendix 2.B. By the Ergodic Theorem, matrix $\boldsymbol{\Phi}$ is consistently estimated by its sample counterpart:


\begin{equation*}
\widehat{\boldsymbol{\Phi}} \equiv\left(\hat{\phi}_{j k}\right), \hat{\phi}_{j k} \equiv \overline{\boldsymbol{\mu}}_{j}^{\prime} \mathbf{S}_{\mathbf{x x}}^{-1} \overline{\boldsymbol{\mu}}_{k} / s^{2} \quad(j, k=1,2, \ldots, p) \tag{2.10.19}
\end{equation*}


where

$$
s^{2} \equiv \frac{1}{n-K} \sum_{t=1}^{n} e_{t}^{2}, \bar{\mu}_{j} \equiv \frac{1}{n} \sum_{t=j+1}^{n} \mathbf{x}_{t} \cdot e_{t-j}
$$

It follows from this and Proposition 2.10 that


\begin{equation*}
\text { modified Box-Pierce } Q \equiv n \cdot \hat{\boldsymbol{\rho}}^{\prime}\left(\mathbf{I}_{p}-\widehat{\boldsymbol{\Phi}}\right)^{-1} \hat{\boldsymbol{\rho}} \underset{\mathrm{~d}}{\rightarrow} \chi^{2}(p) \text {. } \tag{2.10.20}
\end{equation*}


As long as the regressors are predetermined and the error is conditionally homoskedastic in the sense of (2.10.15) and (2.10.16), this modified $Q$ statistic can be used even when the regressors are not strictly exogenous.

\section*{An Auxiliary Regression-Based Test}
Although calculating this modified $Q$ statistic is straightforward with matrix-based software, it is useful to find an asymptotically equivalent statistic that can be calculated from regression packages. For this purpose, consider the following auxiliary regression:


\begin{equation*}
\text { regress } e_{t} \text { on } \mathbf{x}_{t}, e_{t-1}, e_{t-2}, \ldots, e_{t-p} \tag{2.10.21}
\end{equation*}


To run this auxiliary regression for $t=1,2, \ldots, n$, we need data on ( $e_{0}, e_{-1}, \ldots$, $e_{-p+1}$ ). It does not matter asymptotically which particular numbers to assign to them, but it seems sensible to set them equal to 0 , their expected value. ${ }^{15}$ From this auxiliary regression, we can calculate the $F$ statistic for the hypothesis that the $p$ coefficients of $e_{t-1}, e_{t-2}, \ldots, e_{t-p}$ are all zero. Given Proposition 2.5(c), it is only natural to wonder whether $p \cdot F$ is asymptotically $\chi^{2}(p)$. This conjecture is indeed

\footnotetext{${ }^{15}$ Another asymptotically equivalent choice is to run the auxiliary regression for $t=p+1, p+2, \ldots, n$.
}
true: under the hypothesis of Proposition 2.10, the modified $Q$ statistic (2.10.20) is asymptotically equivalent to $p \cdot F$ (i.e., the difference between the two converges to zero in probability as $n \rightarrow \infty$ ), so $p \cdot F$, too, is asymptotically chi-squared (showing this is an analytical exercise).

This $p \cdot F$ statistic, in turn, is asymptotically equivalent to $n R^{2}$ from the auxiliary regression. This can be shown as follows. Recall the algebraic result from Chapter 1 that the $F$-ratio can be calculated from the difference in the sum of squared residuals between the unrestricted and restricted regressions. The unrestricted regression in the present context is (2.10.21) while the restricted regression is


\begin{equation*}
\text { regress } e_{t} \text { on } \mathbf{x}_{t} \text {. } \tag{2.10.22}
\end{equation*}


Therefore, if $S S R_{U}$ and $S S R_{R}$ are $S S R$ s from (2.10.21) and (2.10.22), respectively, we have


\begin{equation*}
p \cdot F=\frac{S S R_{R}-S S R_{U}}{S S R_{U} /\left(n-\# \mathbf{x}_{t}-p\right)}=\left(n-\# \mathbf{x}_{t}-p\right) \frac{S S R_{R}-S S R_{U}}{S S R_{U}} \tag{2.10.23}
\end{equation*}


where $\# \mathbf{x}_{t}$ is the number of variables in $\mathbf{x}_{t}$. However, since $e_{t}$ is the residual from the original regression (a regression of $y_{t}$ on $\mathbf{x}_{t}$ ), the regressors $\mathbf{x}_{t}$ in (2.10.22) have no explanatory power. So $S S R_{R}=\mathbf{e}^{\prime} \mathbf{e}$ where $\mathbf{e}$ is the $n$-dimensional vector of residuals from the original regression, and (2.10.23) is numerically identical to

$$
\left(n-\# \mathbf{x}_{t}-p\right) \frac{R_{u c}^{2}}{1-R_{u c}^{2}}
$$

where $R_{u c}^{2}$ is the uncentered $R^{2}$ for the auxiliary regression (2.10.21). ${ }^{16}$ But since the sample mean of $e_{t}$ is by construction zero (this is because $\mathbf{x}_{t}$ includes a constant), $R_{u c}^{2}$ is numerically identical to the $R^{2}$ for the auxiliary regression. Thus, we have proved the algebraic result that

$$
p \cdot F=\left(n-\# \mathbf{x}_{1}-p\right) \frac{R^{2}}{1-R^{2}},
$$

where $R^{2}$ is equal to the $R^{2}$ for the auxiliary regression (2.10.21). Solving this equation for $R^{2}$ and multiplying both sides by $n$, we obtain

\footnotetext{${ }^{16}$ Recall from (1.2.16) that the uncentered $R^{2}$ for a regression of $\mathbf{y}$ on $\mathbf{x}$ is defined as

$$
R_{u c}^{2} \equiv \frac{\mathbf{y}^{\prime} \mathbf{y}-\mathbf{e}^{\prime} \mathbf{e}}{\mathbf{y}^{\prime} \mathbf{y}}
$$

In the present context, the $\mathbf{y}^{\prime} \mathbf{y}$ in this formula is $\mathbf{e}^{\prime} \mathbf{e}$ while the $\mathbf{e}^{\prime} \mathbf{e}$ in the formula is $S S R_{U}$.
}
$$
n R^{2}=\frac{n}{n-\# \mathbf{x}_{t}-p} \cdot \frac{1}{1+\frac{p \cdot F}{n-\# \mathbf{x}_{t}-p}} \cdot p \cdot F .
$$

Since $\frac{p \cdot F}{n-\mathrm{H}_{t}-p} \rightarrow_{\mathrm{p}} 0$ as $n \rightarrow \infty$ (an implication of Lemma 2.4(b)), this equation shows that $p \cdot F-n R^{2} \rightarrow_{\mathrm{p}} 0$. Therefore, $n R^{2}$ from the auxiliary regression (2.10.21) is asymptotically $\chi^{2}(p)$. The test based on $n R^{2}$ is called the BreuschGodfrey test for serial correlation. When $p=1$, the test statistic is asymptotically equivalent to the square of what is known as Durbin's $h$ statistic. ${ }^{17}$

\section*{QUESTIONS FOR REVIEW}
\begin{enumerate}
  \item (Consistency of sample autocovariances) Show: if $\left\{z_{t}\right\}$ is ergodic stationary, then the sample autocovariance $\hat{\gamma}_{j}$ given in (2.10.1) is consistent for $\gamma_{j}$, the population autocovariance. Hint: $\gamma_{j}=\mathrm{E}\left(z_{t} z_{t-j}\right)-\mathrm{E}\left(z_{t}\right) \mathrm{E}\left(z_{t-j}\right)$. Rewrite $\hat{\gamma}_{j}$ as
\end{enumerate}

$$
\hat{\gamma}_{j}=\frac{1}{n} \sum_{t=j+1}^{n} z_{t} z_{t-j}-\bar{z}_{n} \frac{1}{n} \sum_{t=j+1}^{n} z_{t-j}-\bar{z}_{n} \frac{1}{n} \sum_{t=j+1}^{n} z_{t}+\frac{n-j}{n}\left(\bar{z}_{n}\right)^{2} .
$$

\begin{enumerate}
  \setcounter{enumi}{1}
  \item (Asymptotic equivalence of two $Q$ 's) Prove that the difference between the Box-Pierce $Q$ and the Ljung-Box $Q$ converges to 0 in probability as $n \rightarrow \infty$ (so their asymptotic distributions are the same). Hint: Let
\end{enumerate}

$$
\mathbf{x}_{n} \equiv\left(\left(\sqrt{n} \hat{\rho}_{1}\right)^{2}, \ldots,\left(\sqrt{n} \hat{\rho}_{p}\right)^{2}\right)^{\prime}
$$

Find a $p$-dimensional vector $\mathbf{a}_{n}$ such that the difference between two $Q$ 's is $\mathbf{a}_{n}^{\prime} \mathbf{x}_{n}$. Show that $\mathbf{a}_{n} \rightarrow \mathbf{0}$. The asymptotic property to use is Lemma 2.4(b).\\
3. Consider $\sqrt{n} \tilde{\rho}_{j}$ and $\sqrt{n} \hat{\rho}_{j}$ in (2.10.11). We have shown that $\hat{\gamma}_{0} \rightarrow_{\mathrm{p}} \sigma^{2}$ and $\tilde{\gamma}_{0} \rightarrow_{\mathrm{p}} \sigma^{2}$. By Proposition 2.9, $\sqrt{n} \tilde{\gamma}_{j} \rightarrow_{\mathrm{d}} N\left(0, \sigma^{4}\right)$. Taking these for granted, show that $\sqrt{n} \tilde{\rho}_{j}-\sqrt{n} \hat{\rho}_{j} \rightarrow_{\mathrm{p}} 0$ if $\sqrt{n} \tilde{\gamma}_{j}-\sqrt{n} \hat{\gamma}_{j} \rightarrow_{\mathrm{p}} 0$. Hint: If $\sqrt{n} \tilde{\gamma}_{j}-\sqrt{n} \hat{\gamma}_{j} \rightarrow_{\mathrm{p}}$ 0 , then by Lemma 2.4(a), $\sqrt{n} \hat{\gamma}_{j} \rightarrow_{\mathrm{d}} N\left(0, \sigma^{4}\right)$.

$$
\frac{\tilde{\gamma}_{j}}{\tilde{\gamma}_{0}}-\frac{\hat{\gamma}_{j}}{\hat{\gamma}_{0}}=\tilde{\gamma}_{j} \cdot\left(\frac{1}{\tilde{\gamma}_{0}}-\frac{1}{\sigma^{2}}\right)-\hat{\gamma}_{j} \cdot\left(\frac{1}{\hat{\gamma}_{0}}-\frac{1}{\sigma^{2}}\right)+\frac{1}{\sigma^{2}} \cdot\left(\tilde{\gamma}_{j}-\hat{\gamma}_{j}\right) .
$$

Multiply both sides by $\sqrt{n}$ and apply Lemma 2.4(b).

\footnotetext{${ }^{17}$ See Breusch (1978) and Godfrey (1978). The Breusch-Godfrey test was originally derived as a Lagrange Multiplier test for the case where the error is normally distributed and $\mathbf{x}_{t}$ consists of fixed regressors and the lagged dependent variable. The discussion in the text shows that the test is applicable to the more general case considered in Proposition 2.10.
}\subsection*{2.11 Application: Rational Expectations Econometrics}
In dynamic economic models and other contexts, expectations about the future naturally play an important role. Rational expectations econometrics concerns estimating equations involving expectations. Expectations are usually not observable, but, if we assume that economic agents form expectations rationally, we can overcome the problem of unobservable expectations using the techniques of this chapter. The application we consider is Fama's efficient market hypothesis. In his own words, "An efficient capital market is a market that is efficient in processing information. . . . In an efficient market, prices 'fully reflect' available information" (Fama, 1976, p. 133). The words "fully reflect" can be formalized precisely. We will do so for a particular capital market and test its implication.

\section*{The Efficient Market Hypotheses}
The capital market studied in Fama (1975) is the market for U.S. Treasury bills. In this section we focus on the one-month Treasury bill rates observed on a monthly basis. (Later in the book we will also consider Treasury bills of different maturities.) To formalize market efficiency for the bills market, we need to introduce a few concepts from macroeconomics and finance. Since in this section we deal exclusively with time-series data, we use " $t$ " rather than " $i$ " for the subscript. Define

$$
\begin{aligned}
v_{t} \equiv & \text { price of a one-month Treasury bill at the beginning of month } t, \\
R_{t} \equiv & \text { one-month nominal interest rate over month } t, \text { i.e., nominal return } \\
& \text { on the bill over the month }=\left(1-v_{t}\right) / v_{t}, \text { so } v_{t}=1 /\left(1+R_{t}\right), \\
P_{t} \equiv & \text { value of the Consumer Price Index (CPI), which is our measure of } \\
& \text { the price level at the beginning of month } t, \\
\pi_{t+1} \equiv & \text { inflation rate over month } t \text { (i.e., from the beginning of month } t \text { to } \\
& \text { the beginning of month } t+1)=\left(P_{t+1}-P_{t}\right) / P_{t}, \\
{ }_{t} \pi_{t+1} \equiv & \text { expected inflation rate over month } t, \text { expectation formed at the } \\
& \text { beginning of month } t, \\
\eta_{t+1} \equiv & \text { inflation forecast error }=\pi_{t+1}-{ }_{t} \pi_{t+1}, \\
r_{t+1} \equiv & \text { ex-post real interest rate over month } t=\frac{1 / P_{t+1}-v_{t} / P_{t}}{v_{t} / P_{t}} \\
& =\frac{1+R_{t}}{1+\pi_{t+1}}-1 \approx R_{t}-\pi_{t+1}, \\
t_{t+1} \equiv & \text { ex-ante real interest rate }=\frac{1+R_{t}}{1+{ }_{t} \pi_{t+1}}-1 \approx R_{t}-{ }_{t} \pi_{t+1}
\end{aligned}
$$

\begin{figure}[h]
\begin{center}
  \includegraphics[alt={},max width=\textwidth]{14fb0795-6d81-4e76-9788-792298f1ddae-168_401_1138_242_267}
\captionsetup{labelformat=empty}
\caption{Figure 2.2: What's Observed When}
\end{center}
\end{figure}

Timing - the value of which variable is observed at what point in time-is very important in rational expectations econometrics. Our rule for dating variables (which is different from Fama's) is that the variable has subscript $t$ if its value is first observed at the beginning of period (month) $t$. ${ }^{18}$ For example, because $\pi_{t+1}$, the inflation rate over month $t$, depends on $P_{t+1}$, its subscript is $t+1$ rather than $t$. For the same reason, $r_{t+1}$, the ex-post real interest rate over month $t$ has subscript $t+1$. In contrast, the nominal interest rate $R_{t}$ over month $t$ has subscript $t$, not $t+1$, because it is determined by the price of the Treasury bill at the beginning of month $t, v_{t}$. Figure 2.2 shows the value of which variable is observed at what point in time.

We take a period to be a month, only because our data are monthly. Note that the maturity of the security (one month) coincides with the sampling interval (a month). If the sampling interval is finer than the maturity, which is the case if, for example, we have monthly data on the three-month Treasury bill rate, then maturities overlap and we need a more sophisticated technique to be introduced in Chapter 6.

The efficient market hypothesis is a joint hypothesis combining:\\
Rational Expectations. Inflationary expectations are rational: ${ }_{t} \pi_{t+1}=\mathrm{E}\left(\pi_{t+1} \mid\right. I_{t}$ ), where $I_{t}$ is information available at the beginning of month $t$ and includes $\left\{R_{t}, R_{t-1}, \ldots, \pi_{t}, \pi_{t-1}, \ldots\right\}$. Also, $I_{t} \supseteq I_{t-1} \supseteq I_{t-2} \ldots$. That is, agents participating in the market do not forget.

Constant Real Rates. The ex-ante real interest rate is constant: ${ }_{t} r_{t+1}=r$.

\footnotetext{${ }^{18}$ In Fama's (1975) notation, our $\pi_{t+1}$ is $-\widetilde{\Delta}_{t}, r_{t+1}$ is $\tilde{r}_{t+1}$, and $I_{t}$ is $\phi_{t}$.
}\section*{Testable Implications}
We derive two testable implications of the efficient market hypothesis by utilizing the following key observations about the inflation forecast error under rational expectations.\\
(a) $\mathrm{E}\left(\eta_{t+1} \mid I_{t}\right)=0$, that is, the inflation forecast error is a martingale difference sequence with respect to the information set.\\
(b) $\left\{\eta_{t}, \eta_{t-1}, \ldots\right\}$ as well as $\left\{R_{t}, R_{t-1}, \ldots, \pi_{t}, \pi_{t-1}, \ldots\right\}$ are included in $I_{t}$ (known at the beginning of month $t$ ). That is, agents remember all past mistakes.\\
(a) makes intuitive sense; if people use all available information to forecast the inflation rate, the forecast error, which is known only after the fact, will not have any systematic relation to what people knew when they formed expectations. It can be proved easily:


\begin{align*}
& \mathrm{E}\left(\eta_{t+1} \mid I_{t}\right) \\
& =\mathrm{E}\left(\pi_{t+1}-{ }_{t} \pi_{t+1} \mid I_{t}\right) \quad\left(\text { since } \eta_{t+1} \equiv \pi_{t+1}-{ }_{t} \pi_{t+1}\right) \\
& =\mathrm{E}\left(\pi_{t+1} \mid I_{t}\right)-\mathrm{E}\left(\pi_{t+1} \mid I_{t}\right) \\
& =\mathrm{E}\left(\pi_{t+1} \mid I_{t}\right)-\mathrm{E}\left[\mathrm{E}\left(\pi_{t+1} \mid I_{t}\right) \mid I_{t}\right] \\
& \quad\left(\text { since }{ }_{t} \pi_{t+1}=\mathrm{E}\left(\pi_{t+1} \mid I_{t}\right) \text { by rational expectations }\right) \\
& =\mathrm{E}\left(\pi_{t+1} \mid I_{t}\right)-\mathrm{E}\left(\pi_{t+1} \mid I_{t}\right)=0 \tag{2.11.1}
\end{align*}


(b) follows because agents remember past inflation forecasts. Since ${ }_{t-j-1} \pi_{t-j}= \mathrm{E}\left(\pi_{t-j} \mid I_{t-j-1}\right)$ is a function of $I_{t-j-1}$, it is included in $I_{t-j-1}$ (i.e., known at the beginning of month $t-j-1$ ) and hence in $I_{t}$ (known at the beginning of month $t)$ for $j \geq 0$. Thus $\eta_{t-j} \equiv \pi_{t-j}-t_{t-j-1} \pi_{t-j}$ is included in $I_{t}$ for all $j \geq 0$. Observations (a) and (b), together with the Law of Total Expectations, imply\\
(c) $\left\{\eta_{t}\right\}$ is m.d.s.

So, it is a zero mean and serially uncorrelated process.\\
These are not testable because expectations are unobservable, but, when combined with the constant-real-rate assumption, they imply two testable implications about the inflation rate and the interest rate which are observable.

Implication 1: The ex-post real interest rate has a constant mean and is serially uncorrelated.

This follows because


\begin{align*}
r_{t+1} & \equiv R_{t}-\pi_{t+1} \\
& =\left(R_{t}-{ }_{t} \pi_{t+1}\right)-\left(\pi_{t+1}-{ }_{t} \pi_{t+1}\right) \\
& ={ }_{t} r_{t+1}-\eta_{t+1} \quad \text { (by definition) } \\
& =r-\eta_{t+1} \quad \text { (since }{ }_{t} r_{t+1}=r \text { by assumption) } \tag{2.11.2}
\end{align*}


By (c), $\left\{r_{t}\right\}$ has mean $r$ and is serially uncorrelated.

Implication 2: $\mathrm{E}\left(\pi_{t+1} \mid I_{t}\right)=-r+R_{t}$.

That is, the best (in the sense of minimum mean squared error) inflation forecast on the basis of all the available information. is the nominal interest rate; the nominal rate $R_{t}$, which is determined by the price of a Treasury bill $v_{t}$, summarizes all the currently available information relevant for predicting future inflation. This is the formalization of asset prices "fully reflecting" available information. This, too, can be derived easily. Solve (2.11.2) for $\pi_{t+1}$ as $\pi_{t+1}=-r+R_{t}+\eta_{t+1}$. So


\begin{align*}
& \mathrm{E}\left(\pi_{t+1} \mid I_{t}\right) \\
& =\mathrm{E}\left(-r+R_{t}+\eta_{t+1} \mid I_{t}\right) \\
& =-r+R_{t}+\mathrm{E}\left(\eta_{t+1} \mid I_{t}\right) \quad\left(\text { since } r \text { is a constant and } R_{t} \text { is included in } I_{t}\right) \\
& =-r+R_{t} \quad(\text { by (a) }) \tag{2.11.3}
\end{align*}


In the rest of this section, we test these two implications of market efficiency.

\section*{Testing for Serial Correlation}
Consider first the implication that the ex-post real rate has no serial correlation, which Fama tests using the result from Proposition 2.9. We use the monthly data set used by Mishkin (1992) on the one-month $T$-bill rate and the monthly CPI inflation rate, stated in percent at annual rates. ${ }^{19}$ Figure 2.3 plots the data. The ex-post real interest rate, defined as the difference between the T-bill rate and the CPI inflation rate, is plotted in Figure 2.4. To duplicate Fama's results, we take the sample period to be the same as in Fama (1975), which is January 1953 through July 1971. The sample size is thus 223 . The time-series properties of the real interest rate are summarized in Table 2.1. In calculating the sample autocorrelations $\hat{\rho}_{j}=\hat{\gamma}_{j} / \hat{\gamma}_{0}$,

\footnotetext{${ }^{19}$ The data set is described in some detail in the empirical exercise. Following Fama, the inflation measure used in the text is calculated from the raw CPI numbers. The inflation measure used by Mishkin (1992) treats the housing component of the CPI consistently. See question ( $k$ ) of Empirical Exercise 1.
}\begin{figure}[h]
\begin{center}
  \includegraphics[alt={},max width=\textwidth]{14fb0795-6d81-4e76-9788-792298f1ddae-171_760_1160_359_311}
\captionsetup{labelformat=empty}
\caption{Figure 2.3: Inflation and Interest Rates}
\end{center}
\end{figure}

\begin{figure}[h]
\begin{center}
  \includegraphics[alt={},max width=\textwidth]{14fb0795-6d81-4e76-9788-792298f1ddae-171_829_1139_1303_320}
\captionsetup{labelformat=empty}
\caption{Figure 2.4: Real Interest Rates}
\end{center}
\end{figure}

\begin{table}[h]
\begin{center}
\captionsetup{labelformat=empty}
\caption{Table 2.1: Real Interest Rates, January 1953-July 1971}
\begin{tabular}{|l|l|l|l|l|l|l|l|l|l|l|l|l|}
\hline
\multirow{2}{*}{} & \multicolumn{12}{|c|}{mean $=0.82 \%$, standard deviation $=2.847 \%$, sample size $=223$} \\
\hline
 & $j=1$ & $j=2$ & $j=3$ & $j=4$ & $j=5$ & $j=6$ & $j=7$ & $j=8$. & $j=9$ & $j=10$ & $j=11$ & $j=12$ \\
\hline
$\hat{\rho}_{j}$ & -0.101 & 0.172 & -0.019 & -0.004 & -0.064 & -0.021 & -0.092 & 0.095 & 0.094 & 0.019 & 0.004 & 0.207 \\
\hline
std. error & 0.067 & 0.067 & 0.067 & 0.067 & 0.067 & 0.067 & 0.067 & 0.067 & 0.067 & 0.067 & 0.067 & 0.067 \\
\hline
Ljung-Box $Q$ & 2.3 & 9.1 & 9.1 & 9.1 & 10.1 & 10.2 & 12.1 & 14.2 & 16.3 & 16.4 & 16.4 & 26.5 \\
\hline
$p$-value (\%) & 12.8\% & 1.1\% & 2.8\% & 5.8\% & 7.3\% & 11.7\% & 9.6\% & 7.6\% & 6.1\% & 8.9\% & 12.8\% & 0.9\% \\
\hline
\end{tabular}
\end{center}
\end{table}

we use the formula (2.10.1) for $\hat{\gamma}_{j}$ where the denominator is the sample size rather than $n-j$. Similarly, the formula for the standard error of $\hat{\rho}_{j}$ is $1 / \sqrt{n}$ rather than $1 / \sqrt{n-j}$ (so the standard error does not depend on $j$ ). The Ljung-Box $Q$ statistic in the table for, say, $j=2$ is (2.10.5) with $p=2$. The results in the table show that none of the autocorrelations are highly significant individually or as a group, in accordance with the first implication of the efficient market hypothesis.

\section*{Is the Nominal Interest Rate the Optimal Predictor?}
We can test Implication 2 by estimating an appropriate regression and carrying out the $t$-test of Proposition 2.3. The remarkable fact about the efficient market hypothesis is that, despite its simplicity, it is specific enough to imply the important parts of the assumptions justifying the $t$-test. ${ }^{20}$ To see this, let $y_{t} \equiv \pi_{t+1}, \mathbf{x}_{t} \equiv \left(1, R_{t}\right)^{\prime}, \varepsilon_{t} \equiv \eta_{t+1}$ (the inflation forecast error), and rewrite (2.11.2) as


\begin{equation*}
\pi_{t+1}=\beta_{1}+\beta_{2} R_{t}+\eta_{t+1} \text { or } y_{t}=\mathbf{x}_{t}^{\prime} \boldsymbol{\beta}+\varepsilon_{t}, \tag{2.11.4}
\end{equation*}


so Assumption 2.1 is obviously satisfied with


\begin{equation*}
\beta=(-r, 1)^{\prime} . \tag{2.11.5}
\end{equation*}


The other assumptions of the model are verified as follows.

\begin{itemize}
  \item Assumption 2.3 (predetermined regressors). The $\mathbf{g}_{t}\left(\equiv \mathbf{x}_{t} \cdot \varepsilon_{t}\right)$ in the present context is $\left(\eta_{t+1}, R_{t} \eta_{t+1}\right)^{\prime}$. What we need to show is that $\mathrm{E}\left(\eta_{t+1}\right)=0$ and $\mathrm{E}\left(R_{t} \eta_{t+1}\right)=0$. The former is implied by (a) and the Law of Total Expectations:
\end{itemize}


\begin{equation*}
\mathrm{E}\left(\eta_{t+1}\right)=\mathrm{E}\left[\mathrm{E}\left(\eta_{t+1} \mid I_{t}\right)\right]=0 . \tag{2.11.6}
\end{equation*}


The latter holds because


\begin{align*}
& \mathrm{E}\left(R_{t} \eta_{t+1}\right) \\
& =\mathrm{E}\left[\mathrm{E}\left(R_{t} \eta_{t+1} \mid I_{t}\right)\right] \quad \text { (by the Law of Total Expectations) } \\
& \left.=\mathrm{E}\left[R_{t} \mathrm{E}\left(\eta_{t+1} \mid I_{t}\right)\right] \quad \text { (by linearity of conditional expectations and } R_{t} \in I_{t}\right) \\
& =0 \quad(\text { by }(\text { a })) . \tag{2.11.7}
\end{align*}


\begin{itemize}
  \item Assumption 2.5. By (b), $I_{t}$ includes $\left(\varepsilon_{t-1}\left(=\eta_{t}\right), \varepsilon_{t-2}, \ldots\right)$ as well as $\left(\mathbf{x}_{t}\right.$, $\mathbf{x}_{t-1}, \ldots$ ). (a) then implies (2.3.1), which is a sufficient condition for $\left\{\mathbf{g}_{t}\right\}$ to be m.d.s.
\end{itemize}

\footnotetext{${ }^{20}$ The technical part - that the fourth-moment matrix $\mathbf{E}\left(\mathbf{g}_{t} \mathbf{g}_{t}^{\prime}\right)$ exists and is finite and that the regressors have finite fourth moments - is not implied by the efficient market hypothesis and needs to be assumed.
}\begin{itemize}
  \item Assumption 2.4. $\mathrm{E}\left(\mathbf{x}_{t} \mathbf{x}_{t}^{\prime}\right)$ here is
\end{itemize}

\[
\mathrm{E}\left(\mathbf{x}_{t} \mathbf{x}_{t}^{\prime}\right)=\left[\begin{array}{cc}
1 & \mathrm{E}\left(R_{t}\right)  \tag{2.11.8}\\
\mathrm{E}\left(R_{t}\right) & \mathrm{E}\left(R_{t}^{2}\right)
\end{array}\right]
\]

Its determinant is $\mathrm{E}\left(R_{t}^{2}\right)-\left[\mathrm{E}\left(R_{t}\right)\right]^{2}=\operatorname{Var}\left(R_{t}\right)$, which must be positive; if the variance were zero, we would not observe the interest rate $R_{t}$ fluctuating.

\begin{itemize}
  \item Assumption 2.2 (ergodic stationarity). It requires $\left\{\pi_{t+1}, R_{t}\right\}$ to be ergodic stationary. The plot in Figure 2.3 shows a stretch of upward movements, especially for the nominal interest rate, although the stretch is followed by a reversion to a lower level. Testing stationarity of the series will be covered in Example 9.2 of Chapter 9. For now, despite this rather casual evidence to the contrary, we proceed under the assumption of stationarity (but the implication of having trending series will be touched upon in the final subsection). For ergodicity, we just have to assume it.
\end{itemize}

What we have shown is that, under the efficient market hypothesis, $\left\{y_{t}, \mathbf{x}_{t}\right\}$ (where $y_{t}=\pi_{t+1}$ and $\left.\mathbf{x}_{t}=\left(1, R_{t}\right)^{\prime}\right)$ belongs to the set of DGPs satisfying Assumptions 2.1-2.5.

Having verified the required assumptions, we now proceed to estimation. Using our data, the estimated regression is, for $t=1 / 53, \ldots, 7 / 71$,


\begin{gather*}
\pi_{t+1}=\underset{(0.431)}{-0.868}+\underset{(0.112)}{1.015} R_{t}, \\
R^{2}=0.24, \text { mean of dependent variable }=2.35, S E R=2.84 \%, n=223 . \tag{2.11.9}
\end{gather*}


The heteroskedasticity-robust standard errors calculated as in (2.4.1) (where $\widehat{\mathbf{S}}$ is given by (2.5.1) with $\hat{\varepsilon}_{i}=e_{i}$ ) are shown in parentheses. As shown in (2.11.5), market efficiency implies that the $R_{t}$ coefficient equals $1 .{ }^{21}$ The robust $t$-ratio for the null that the $R_{t}$ coefficient equals 1 is $(1.015-1) / 0.112=0.13$. Its $p$-value is about 90 percent. Thus, we can easily accept the null hypothesis.

There are other ways to test the implication of the nominal interest rate summarizing all the currently available information. We can bring in more explanatory variables in the regression and see if they have an explanatory power over and above the nominal rate. The additional variables to be brought in, however, must be part of the current information set $I_{t}$ because, obviously, the efficient market

\footnotetext{${ }^{21}$ The intercept is known to be $-r$, but it does not produce a testable restriction since the efficient market hypothesis does not specify the value of $r$.
}
hypothesis is consistent with the existence of some variable not currently available predicting future inflation better than the current nominal rate. Formally, if $\mathbf{x}_{t}$ includes variables not included in $I_{t}$, the argument we used for verifying Assumption 2.3 is no longer valid.

\section*{$\boldsymbol{R}_{\boldsymbol{t}}$ Is Not Strictly Exogenous}
As has been emphasized, one important advantage of large-sample theory over finite-sample theory of Chapter 1 is that the regressors do not have to be strictly exogenous as long as they are orthogonal to the error term. We now show by example that $R_{t}$ is not strictly exogenous, so finite-sample theory is not applicable to the Fama regression (2.11.4). Suppose that the process for the inflation rate is an AR(1) process:

$$
\pi_{t}=c+\rho \pi_{t-1}+\eta_{t}, \quad\left\{\eta_{t}\right\} \text { is independent white noise. }
$$

If $\eta_{t+1}$ is independent of any elements of $I_{t}$, then

$$
\begin{aligned}
& \mathrm{E}\left(\pi_{t+1} \mid I_{t}\right) \\
& =c+\rho \pi_{t}+\mathrm{E}\left(\eta_{t+1} \mid I_{t}\right) \quad\left(\text { since } \pi_{t} \text { is in } I_{t}\right) \\
& =c+\rho \pi_{t} \quad\left(\text { since } \mathrm{E}\left(\eta_{t+1} \mid I_{t}\right)=\mathrm{E}\left(\eta_{t+1}\right) \text { and } \mathrm{E}\left(\eta_{t+1}\right)=0 \text { by assumption }\right) .
\end{aligned}
$$

So $\eta_{t+1}$ is indeed the inflation forecast error for $\pi_{t+1}$ and

$$
R_{t}=r+c+\rho \pi_{t}
$$

(since $R_{t}=r+\mathrm{E}\left(\pi_{t+1} \mid I_{t}\right)$ under market efficiency). It is then easy to see that the error term $\varepsilon_{t}$ in the Fama regression (of $\pi_{t+1}$ on a constant and $R_{t}$ ), which is the inflation forecast error $\eta_{t+1}$, can be correlated with future regressors. For example,

$$
\begin{aligned}
& \operatorname{Cov}\left(\varepsilon_{t}, R_{t+1}\right) \\
& =\operatorname{Cov}\left(\eta_{t+1}, r+c+\rho \pi_{t+1}\right) \\
& =\rho \operatorname{Cov}\left(\eta_{t+1}, \pi_{t+1}\right) \\
& =\rho \operatorname{Var}\left(\eta_{t+1}\right) \quad\left(\text { since } \pi_{t+1}=c+\rho \pi_{t}+\eta_{t+1} \text { and } \operatorname{Cov}\left(\eta_{t+1}, \pi_{t}\right)=0\right),
\end{aligned}
$$

which is not zero.\\
Although finite-sample theory is not applicable, its prescription for statistical inference is approximately valid, provided that the error is conditionally homoskedastic and the sample size is large enough, which was the message of Section 2.6.

\section*{Subsequent Developments}
Although Fama's paper represents perhaps the finest example of rational expectations econometrics, events after its publication and a large number of empirical studies it inspired proved that the near one-for-one relation between inflation and interest rates is limited to the postwar period until 1979. Such strong association cannot be found for the prewar period or for many other countries. Furthermore, the increase in the real interest rates after the change in the Fed's operating procedure that took place in October 1979 runs counter to the premise of the Fama regression that the expected real interest rate is constant over time. When the sample is restricted to the post-October 1979 period, the interest rate coefficient in the Fama regression is far below unity. ${ }^{22}$

These findings raise the question of why the strong association occurs only for certain periods but not for others. The explanation by Mishkin (1992) is that the inflation premium gets incorporated into interest rates only gradually. In periods when the inflation rate shows only short-run fluctuations, interest rates are not responsive to inflation. However, sustained movements in inflation will be reflected in interest rates. The period when the strong association is observed (which includes Fama's sample period, see Figure 2.3) is precisely when inflation showed a sustained upward movement.

On January 29, 1997, the U.S. Treasury auctioned \$7 billion in ten-year inflation-indexed bonds, making it possible for researchers to observe the ex-ante real interest rate as the yield on indexed bonds. Evidence from Great Britain, where indexed bonds have been available since the early 1980s, is that the yields on indexed bonds, although much less volatile than the yield on ordinary bonds, are not constant over time. The constant-real-rate assumption that we made was the auxiliary assumption whereby we can test whether bond prices fully reflect available information. Now we can do so without the auxiliary assumption by regressing the actual inflation rate on the yield differential between the conventional and indexed bonds. We can also relax the assumption, made implicitly so far, that the yield differential equals the expected inflation rate. Investors may not be risk-neutral, and the inflation risk premium over and above the expected inflation rate may be needed to induce them to hold ordinary bonds. If so, the expected inflation rate is once again unobservable. Rational expectations econometrics will continue to be a useful tool for dealing with unobservable expectations.

\footnotetext{${ }^{22}$ See, for example, Mishkin (1992, Table 1). We will verify this in part (1) of Empirical Exercise 1.
}\section*{QUESTIONS FOR REVIEW}
\begin{enumerate}
  \item Let $R_{t}=5 \%$ and $\pi_{t+1}=2 \%$. Calculate $r_{t+1}$ using the two formulas, ( $1+ \left.R_{t}\right) /\left(1+\pi_{t+1}\right)-1$ and $R_{t}-\pi_{t+1}$. Is the difference small? What if $R_{t}=20 \%$ and $\pi_{t+1}=17 \%$ ?
  \item Show that $\widehat{\mathrm{E}}^{*}\left(\pi_{t+1} \mid 1, R_{t}\right)=-r+R_{t}$ under the efficient market hypothesis.
\end{enumerate}

Hint: Use (2.9.7).\\
3. If the inflation forecast error $\eta_{t+1}$ were observable, how would you test market efficiency? Do we need the constant ex-ante real rate assumption?\\
4. If the inflation rates and the interest rates were measured in fractions rather than percent (e.g., 0.08 instead of $8 \%$ ), how would the regression result (2.11.9) change? If the inflation rate is in percent but per month and the interest rate is in percent per year? Hint: If $x$ is the inflation rate per month and $y$ is its annual rate, then $1+y=(1+x)^{12}$ so that $y \approx 12 x$.\\
5. Suppose in (2.11.4) $\mathbf{x}_{t}$ includes a third variable which is not in $I_{t}$. Which part of the argument deriving Assumption 2.3 fails? Hint: The third element of $\mathbf{g}_{t}$ is the third variable times $\eta_{t+1}$.\\
6. Provide an example of market inefficiency such that Implication 1 holds but Implication 2 does not. Hint: Suppose $\pi_{t}$ and $R_{t}$ are serially independent and mutually independent processes.

\subsection*{2.12 Time Regressions}
Throughout this chapter, we have assumed that $\left\{y_{t}, \mathbf{x}_{t}\right\}$ is stationary. This assumption, however, is not always satisfied in time-series regressions. In this book, we examine two cases where the regressors are not stationary and yet OLS is applicable. The first is examined here, while the other case, called "cointegrating regressions," is covered in Chapter 10.

The regression we consider is written as


\begin{equation*}
y_{t}=\alpha+\delta \cdot t+\varepsilon_{t}, \tag{2.12.1}
\end{equation*}


where $\left\{\varepsilon_{t}\right\}$ is independent white noise. This regression can be written as $y_{t}= \mathbf{x}_{t}^{\prime} \boldsymbol{\beta}+\varepsilon_{t}$ with


\begin{equation*}
\mathbf{x}_{t}=(1, t)^{\prime}, \boldsymbol{\beta}=(\alpha, \delta)^{\prime} . \tag{2.12.2}
\end{equation*}


Clearly, $\mathbf{x}_{t}$ is not stationary because the mean of the second element is $t$, which increases with time. Similarly, $y_{t}$ is not stationary. The linear function, $\alpha+\delta \cdot t$, is called the time trend of $y_{t}$. We say that a process is trend stationary if it can be written as the sum of a time trend and a stationary process. The process $\left\{y_{t}\right\}$ here is a special trend-stationary process where the stationary component is independent white noise.

\section*{The Asymptotic Distribution of the OLS Estimator}
Let $\mathbf{b}$ be the OLS estimate of $\boldsymbol{\beta}$ based on a sample of size $n$ :

\[
\mathbf{b} \equiv\left[\begin{array}{c}
\hat{\alpha}  \tag{2.12.3}\\
\hat{\delta}
\end{array}\right]=\left(\sum_{t=1}^{n} \mathbf{x}_{t} \mathbf{x}_{t}^{\prime}\right)^{-1}\left(\sum_{t=1}^{n} \mathbf{x}_{t} \cdot y_{t}\right) .
\]

The sampling error can be written as

\[
\mathbf{b}-\boldsymbol{\beta}=\left[\begin{array}{c}
\hat{\alpha}-\alpha  \tag{2.12.4}\\
\hat{\delta}-\delta
\end{array}\right]=\left(\sum_{t=1}^{n} \mathbf{x}_{t} \mathbf{x}_{t}^{\prime}\right)^{-1}\left(\sum_{t=1}^{n} \mathbf{x}_{t} \cdot \varepsilon_{t}\right) .
\]

Using the algebraic result that $\sum_{t=1}^{n} t=n \cdot(n+1) / 2, \sum_{t=1}^{n} t^{2}=n \cdot(n+1)(2 n+$ 1)/6, it is easy to show

\[
\sum_{t=1}^{n} \mathbf{x}_{t} \mathbf{x}_{t}^{\prime}=\left[\begin{array}{cc}
n & n \cdot(n+1) / 2  \tag{2.12.5}\\
n \cdot(n+1) / 2 & n \cdot(n+1)(2 n+1) / 6
\end{array}\right] .
\]

So, unlike in the stationary case, $\sum_{t=1}^{n} \mathbf{x}_{t} \mathbf{x}_{t}^{\prime} / n$ does not converge in probability to a nonsingular matrix; it actually diverges.

It turns out that the OLS estimates $\hat{\alpha}$ and $\hat{\delta}$ are consistent but have different rates of convergence. As in the stationary case, the rate of convergence for $\hat{\alpha}$ is $\sqrt{n}$. In contrast, the rate for $\hat{\delta}$ is $n^{3 / 2}$. That is, $n^{3 / 2}(\hat{\delta}-\delta)$ converges to a nondegenerate distribution (the distribution of a nonconstant random variable). To assign those different rates of convergence to the elements of $\mathbf{b}-\boldsymbol{\beta}$, consider multiplying both sides of (2.12.4) by the matrix

\[
\mathbf{\Upsilon}_{n}=\left[\begin{array}{cc}
\sqrt{n} & 0  \tag{2.12.6}\\
0 & n^{3 / 2}
\end{array}\right]
\]

This produces


\begin{align*}
\boldsymbol{\Upsilon}_{n}(\mathbf{b}-\boldsymbol{\beta}) & =\left[\begin{array}{c}
\sqrt{n}(\hat{\alpha}-\alpha) \\
n^{3 / 2}(\hat{\delta}-\delta)
\end{array}\right]=\boldsymbol{\Upsilon}_{n}\left(\sum_{t=1}^{n} \mathbf{x}_{t} \mathbf{x}_{t}^{\prime}\right)^{-1}\left(\sum_{t=1}^{n} \mathbf{x}_{t} \cdot \varepsilon_{t}\right) \\
& =\boldsymbol{\Upsilon}_{n}\left(\sum_{t=1}^{n} \mathbf{x}_{t} \mathbf{x}_{t}^{\prime}\right)^{-1} \boldsymbol{\Upsilon}_{n} \boldsymbol{\Upsilon}_{n}^{-1}\left(\sum_{t=1}^{n} \mathbf{x}_{t} \cdot \varepsilon_{t}\right) \\
& =\left[\boldsymbol{\Upsilon}_{n}^{-1}\left(\sum_{t=1}^{n} \mathbf{x}_{t} \mathbf{x}_{t}^{\prime}\right) \boldsymbol{\Upsilon}_{n}^{-1}\right]^{-1}\left(\boldsymbol{\Upsilon}_{n}^{-1} \sum_{t=1}^{n} \mathbf{x}_{t} \cdot \varepsilon_{t}\right) \\
& \equiv \mathbf{Q}_{n}^{-1} \mathbf{v}_{n} \tag{2.12.7}
\end{align*}


where


\begin{equation*}
\mathbf{Q}_{n} \equiv \mathbf{\Upsilon}_{n}^{-1}\left(\sum_{t=1}^{n} \mathbf{x}_{t} \mathbf{x}_{t}^{\prime}\right) \mathbf{\Upsilon}_{n}^{-1} \quad \text { and } \quad \mathbf{v}_{n} \equiv \mathbf{\Upsilon}_{n}^{-1} \sum_{t=1}^{n} \mathbf{x}_{t} \cdot \varepsilon_{t} \tag{2.12.8}
\end{equation*}


Substituting (2.12.5) and (2.12.6) into these expressions for $\mathbf{Q}_{n}$ and $\mathbf{v}_{n}$, we obtain


\begin{gather*}
\mathbf{Q}_{n}=\left[\begin{array}{cc}
1 & (n+1) /(2 n) \\
(n+1) /(2 n) & (n+1)(2 n+1) /\left(6 n^{2}\right)
\end{array}\right], \\
\mathbf{v}_{n}=\left[\begin{array}{c}
\frac{1}{\sqrt{n}} \sum_{t=1}^{n} \varepsilon_{t} \\
\frac{1}{\sqrt{n}} \sum_{t=1}^{n}(t / n) \varepsilon_{t}
\end{array}\right] . \tag{2.12.9}
\end{gather*}


Clearly, we have

\[
\mathbf{Q}_{n} \rightarrow \mathbf{Q} \equiv\left[\begin{array}{cc}
1 & 1 / 2  \tag{2.12.10}\\
1 / 2 & 1 / 3
\end{array}\right]
\]

For $\mathbf{v}_{n}$, it can be shown (see, e.g., Hamilton, 1994, pp. 458-460) that


\begin{equation*}
\mathbf{v}_{n} \underset{\mathrm{~d}}{\rightarrow} N\left(\mathbf{0}, \sigma^{2} \mathbf{Q}\right) \tag{2.12.11}
\end{equation*}


if $\left\{\varepsilon_{t}\right\}$ is independent white noise with $\mathrm{E}\left(\varepsilon_{t}^{2}\right)=\sigma^{2}$ and $\mathrm{E}\left(\varepsilon_{t}^{4}\right)<\infty$. Thus the asymptotic distribution of (2.12.7) is normal with mean $\mathbf{0}$ and variance $\mathbf{Q}^{-1}\left(\sigma^{2} \mathbf{Q}\right)$. $\mathbf{Q}^{-1}=\sigma^{2} \mathbf{Q}^{-1}$. Summarizing the discussion so far,

Proposition 2.11 (OLS estimation of the time regression): Consider the time regression (2.12.1) where $\varepsilon_{t}$ is independent white noise with $\mathrm{E}\left(\varepsilon_{t}^{2}\right)=\sigma^{2}$ and\\
$\mathrm{E}\left(\varepsilon_{t}^{4}\right)<\infty$ and let $\hat{\alpha}$ and $\hat{\delta}$ be the OLS estimate of $\alpha$ and $\delta$. Then

$$
\left[\begin{array}{c}
\sqrt{n}(\hat{\alpha}-\alpha) \\
n^{3 / 2}(\hat{\delta}-\delta)
\end{array}\right] \rightarrow N\left(\mathbf{0}, \sigma^{2} \cdot\left[\begin{array}{cc}
1 & 1 / 2 \\
1 / 2 & 1 / 3
\end{array}\right]^{-1}\right) .
$$

As in the stationary case, $\hat{\alpha}$ is $\sqrt{n}$-consistent because $\sqrt{n}(\hat{\alpha}-\alpha)$ converges to a (normal) random variable. The OLS estimate of the time coefficient, $\hat{\delta}$, too is consistent, but the speed of convergence is faster: it is $n^{3 / 2}$-consistent in that $n^{3 / 2}(\hat{\delta}-\delta)$ converge to a random variable. In this sense, $\hat{\delta}$ is hyperconsistent. ${ }^{23}$

\section*{Hypothesis Testing for Time Regressions}
Thus, the OLS coefficient estimates of the time regression are asymptotically normal, provided the sampling error is properly scaled. We now show that the deflation of the sampling error by the standard error provides a scaling that makes the resulting ratio - the $t$-value - asymptotically standard normal.

First consider the $t$-value for the null $\alpha=\alpha_{0}$. Noting that

\[
(1,1) \text { element of }\left(\sum_{t=1}^{n} \mathbf{x}_{t} \mathbf{x}_{t}^{\prime}\right)^{-1}=\left[\begin{array}{ll}
1 & 0
\end{array}\right]\left(\sum_{t=1}^{n} \mathbf{x}_{t} \mathbf{x}_{t}^{\prime}\right)^{-1}\left[\begin{array}{l}
1  \tag{2.12.12}\\
0
\end{array}\right],
\]

the $t$-value can be written as


\begin{align*}
t & =\frac{\hat{\alpha}-\alpha_{0}}{\sqrt{s^{2} \cdot\left[\begin{array}{ll}
1 & 0
\end{array}\right]\left(\sum_{t=1}^{n} \mathbf{x}_{t} \mathbf{x}_{t}^{\prime}\right)^{-1}\left[\begin{array}{l}
1 \\
0
\end{array}\right]}} \\
& =\frac{\sqrt{n}\left(\hat{\alpha}-\alpha_{0}\right)}{\sqrt{s^{2} \cdot\left[\begin{array}{ll}
\sqrt{n} & 0]\left(\sum_{t=1}^{n} \mathbf{x}_{t} \mathbf{x}_{t}^{\prime}\right)^{-1}\left[\begin{array}{c}
\sqrt{n} \\
0
\end{array}\right] \\
\sqrt{n}\left(\hat{\alpha}-\alpha_{0}\right)
\end{array}\right.}} \\
& =\frac{\sqrt{\sqrt{s^{2} \cdot\left[\begin{array}{ll}
1 & 0
\end{array}\right] \Upsilon_{n}\left(\sum_{t=1}^{n} \mathbf{x}_{t} \mathbf{x}_{t}^{\prime}\right)^{-1} \Upsilon_{n}\left[\begin{array}{l}
1 \\
0
\end{array}\right]}}}{}\left(\text { since }\left[\begin{array}{ll}
\sqrt{n} & 0
\end{array}\right]=\left[\begin{array}{ll}
1 & 0
\end{array}\right] \Upsilon_{n}\right) \\
& =\frac{\sqrt{n}\left(\hat{\alpha}-\alpha_{0}\right)}{\sqrt{s^{2} \cdot\left[\begin{array}{ll}
1 & 0
\end{array}\right]\left\{\Upsilon_{n}^{-1}\left(\sum_{t=1}^{n} \mathbf{x}_{t} \mathbf{x}_{t}^{\prime}\right) \Upsilon_{n}^{-1}\right\}^{-1}\left[\begin{array}{l}
1 \\
0
\end{array}\right]}} \\
& =\frac{\sqrt{n}\left(\hat{\alpha}-\alpha_{0}\right)}{\sqrt{s^{2} \cdot\left[\begin{array}{ll}
1 & 0
\end{array}\right] \mathbf{Q}_{n}^{-1}\left[\begin{array}{l}
1 \\
0
\end{array}\right]}}(\text { by }(2.12 .8)) \tag{2.12.13}
\end{align*}


\footnotetext{${ }^{23}$ Estimators that are $n$-consistent are usually called superconsistent. A prominent example of superconsistent estimators arises in the estimation of unit-root processes; see Chapter 9. Some authors use the term "superconsistent" more broadly, to refer to estimators that are $n^{\gamma}$-consistent with $\gamma>\frac{1}{2}$. In their language, $\hat{\delta}$ in the text is superconsistent.
}It is straightforward (see review question 1) to show that $s^{2}$ is consistent for $\sigma^{2}$. And by (2.12.10), $\mathbf{Q}_{n} \rightarrow_{\mathrm{p}} \mathbf{Q}$. Thus by Lemma 2.4(c) and Proposition 2.11,


\begin{equation*}
t \text {-value for the null of } \alpha=\alpha_{0} \rightarrow \frac{1 \text { st element of } \mathbf{z}}{\sqrt{\sigma^{2} \times(1,1) \text { element of } \mathbf{Q}^{-1}}} \tag{2.12.14}
\end{equation*}


where $\mathbf{z} \sim N\left(\mathbf{0}, \sigma^{2} \mathbf{Q}^{-1}\right)$. So the $t$-value for $\alpha$ is asymptotically $N(0,1)$, as in the stationary case.

Using the same trick and noting that $\left[\begin{array}{ll}0 & n^{3 / 2}\end{array}\right]=\left[\begin{array}{ll}0 & 1\end{array}\right] \Upsilon_{n}$, we can write the $t$-value for the null of $\delta=\delta_{0}$ as

\[
t=\frac{n^{3 / 2}\left(\hat{\delta}-\delta_{0}\right)}{\left.\sqrt{s^{2} \cdot[0} 1\right] \mathbf{Q}_{n}^{-1}\left[\begin{array}{l}
0  \tag{2.12.15}\\
1
\end{array}\right]}
\]

So the $t$-value for the coefficient of time, too, is asymptotically $N(0,1)!$ Inference about $\alpha$ or $\delta$ can be done exactly as in the stationary case.

\section*{QUESTIONS FOR REVIEW}
\begin{enumerate}
  \item ( $s^{2}$ is consistent for $\sigma^{2}$ ) Show that $s^{2} \rightarrow_{\mathrm{p}} \sigma^{2}$. Hint: From (2.3.10), derive
\end{enumerate}

$$
e_{t}^{2}=\varepsilon_{t}^{2}-2(\mathbf{b}-\boldsymbol{\beta})^{\prime} \boldsymbol{\Upsilon}_{n} \boldsymbol{\Upsilon}_{n}^{-1} \mathbf{x}_{t} \cdot \varepsilon_{t}+(\mathbf{b}-\boldsymbol{\beta})^{\prime} \boldsymbol{\Upsilon}_{n} \boldsymbol{\Upsilon}_{n}^{-1} \mathbf{x}_{t} \mathbf{x}_{t}^{\prime} \boldsymbol{\Upsilon}_{n}^{-1} \boldsymbol{\Upsilon}_{n}(\mathbf{b}-\boldsymbol{\beta})
$$

Sum over $t$ to obtain

$$
\frac{1}{n} \sum_{t=1}^{n} e_{t}^{2}=\frac{1}{n} \sum_{t=1}^{n} \varepsilon_{t}^{2}-\frac{2}{n}(\mathbf{b}-\boldsymbol{\beta})^{\prime} \Upsilon_{n} \mathbf{v}_{n}+\frac{1}{n}(\mathbf{b}-\boldsymbol{\beta})^{\prime} \Upsilon_{n} \mathbf{Q}_{n} \Upsilon_{n}(\mathbf{b}-\boldsymbol{\beta})
$$

\begin{enumerate}
  \setcounter{enumi}{1}
  \item (Serially uncorrelated errors) Suppose $\left\{\varepsilon_{t}\right\}$ is a general stationary process, rather than an independent white noise process, but suppose that $\mathbf{v}_{n} \rightarrow_{\mathrm{d}} N(\mathbf{0}, \mathbf{V})$ where $\mathbf{V} \neq \sigma^{2} \mathbf{Q}$. Are the OLS estimators, $\hat{\alpha}$ and $\hat{\delta}$, consistent? Does $\boldsymbol{\Upsilon}_{n}(\mathbf{b}- \boldsymbol{\beta})$ converge to a normal random vector with mean zero? If so, what is the variance of the normal distribution? [Answer: $\mathbf{Q}^{-1} \mathbf{V} \mathbf{Q}^{-1}$.] Are the $t$-values asymptotically normal? [Answer: Yes, but the variance will not be unity.]
\end{enumerate}

\section*{Appendix 2.A: Asymptotics with Fixed Regressors}
To prove results similar in conclusion to Proposition 2.5, we add to Assumption 2.1 the following.

Assumption 2.A.1: $\mathbf{X}$ is deterministic. As $n \rightarrow \infty, \mathbf{S}_{\mathbf{x x}} \rightarrow \boldsymbol{\Sigma}_{\mathbf{x x}}$, a nonsingular matrix (conventional convergence here, because $\mathbf{X}$ is nonrandom).

Assumption 2.A.2: $\left\{\varepsilon_{i}\right\}$ is i.i.d. with $\mathrm{E}\left(\varepsilon_{i}\right)=0$ and $\operatorname{Var}\left(\varepsilon_{i}\right)=\sigma^{2}$.

Proposition 2.A. 1 (asymptotics with fixed regressors): Under Assumptions 2.1, 2.A.1, and 2.A.2, plus an assumption called the Lindeberg condition on $\left\{\mathbf{g}_{i}\right\}$ (where $\left.\mathbf{g}_{i} \equiv \mathbf{x}_{i} \cdot \varepsilon_{i}\right)$, the same conclusions as in Proposition 2.5 hold.

Rather than stating the Lindeberg condition, we explain why such condition is needed. To prove the asymptotic normality of $\mathbf{b}$, we would take (2.3.8) and show that $\sqrt{n} \overline{\mathbf{g}}$ converges to a normal distribution. The technical difficulty is that the sequence $\left\{\mathbf{g}_{i}\right\}$ is not stationary thanks to the assumption that $\left\{\mathbf{x}_{i}\right\}$ is a sequence of nonrandom vectors. For example,

$$
\operatorname{Var}\left(\mathbf{g}_{i}\right)=\mathrm{E}\left(\varepsilon_{i}^{2} \mathbf{x}_{i} \mathbf{x}_{i}^{\prime}\right)=\mathrm{E}\left(\varepsilon_{i}^{2}\right) \mathbf{x}_{i} \mathbf{x}_{i}^{\prime}=\sigma^{2} \mathbf{x}_{i} \mathbf{x}_{i}^{\prime}
$$

is not constant across observations. Consequently, neither the Lindeberg-Levy CLT nor the Martingale Differences CLT is applicable. Fortunately, however, there is a generalization of the Lindeberg-Levy CLT to nonstationary processes with nonconstant variance, called the Lindeberg-Feller CLT, which places a technical restriction on the tail distribution of $\mathbf{g}_{i}$ to prevent observations with large $\mathbf{x}_{i}$ from dominating the distribution of $\sqrt{n} \overline{\mathbf{g}}$.

\section*{Appendix 2.B: Proof of Proposition 2.10}
This appendix provides a proof of Proposition 2.10. We start with the expression for $\hat{\gamma}_{j}$ given in (2.10.12). The first part of the proof is to find an expression that is asymptotically equivalent.

$$
\begin{gathered}
\sqrt{n} \hat{\gamma}_{j}=\sqrt{n} \tilde{\gamma}_{j}-\frac{1}{n} \sum_{t=j+1}^{n}\left(\mathbf{x}_{t-j} \cdot \varepsilon_{t}+\mathbf{x}_{t} \cdot \varepsilon_{t-j}\right)^{\prime} \sqrt{n}(\mathbf{b}-\boldsymbol{\beta}) \\
+\sqrt{n}(\mathbf{b}-\boldsymbol{\beta})^{\prime}\left(\frac{1}{n} \sum_{t=j+1}^{n} \mathbf{x}_{t} \mathbf{x}_{t-j}^{\prime}\right)(\mathbf{b}-\boldsymbol{\beta}) \\
\sim \sqrt{n} \tilde{\gamma}_{j}-\frac{1}{n} \sum_{t=j+1}^{n}\left(\mathbf{x}_{t-j} \cdot \varepsilon_{t}+\mathbf{x}_{t} \cdot \varepsilon_{t-j}\right)^{\prime} \sqrt{n}(\mathbf{b}-\boldsymbol{\beta})
\end{gathered}
$$

(since the last term above vanishes)

$$
\begin{aligned}
& \sim \sqrt{n} \tilde{\gamma}_{j}-\mathrm{E}\left(\mathbf{x}_{t-j} \cdot \varepsilon_{t}+\mathbf{x}_{t} \cdot \varepsilon_{t-j}\right)^{\prime} \sqrt{n}(\mathbf{b}-\boldsymbol{\beta}) \\
& \quad\left(\text { since } \frac{1}{n} \sum_{t=j+1}^{n}\left(\mathbf{x}_{t-j} \cdot \varepsilon_{t}+\mathbf{x}_{t} \cdot \varepsilon_{t-j}\right) \rightarrow \mathrm{E}\left(\mathbf{x}_{t-j} \cdot \varepsilon_{t}+\mathbf{x}_{t} \cdot \varepsilon_{t-j}\right)\right) \\
&= \sqrt{n} \tilde{\gamma}_{j}-\boldsymbol{\mu}_{j}^{\prime} \sqrt{n}(\mathbf{b}-\boldsymbol{\beta}) \quad \text { where } \boldsymbol{\mu}_{j} \equiv \mathrm{E}\left(\mathbf{x}_{t} \cdot \varepsilon_{t-j}\right) \\
& \quad\left(\text { since } \mathrm{E}\left(\mathbf{x}_{t-j} \cdot \varepsilon_{t}\right)=\mathbf{0} \text { by }(2.10 .15)\right) \\
&= \sqrt{n} \tilde{\gamma}_{j}-\boldsymbol{\mu}_{j}^{\prime} \sqrt{n} \mathbf{S}_{\mathbf{x x}}^{-1} \frac{1}{n} \sum_{t=1}^{n} \mathbf{x}_{t} \cdot \varepsilon_{t} \\
&= \sqrt{n} \frac{1}{n} \sum_{t=j+1}^{n} \varepsilon_{t} \varepsilon_{t-j}-\boldsymbol{\mu}_{j}^{\prime} \sqrt{n} \mathbf{S}_{\mathbf{x x}}^{-1} \frac{1}{n} \sum_{t=1}^{n} \mathbf{x}_{t} \cdot \varepsilon_{t} \quad(\text { by }(2.10 .7)) \\
& \widetilde{a} \sqrt{n} \frac{1}{n} \sum_{t=1}^{n} \varepsilon_{t} \varepsilon_{t-j}-\boldsymbol{\mu}_{j}^{\prime} \sqrt{n} \mathbf{S}_{\mathbf{x x}}^{-1} \frac{1}{n} \sum_{t=1}^{n} \mathbf{x}_{t} \cdot \varepsilon_{t} \\
& \quad\left(\text { since } \sqrt{n} \frac{1}{n} \sum_{t=1}^{n} \varepsilon_{t} \varepsilon_{t-j}-\sqrt{n} \frac{1}{n} \sum_{t=j+1}^{n} \varepsilon_{t} \varepsilon_{t-j}\right. \\
&\left.=\frac{1}{\sqrt{n}} \sum_{t=1}^{j} \varepsilon_{t} \varepsilon_{t-j} \rightarrow 0 \text { for each } j\right) \\
&\left.\widetilde{\mathrm{p}} \sqrt{n} \frac{1}{n} \sum_{t=1}^{n} \varepsilon_{t} \varepsilon_{t-j}-\boldsymbol{\mu}_{j}^{\prime} \sqrt{n} \mathbf{\Sigma}_{\mathbf{x x}}^{-1} \frac{1}{n} \sum_{t=1}^{n} \mathbf{x}_{t} \cdot \varepsilon_{t} \quad \text { (since } \mathbf{S}_{\mathbf{x x}} \rightarrow \mathbf{\Sigma}_{\mathbf{x x}}\right) \\
&=\mathbf{c}_{j}^{\prime} \sqrt{n} \overline{\mathbf{g}}_{j},
\end{aligned}
$$

where

\[
\underset{((K+1) \times 1)}{\mathbf{c}_{j}} \equiv\left[\begin{array}{c}
1  \tag{2.B.1}\\
-\boldsymbol{\Sigma}_{\mathbf{x x}}^{-1} \boldsymbol{\mu}_{j}
\end{array}\right], \quad \underset{((K+1) \times 1)}{\overline{\mathbf{g}}_{j}} \equiv \frac{1}{n} \sum_{t=1}^{n} \mathbf{g}_{j t}, \quad \underset{((K+1) \times 1)}{\mathbf{g}_{j t}} \equiv\left[\begin{array}{l}
\varepsilon_{t-j} \varepsilon_{t} \\
\mathbf{x}_{t} \cdot \varepsilon_{t}
\end{array}\right] .
\]

Thus, letting

$$
\underset{(p \times 1)}{\hat{\gamma}} \equiv\left[\begin{array}{c}
\hat{\gamma}_{1} \\
\vdots \\
\hat{\gamma}_{p}
\end{array}\right]
$$

we have proved that

$$
\sqrt{n} \hat{\boldsymbol{\gamma}} \underset{a}{\sim} \mathbf{C}^{\prime} \sqrt{n} \overline{\mathbf{g}}
$$

where

$$
\begin{aligned}
& \underset{(p(K+1) \times p)}{\mathbf{C}} \equiv\left[\begin{array}{lll}
\mathbf{c}_{1} & & \\
& \ddots & \\
& & \mathbf{c}_{p}
\end{array}\right] \\
& \underset{(p(K+1) \times 1)}{\overline{\mathbf{g}}} \equiv \frac{1}{n} \sum_{t=1}^{n} \mathbf{g}_{t}, \\
& \underset{(p(K+1) \times 1)}{\mathbf{g}_{t}} \equiv\left[\begin{array}{c}
\mathbf{g}_{1 t} \\
\vdots \\
\mathbf{g}_{p t}
\end{array}\right]
\end{aligned}
$$

The second part of the proof is to show that $\mathbf{g}_{t}$ is a martingale difference sequence, namely, that

$$
\mathrm{E}\left(\mathbf{g}_{j t} \mid \mathbf{g}_{t-1}, \mathbf{g}_{t-2}, \ldots\right)=\mathbf{0} \quad(\text { for } j=1,2, \ldots, p)
$$

But this easily follows from the Law of Total Expectations, the fact that ( $\mathbf{g}_{t-1}$, $\left.\mathbf{g}_{t-2}, \ldots\right)$ has less information than $\left(\mathbf{x}_{t}, \mathbf{x}_{t-1}, \ldots, \varepsilon_{t-1}, \varepsilon_{t-2}, \ldots\right)$, and (2.10.15).

Clearly, $\mathbf{g}_{t}$ is ergodic stationary. Therefore, by the ergodic stationary Martingale Differences CLT, $\sqrt{n} \overline{\mathbf{g}} \rightarrow_{\mathrm{d}} N\left(\mathbf{0}, \mathrm{E}\left(\mathbf{g}_{t} \mathbf{g}_{t}^{\prime}\right)\right)$. The next step is to calculate $\mathrm{E}\left(\mathbf{g}_{t} \mathbf{g}_{t}^{\prime}\right)$. Its $(j, k)$ block is

$$
\underset{((K+1) \times(K+1))}{\mathrm{E}\left(\mathbf{g}_{j t} \mathbf{g}_{k t}^{\prime}\right)}=\left[\begin{array}{cc}
\mathrm{E}\left(\varepsilon_{t}^{2} \varepsilon_{t-j} \varepsilon_{t-k}\right) & \mathrm{E}\left(\mathbf{x}_{t}^{\prime} \varepsilon_{t}^{2} \varepsilon_{t-j}\right) \\
\mathrm{E}\left(\mathbf{x}_{t} \varepsilon_{t}^{2} \varepsilon_{t-k}\right) & \mathrm{E}\left(\varepsilon_{t}^{2} \mathbf{x}_{t} \mathbf{x}_{t}^{\prime}\right)
\end{array}\right]
$$

By using the Law of Total Expectations and (2.10.16), it is easy to show that

\[
\underset{((K+1) \times(K+1))}{\mathrm{E}\left(\mathbf{g}_{j t} \mathbf{g}_{k t}^{\prime}\right)}=\left[\begin{array}{cc}
\sigma^{4} \delta_{j k} & \sigma^{2} \boldsymbol{\mu}_{j}^{\prime}  \tag{2.B.2}\\
\sigma^{2} \boldsymbol{\mu}_{k} & \sigma^{2} \mathbf{\Sigma}_{\mathbf{x x}}
\end{array}\right]
\]

where $\delta_{j k}$ is "Kronecker's delta," which is 1 if $j=k$ and 0 otherwise.\\
Since, as shown above, $\sqrt{n} \hat{\boldsymbol{\gamma}} \underset{a}{\sim} \mathbf{C}^{\prime} \sqrt{n} \overline{\mathbf{g}}$ and $\sqrt{n} \overline{\mathbf{g}} \rightarrow_{\mathrm{d}} N\left(\mathbf{0}, \mathrm{E}\left(\mathbf{g}_{t} \mathbf{g}_{t}^{\prime}\right)\right)$, we have

$$
\operatorname{Avar}(\hat{\gamma})(\equiv \text { variance of limiting distribution of } \sqrt{n} \hat{\gamma})=\mathbf{C}^{\prime} \mathrm{E}\left(\mathbf{g}_{t} \mathbf{g}_{t}^{\prime}\right) \mathbf{C} .
$$

Its $(j, k)$ element is

$$
(j, k) \text { element of } \operatorname{Avar}(\hat{\boldsymbol{\gamma}})=\mathbf{c}_{j}^{\prime} \mathrm{E}\left(\mathbf{g}_{j t} \mathbf{g}_{k t}^{\prime}\right) \mathbf{c}_{k}
$$

Substituting (2.B.2) into this and using the definition of $\mathbf{c}_{j}$ in (2.B.1), we obtain, after some simple calculations,

$$
(j, k) \text { element of } \operatorname{Avar}(\hat{\boldsymbol{\gamma}})=\sigma^{4} \cdot\left[\delta_{j k}-\boldsymbol{\mu}_{j}^{\prime} \boldsymbol{\Sigma}_{\mathbf{x x}}^{-1} \boldsymbol{\mu}_{k} / \sigma^{2}\right] .
$$

So

$$
\sqrt{n} \hat{\boldsymbol{\gamma}} \underset{\mathrm{~d}}{\rightarrow} N\left(\mathbf{0}, \sigma^{4} \mathbf{I}_{p}-\sigma^{4} \mathbf{\Phi}\right),
$$

where $\boldsymbol{\Phi}$ is defined in Proposition 2.10. Since $\sqrt{n} \hat{\boldsymbol{\rho}} \underset{a}{\sim} \sqrt{n} \hat{\boldsymbol{\gamma}} / \sigma^{2}$, the limiting distribution of $\sqrt{n} \hat{\boldsymbol{\rho}}$ is the same as that of $\sqrt{n} \hat{\boldsymbol{\gamma}} / \sigma^{2}$, which is $N\left(\mathbf{0}, \mathbf{I}_{p}-\boldsymbol{\Phi}\right)$.

\section*{PROBLEM SET FOR CHAPTER 2}
\section*{ANALYTICAL EXERCISES}
\begin{enumerate}
  \item (Taken from Example 3.4.2 of Amemiya (1985)) Let $z_{n}$ be defined by
\end{enumerate}

$$
z_{n}= \begin{cases}0 & \text { with probability }(n-1) / n \\ n^{2} & \text { with probability } 1 / n\end{cases}
$$

Show that $\operatorname{plim}_{n \rightarrow \infty} z_{n}=0$ but $\lim _{n \rightarrow \infty} \mathrm{E}\left(z_{n}\right)=\infty$.\\
2. Prove Chebychev's weak LLN, by showing that $\bar{z}_{n} \rightarrow_{\text {m.s. }} \mu$. Hint:

$$
\bar{z}_{n}-\mu=\left(\bar{z}_{n}-\mathrm{E}\left(\bar{z}_{n}\right)\right)+\left(\mathrm{E}\left(\bar{z}_{n}\right)-\mu\right) .
$$

So

$$
\left(\bar{z}_{n}-\mu\right)^{2}=\left(\bar{z}_{n}-\mathrm{E}\left(\bar{z}_{n}\right)\right)^{2}+2\left(\bar{z}_{n}-\mathrm{E}\left(\bar{z}_{n}\right)\right)\left(\mathrm{E}\left(\bar{z}_{n}\right)-\mu\right)+\left(\mathrm{E}\left(\bar{z}_{n}\right)-\mu\right)^{2} .
$$

Show that $\mathrm{E}\left[\left(\bar{z}_{n}-\mathrm{E}\left(\bar{z}_{n}\right)\right)\left(\mathrm{E}\left(\bar{z}_{n}\right)-\mu\right)\right]=0$.\\
3. (Consistency and Asymptotic Normality of OLS for Random Samples) Consider replacing Assumption 2.2 by

Assumption 2.2': $\left\{y_{i}, \mathbf{x}_{i}\right\}$ is a random sample.\\
and Assumption 2.5 by

Assumption 2.5': $\mathbf{S} \equiv \mathrm{E}\left(\mathbf{g}_{i} \mathbf{g}_{i}^{\prime}\right)$ exists and is finite.

Prove the following simplified version of Proposition 2.1:

Proposition 2.1 for Random Samples: Under Assumptions 2.1, 2.3, 2.4, and 2.2', the OLS estimator $\mathbf{b}$ is consistent. Under the additional assumption of Assumption 2.5', the OLS estimator $\mathbf{b}$ is asymptotically normal with $\operatorname{Avar}(\mathbf{b})$ given by (2.3.4).\\
(a) Show that this is a special case of Proposition 2.1 of the text.\\
(b) Give a direct proof of this proposition.

Hint: Use Kolmogorov and Linderberg-Levy.\\
4. (Proof of Proposition 2.4) We wish to prove Proposition 2.4. To avoid inessential complications, we assume as in the text that $K=1$ (only one regressor). What remains to be shown is that the sample mean of the first term on the RHS of (2.5.3) converges in probability to some finite number. Prove it. Hint: Let $f \equiv x_{i} \varepsilon_{i}$ and $h \equiv x_{i}^{2}$. The Cauchy-Schwartz inequality states:

$$
\mathrm{E}(|f \cdot h|) \leq \sqrt{\mathrm{E}\left(f^{2}\right) \mathrm{E}\left(h^{2}\right)} .
$$

So

$$
\mathrm{E}\left(\left|x_{i}^{3} \varepsilon_{i}\right|\right) \leq \sqrt{\mathrm{E}\left(x_{i}^{2} \varepsilon_{i}^{2}\right) \mathrm{E}\left(x_{i}^{4}\right)}
$$

By assumption, $\mathrm{E}\left(x_{i}^{2} \varepsilon_{i}^{2}\right)\left(=\mathrm{E}\left(g_{i}^{2}\right)\right)$ is finite, and by Assumption 2.6 $\mathrm{E}\left(x_{i}^{4}\right)$ is also finite. So $\mathrm{E}\left(\varepsilon_{i} x_{i}^{3}\right)$ is finite and $\frac{1}{n} \sum_{i} \varepsilon_{i} x_{i}^{3} \rightarrow_{\mathrm{p}}$ some finite number by ergodic stationarity.\\
5. (Direct proof that the change in $\operatorname{SSR}$ divided by $\sigma^{2}$ is asymptotically $\chi^{2}$ ) In Section 2.6, we proved that

$$
\left(S S R_{R}-S S R_{U}\right) / s^{2} \underset{\mathrm{~d}}{\rightarrow} \chi^{2}(\# \mathbf{r})
$$

as a corollary to Proposition 2.3. Give a direct proof, first by showing that

$$
\begin{gathered}
S S R_{R}-S S R_{U}=(\sqrt{n} \overline{\mathbf{g}})^{\prime} \mathbf{S}_{\mathbf{x x}}^{-1} \mathbf{R}^{\prime}\left(\mathbf{R} \mathbf{S}_{\mathbf{x x}}^{-1} \mathbf{R}^{\prime}\right)^{-1} \mathbf{R} \mathbf{S}_{\mathbf{x x}}^{-1}(\sqrt{n} \overline{\mathbf{g}}), \\
\overline{\mathbf{g}} \equiv \frac{1}{n} \sum_{i=1}^{n} \mathbf{x}_{i} \cdot \varepsilon_{i}
\end{gathered}
$$

\begin{enumerate}
  \setcounter{enumi}{5}
  \item (Optional, consistency of $\hat{\boldsymbol{\alpha}}$ ) In this question we prove the claim made in Section 2.8 that $\hat{\boldsymbol{\alpha}}$, the OLS estimate for the regression of $e_{i}^{2}$ on $\mathbf{z}_{i}$, is consistent.
\end{enumerate}

Let $\tilde{\alpha}$ be the OLS estimate from


\begin{equation*}
\varepsilon_{i}^{2}=\mathbf{z}_{i}^{\prime} \boldsymbol{\alpha}+\eta_{i}, \eta_{i} \equiv \varepsilon_{i}^{2}-\mathrm{E}\left(\varepsilon_{i}^{2} \mid \mathbf{x}_{i}\right) \tag{2.8.8}
\end{equation*}


(a) Make appropriate assumptions in addition to Assumptions 2.1 and 2.2 to show that $\tilde{\alpha}$ is consistent. Hint: You need a rank condition for $\mathbf{z}_{i}$.\\
(b) Derive:


\begin{equation*}
\hat{\alpha}-\tilde{\alpha}=\left(\frac{1}{n} \sum_{i=1}^{n} \mathbf{z}_{i} \mathbf{z}_{i}^{\prime}\right)^{-1} \frac{1}{n} \sum_{i=1}^{n} \mathbf{z}_{i} \cdot v_{i} \tag{*}
\end{equation*}


with $v_{i} \equiv-2(\mathbf{b}-\boldsymbol{\beta})^{\prime} \mathbf{x}_{i} \cdot \varepsilon_{i}+(\mathbf{b}-\boldsymbol{\beta})^{\prime} \mathbf{x}_{i} \mathbf{x}_{i}^{\prime}(\mathbf{b}-\boldsymbol{\beta})$. Hint: The discrepancy between $\varepsilon_{i}^{2}$ and the squared OLS residual $e_{i}^{2}$ from the OLS estimation of the original equation $y_{i}=\mathbf{x}_{i}^{\prime} \boldsymbol{\beta}+\varepsilon_{i}$ is given by (2.3.10). Substituting it into (2.8.8) gives: $e_{i}^{2}=\mathbf{z}_{i}^{\prime} \boldsymbol{\alpha}+\left(\eta_{i}+v_{i}\right)$.\\
(c) To avoid inessential complications, suppose that $\mathbf{x}_{i}$ is a scalar $x_{i}$. Then $\frac{1}{n} \sum_{i} \mathbf{z}_{i} \cdot v_{i}$ in ( $*$ ) becomes


\begin{equation*}
-2(b-\beta) \frac{1}{n} \sum_{i=1}^{n} x_{i} \varepsilon_{i} \mathbf{z}_{i}+(b-\beta)^{2} \frac{1}{n} \sum_{i=1}^{n} x_{i}^{2} \mathbf{z}_{i} \tag{**}
\end{equation*}


Show that the plim of the first term is zero. Hint:

$$
\mathrm{E}\left(x_{i} \varepsilon_{i} \mathbf{z}_{i}\right)=\mathrm{E}\left[\mathbf{z}_{i} \cdot x_{i} \cdot \mathrm{E}\left(\varepsilon_{i} \mid x_{i}\right)\right]
$$

Show that the plim of the second term is zero if $\mathrm{E}\left(x_{i}^{2} \mathbf{z}_{i}\right)$ exists and is finite.\\
(d) Thus we have proved that $\hat{\boldsymbol{\alpha}}-\tilde{\boldsymbol{\alpha}}$ vanishes. Prove the stronger result that $\sqrt{n}(\hat{\boldsymbol{\alpha}}-\tilde{\boldsymbol{\alpha}})$ vanishes. Hint: Use Lemma 2.4(b).\\
7. (Optional, proof that WLS is asymptotically more efficient than OLS) In Section 2.8, the WLS estimator is denoted $\widehat{\boldsymbol{\beta}}(\widehat{\mathbf{V}})$. We wish to prove that it is asymptotically more efficient than OLS, namely,

$$
\operatorname{Avar}(\widehat{\boldsymbol{\beta}}(\widehat{\mathbf{V}})) \leq \operatorname{Avar}(\mathbf{b})
$$

But since $\operatorname{Avar}(\widehat{\boldsymbol{\beta}}(\widehat{\mathbf{V}}))=\operatorname{Avar}(\widehat{\boldsymbol{\beta}}(\mathbf{V}))$, it is sufficient to prove that

$$
\operatorname{Avar}(\widehat{\boldsymbol{\beta}}(\mathbf{V})) \leq \operatorname{Avar}(\mathbf{b})
$$

Prove this last matrix inequality. Hint: In Chapter 1, we proved the algebraic\\
result that


\begin{equation*}
\left(\mathbf{X}^{\prime} \mathbf{X}\right)^{-1}\left(\mathbf{X}^{\prime} \mathbf{V} \mathbf{X}\right)\left(\mathbf{X}^{\prime} \mathbf{X}\right)^{-1} \geq\left(\mathbf{X}^{\prime} \mathbf{V}^{-1} \mathbf{X}\right)^{-1} \tag{*}
\end{equation*}


for any positive definite matrix $\mathbf{V}$. (This follows from the fact that GLS is more efficient than OLS; the LHS is ( $\sigma^{2}$ times) the variance of the OLS estimator for the generalized regression model while the RHS is ( $\sigma^{2}$ times) the variance of the GLS estimator.) Dividing both sides of $(*)$ by $n$ and taking probability limits yields:


\begin{equation*}
\left(\operatorname{plim} \frac{1}{n} \mathbf{X}^{\prime} \mathbf{X}\right)^{-1}\left(\operatorname{plim} \frac{1}{n} \mathbf{X}^{\prime} \mathbf{V} \mathbf{X}\right)\left(\operatorname{plim} \frac{1}{n} \mathbf{X}^{\prime} \mathbf{X}\right)^{-1} \geq \operatorname{plim}\left(\frac{1}{n} \mathbf{X}^{\prime} \mathbf{V}^{-1} \mathbf{X}\right)^{-1} \tag{**}
\end{equation*}


The RHS is $\operatorname{Avar}(\widehat{\boldsymbol{\beta}}(\mathbf{V}))$ (see (2.8.7)). Show that the LHS equals Avar(b) $= \boldsymbol{\Sigma}_{\mathbf{x x}}^{-1} \mathbf{S} \boldsymbol{\Sigma}_{\mathbf{x x}}^{-1}$.\\
8. Prove (2.9.7). Hint: If $(\mu, \boldsymbol{\gamma})$ is the least squares projection coefficients, it satisfies

\[
\mathrm{E}\left(\mathbf{x x}^{\prime}\right)\left[\begin{array}{l}
\mu  \tag{*}\\
\boldsymbol{\gamma}
\end{array}\right]=\mathrm{E}(\mathbf{x} \cdot y)
\]

or

\[
\left[\begin{array}{l}
\mu  \tag{**}\\
\boldsymbol{\gamma}
\end{array}\right]=\left[\mathrm{E}\left(\mathbf{x} \mathbf{x}^{\prime}\right)\right]^{-1} \mathrm{E}(\mathbf{x} \cdot \boldsymbol{y}),
\]

where

$$
\mathrm{E}\left(\mathbf{x}^{\prime}\right)=\left[\begin{array}{cc}
1 & \mathrm{E}(\tilde{\mathbf{x}})^{\prime} \\
\mathrm{E}(\tilde{\mathbf{x}}) & \mathrm{E}\left(\tilde{\mathbf{x}} \tilde{\mathbf{x}}^{\prime}\right)
\end{array}\right], \mathrm{E}(\mathbf{x} \cdot y)=\left[\begin{array}{c}
\mathrm{E}(y) \\
\mathrm{E}(\tilde{\mathbf{x}} \cdot y)
\end{array}\right] .
$$

The first equation of $(*)$ is $\mu+\mathrm{E}(\tilde{\mathbf{x}})^{\prime} \boldsymbol{\gamma}=\mathrm{E}(y)$. Use this to eliminate $\mu$ from the rest of the equations of (*) and then solve for $\boldsymbol{\gamma}$. In the process, use $\operatorname{Var}(\tilde{\mathbf{x}})= \mathrm{E}\left(\tilde{\mathbf{x}} \tilde{\mathbf{x}}^{\prime}\right)-\mathrm{E}(\tilde{\mathbf{x}}) \mathrm{E}\left(\tilde{\mathbf{x}}^{\prime}\right), \operatorname{Cov}(\tilde{\mathbf{x}}, y)=\mathrm{E}(\tilde{\mathbf{x}} \cdot y)-\mathrm{E}(\tilde{\mathbf{x}}) \mathrm{E}(y)$. Alternatively, use the formula for the partitioned inverse (see (A.10) of Appendix A) to obtain

$$
\mathrm{E}\left(\mathbf{x} \mathbf{x}^{\prime}\right)^{-1}=\left[\begin{array}{cc}
1+\mathrm{E}(\tilde{\mathbf{x}})^{\prime} \operatorname{Var}(\tilde{\mathbf{x}})^{-1} \mathrm{E}(\tilde{\mathbf{x}}) & -\mathrm{E}(\tilde{\mathbf{x}})^{\prime} \operatorname{Var}(\tilde{\mathbf{x}})^{-1} \\
-\operatorname{Var}(\tilde{\mathbf{x}})^{-1} \mathrm{E}(\tilde{\mathbf{x}}) & \operatorname{Var}(\tilde{\mathbf{x}})^{-1}
\end{array}\right] .
$$

Then substitute this into (**).\\
9. (Proof of Proposition 2.9 for $p=1$ ) Assume that $\left\{\varepsilon_{t}\right\}$ is an ergodic stationary martingale difference sequence and that $\mathrm{E}\left(\varepsilon_{t}^{2} \mid \varepsilon_{t-1}, \varepsilon_{t-2}, \ldots, \varepsilon_{1}\right)=\sigma^{2}<\infty$.

Let $g_{t}=\varepsilon_{t} \cdot \varepsilon_{t-1}$ and

$$
\hat{\gamma}_{j}=\frac{1}{n} \sum_{t=j+1}^{n} \varepsilon_{t} \cdot \varepsilon_{t-j} \quad(j=0,1), \hat{\rho}_{1}=\frac{\hat{\gamma}_{1}}{\hat{\gamma}_{0}} .
$$

(a) Show that $\left\{g_{t}\right\}(t=2,3, \ldots)$ is a martingale difference sequence.\\
(b) Show that $\mathrm{E}\left(g_{t}^{2}\right)=\sigma^{4}$.\\
(c) Show that $\sqrt{n} \hat{\gamma}_{1} \rightarrow_{\mathrm{d}} N\left(0, \sigma^{4}\right)$ as $n \rightarrow \infty$.\\
(d) Show that $\sqrt{n} \hat{\rho}_{1} \rightarrow_{\mathrm{d}} N(0,1)$. Hint: First show that $\sqrt{n} \hat{\gamma}_{1} / \sigma^{2} \rightarrow_{\mathrm{d}} N(0,1)$. Then use Lemma 2.4(c).\\
10. (Asymptotics of sample mean of MA(2)) Let ( $\varepsilon_{-1}, \varepsilon_{0}, \varepsilon_{1}, \varepsilon_{2}, \ldots$ ) be i.i.d. with mean zero and variance $\sigma_{\varepsilon}^{2}$. Consider a process $\left(y_{-1}, y_{0}, y_{1}, y_{2}, \ldots\right)$ generated by

$$
y_{-1}=\varepsilon_{-1}, y_{0}=\varepsilon_{0}+\theta_{1} \varepsilon_{-1}, y_{t}=\varepsilon_{t}+\theta_{1} \varepsilon_{t-1}+\theta_{2} \varepsilon_{t-2} \quad(t=1,2, \ldots) .
$$

(This process is called a second-order moving average process (MA(2)); see Chapter 6.\\
(a) Show that $\left(y_{1}, y_{2}, \ldots\right)$ is covariance stationary. Derive the expressions for the autocovariances $\gamma_{j}(j=0,1,2, \ldots)$.\\
(b) Let

$$
\begin{aligned}
& r_{t j} \equiv \mathrm{E}\left(y_{t} \mid y_{t-j}, y_{t-j-1}, \ldots, y_{0}, y_{-1}\right) \\
& -\mathrm{E}\left(y_{t} \mid y_{t-j-1}, y_{t-j-2}, \ldots, y_{0}, y_{-1}\right) \\
& \quad(t=j, j+1, \ldots ; j=0,1,2, \ldots)
\end{aligned}
$$

Show that

$$
r_{t 0}=\varepsilon_{t}, r_{t 1}=\theta_{1} \varepsilon_{t-1}, r_{t 2}=\theta_{2} \varepsilon_{t-2}, r_{t 3}=0, r_{t 4}=0, \ldots
$$

Hint: There is a one-to-one mapping between ( $\varepsilon_{-1}, \varepsilon_{0}, \varepsilon_{1}, \ldots, \varepsilon_{t}$ ) and $\left(y_{-1}, y_{0}, y_{1}, \ldots, y_{t}\right)$ so $\mathrm{E}\left(y_{t} \mid y_{t-j}, \ldots, y_{-1}\right)=\mathrm{E}\left(y_{t} \mid \varepsilon_{t-j}, \ldots, \varepsilon_{-1}\right)$.\\
(c) Let

$$
\bar{y}_{n} \equiv \frac{1}{n}\left(y_{1}+y_{2}+\cdots+y_{n}\right) .
$$

Show that

$$
\operatorname{Var}\left(\sqrt{n} \bar{y}_{n}\right)=\gamma_{0}+2\left[\left(1-\frac{1}{n}\right) \gamma_{1}+\left(1-\frac{2}{n}\right) \gamma_{2}\right] .
$$

(d) Because $\left\{y_{t}\right\}$ is not i.i.d, the Lindeberg-Levy CLT is not applicable. However, it will be shown in Chapter 6 that $\sqrt{n} \bar{y}_{n}$ converges in distribution to a normal random variable. What is the mean of the limiting distribution? What is the variance of the limiting distribution (i.e., $\operatorname{Avar}\left(\bar{y}_{n}\right)$ )? Hint: In Lemma 2.1, set $z_{n}=\sqrt{n} \bar{y}_{n}$, and $\lim _{n \rightarrow \infty}(1-q / n)=1$ for any fixed $q$.\\
11. (Optional, Breusch-Godfrey test for serial correlation) In this question, we prove that the modified Box-Pierce $Q$ is asymptotically equivalent to $p F$ from auxiliary regression (2.10.21), where $p$ is the number of lags and $F$ is the $F$ ratio for the hypothesis that the coefficients of $e_{t-1}, e_{t-2}, \ldots, e_{t-p}$ are all zero.

First we establish the notation. Let

$$
\begin{gathered}
\underset{(n \times K)}{\mathbf{X}}=\left[\begin{array}{c}
\mathbf{x}_{1}^{\prime} \\
\vdots \\
\mathbf{x}_{n}^{\prime}
\end{array}\right], \underset{(n \times p)}{\mathbf{E}}=\left[\begin{array}{cccc}
0 & 0 & \cdots & 0 \\
e_{1} & 0 & \cdots & \vdots \\
e_{2} & e_{1} & & 0 \\
e_{3} & e_{2} & & e_{1} \\
\vdots & \vdots & & \vdots \\
e_{n-1} & e_{n-2} & \cdots & e_{n-p}
\end{array}\right], \underset{(n \times 1)}{\mathbf{e}}=\left[\begin{array}{c}
e_{1} \\
\vdots \\
e_{n}
\end{array}\right], \\
\underset{((K+p) \times(K+p))}{\widehat{\mathbf{B}}}=\left[\begin{array}{cc}
\frac{1}{n} \mathbf{X}^{\prime} \mathbf{X} & \frac{1}{n} \mathbf{X}^{\prime} \mathbf{E} \\
\frac{1}{n} \mathbf{E}^{\prime} \mathbf{X} & \frac{1}{n} \mathbf{E}^{\prime} \mathbf{E}
\end{array}\right], \widehat{\mathbf{B}}^{-1}=\left[\begin{array}{cc}
\widehat{\mathbf{B}}^{11} & \widehat{\mathbf{B}}^{12} \\
(K \times K) & (K \times p) \\
\widehat{\mathbf{B}}^{21} & \underset{(p \times K)}{\widehat{\mathbf{B}^{22}}}
\end{array}\right], \underset{(p \times p)}{\hat{\gamma}}=\left[\begin{array}{c}
\hat{\gamma}_{1} \\
\vdots \\
\hat{\gamma}_{p}
\end{array}\right],
\end{gathered}
$$

where $e_{t}$ is the residual from the original regression $y_{t}=\mathbf{x}_{t}^{\prime} \boldsymbol{\beta}+\varepsilon_{t}$ and $\hat{\gamma}_{j}$ is the sample $j$-th order autocovariance, defined in (2.10.8), calculated from the residuals.\\
(a) The auxiliary regression (2.10.21) has $K+p$ regressors. Let $\hat{\boldsymbol{\alpha}}$ be the vector of the $K+p$ coefficients. Show that

$$
\hat{\boldsymbol{\alpha}}=\widehat{\mathbf{B}}^{-1}\left[\begin{array}{c}
\underset{(K \times 1)}{\mathbf{0}} \\
\underset{(p \times 1)}{\hat{\gamma}}
\end{array}\right] .
$$

Hint: Since $\mathbf{e}$ is the vector of residuals from the original regression with $\mathbf{X}$ as regressors, $\mathbf{X}^{\prime} \mathbf{e}=\mathbf{0}$.\\
(b) Show:

$$
\widehat{\mathbf{B}} \underset{\mathrm{p}}{\rightarrow} \mathbf{B} \equiv\left[\begin{array}{cc}
\boldsymbol{\Sigma}_{\mathbf{x x}} & \mathbf{H} \\
\mathbf{H}^{\prime} & \sigma^{2} \mathbf{I}_{p}
\end{array}\right],
$$

where

$$
\underset{(K \times p)}{\mathbf{H}} \equiv\left[\begin{array}{lll}
\mathrm{E}\left(\mathbf{x}_{t} \cdot \varepsilon_{t-1}\right) & \cdots & \mathrm{E}\left(\mathbf{x}_{t} \cdot \varepsilon_{t-p}\right)
\end{array}\right]
$$

Hint: The $j$-th column of $\frac{1}{n} \mathbf{X}^{\prime} \mathbf{E}$ is $\frac{1}{n} \sum_{t=j+1}^{n} \mathbf{x}_{t} \cdot e_{t-j}$. Use (2.3.9) to show that it converges in probability to $\mathrm{E}\left(\mathbf{x}_{t} \cdot \varepsilon_{t-j}\right)$.\\
(c) (very easy) Show: $\hat{\boldsymbol{\alpha}} \rightarrow_{\mathrm{p}} \mathbf{0}$.\\
(d) Let $S S R$ here be the sum of squared residuals from the auxiliary, not the original, regression. Show:

$$
\frac{S S R}{n-K-p} \underset{\mathrm{p}}{\rightarrow} \sigma^{2}
$$

where $\sigma^{2}$ is the variance of $\varepsilon_{t}$, the error term from the original regression.\\
Hint: $[\mathbf{X} \vdots \mathbf{E}] \hat{\boldsymbol{\alpha}}=\mathbf{E} \hat{\boldsymbol{\gamma}}$. So $S S R=(\mathbf{e}-\mathbf{E} \hat{\boldsymbol{\gamma}})^{\prime}(\mathbf{e}-\mathbf{E} \hat{\boldsymbol{\gamma}}) . \frac{1}{n} \mathbf{E}^{\prime} \mathbf{e}=\hat{\boldsymbol{\gamma}}$, $\frac{1}{n} \mathbf{E}^{\prime} \mathbf{E} \rightarrow_{\mathrm{p}} \sigma^{2} \mathbf{I}_{p}$.\\
(e) Show:

$$
p F=\frac{n \hat{\boldsymbol{\gamma}}^{\prime} \widehat{\mathbf{B}}^{22} \hat{\boldsymbol{\gamma}}}{S S R /(n-K-p)} .
$$

Hint: Apply the formula for the $F$-ratio in (1.4.9) to the auxiliary regression.\\
(f) Use the formula for the partitioned inverse (see (A.10) of Appendix A) to show

$$
\widehat{\mathbf{B}}^{22}=\left[\frac{1}{n} \mathbf{E}^{\prime} \mathbf{E}-\left(\frac{1}{n} \mathbf{E}^{\prime} \mathbf{X}\right) \mathbf{S}_{\mathbf{x x}}^{-1}\left(\frac{1}{n} \mathbf{X}^{\prime} \mathbf{E}\right)\right]^{-1} .
$$

(g) Let $\hat{\rho}$ and $\widehat{\Phi}$ be as in (2.10.8) and (2.10.19). Show that the modified BoxPierce $Q$ defined in (2.10.20) is asymptotically equivalent to $p F$ (i.e., the difference between the two converges to zero in probability). Hint: Show that both $p F$ and the modified $Q$ are asymptotically equivalent to

$$
n \hat{\boldsymbol{\gamma}}^{\prime}\left(\mathbf{I}_{p}-\boldsymbol{\Phi}\right)^{-1} \hat{\boldsymbol{\gamma}} / \sigma^{2}
$$

where $\boldsymbol{\Phi}$ is defined in (2.10.18).\\
12. (Optional, Chow test for structural change in large samples) Consider applying the Chow test for structural change to the regression model with conditional homoskedasticity described in Proposition 2.5. Since the issue of structural change arises mostly in time-series models, we use " $t$ " for the subscript. We generalize Assumption 2.1 by allowing the coefficient vector to change at the break date $r$ :

$$
y_{t}= \begin{cases}\mathbf{x}_{t}^{\prime} \boldsymbol{\beta}_{1}+\varepsilon_{t} & \text { for } t=1,2, \ldots, r \\ \mathbf{x}_{t}^{\prime} \boldsymbol{\beta}_{2}+\varepsilon_{t} & \text { for } t=r+1, r+2, \ldots, n\end{cases}
$$

We assume that the break date $r$ is known. (When the break date is unknown, the situation is more complicated and has been an object of recent research. See Stock (1994, Section 5) for a survey.) The null hypothesis to test is that $\boldsymbol{\beta}_{1}=\boldsymbol{\beta}_{2}$. This is a set of $K$ restrictions. Let $S S R_{1}$ be the sum of squared residuals from the first period $(t=1,2, \ldots, r), S S R_{2}$ be the $S S R$ from the second period ( $t=r+1, r+2, \ldots, n$ ), and $S S R_{R}$ be the $S S R$ from the entire sample under the constraint $\boldsymbol{\beta}_{1}=\boldsymbol{\beta}_{2}$. From Chapter 1, the Chow statistic is defined as

$$
F=\frac{\left[S S R_{R}-\left(S S R_{1}+S S R_{2}\right)\right] / K}{\left(S S R_{1}+S S R_{2}\right) /(n-2 K)}
$$

Let $\mathbf{b}_{1}$ be the OLS estimate of $\boldsymbol{\beta}_{1}$ obtained from the first period and $\mathbf{b}_{2}$ be the OLS estimate of $\boldsymbol{\beta}_{2}$ from the second period.\\
(a) Show that $K F$ is numerically equal to

$$
\sqrt{n}\left(\mathbf{b}_{1}-\mathbf{b}_{2}\right)^{\prime}\left[\left(\frac{\frac{1}{n} \sum_{t=1}^{r} \mathbf{x}_{t} \mathbf{x}_{t}^{\prime}}{s^{2}}\right)^{-1}+\left(\frac{\frac{1}{n} \sum_{t=r+1}^{n} \mathbf{x}_{t} \mathbf{x}_{t}^{\prime}}{s^{2}}\right)^{-1}\right]^{-1} \sqrt{n}\left(\mathbf{b}_{1}-\mathbf{b}_{2}\right),
$$

where $s^{2}=\left(S S R_{1}+S S R_{2}\right) /(n-2 K)$. Hint: The $n$ equations can be written in matrix form as $\mathbf{y}=\mathbf{X} \boldsymbol{\beta}+\boldsymbol{\varepsilon}$, where

$$
\begin{gathered}
\underset{(n \times 2 K)}{\mathbf{X}}=\left[\begin{array}{cc}
\mathbf{X}_{1} & \mathbf{0} \\
\mathbf{0} & \mathbf{X}_{2}
\end{array}\right], \underset{(r \times K)}{\mathbf{X}_{1}}=\left[\begin{array}{c}
\mathbf{x}_{1}^{\prime} \\
\vdots \\
\mathbf{x}_{r}^{\prime}
\end{array}\right], \\
\underset{((n-r) \times K)}{\mathbf{X}_{2}}=\left[\begin{array}{c}
\mathbf{x}_{r+1}^{\prime} \\
\vdots \\
\mathbf{x}_{n}^{\prime}
\end{array}\right], \quad \underset{(2 K \times 1)}{\boldsymbol{\beta}}=\left[\begin{array}{c}
\boldsymbol{\beta}_{1} \\
\boldsymbol{\beta}_{2}
\end{array}\right] .
\end{gathered}
$$

The null hypothesis that $\boldsymbol{\beta}_{1}=\boldsymbol{\beta}_{2}$ can be written as $\mathbf{R} \boldsymbol{\beta}=\mathbf{r}$ where $\mathbf{R}= \left[\mathbf{I}_{K} \vdots-\mathbf{I}_{K}\right]$ and $\mathbf{r}=\mathbf{0}$. Use Proposition 1.4.\\
(b) Let $\lambda \equiv r / n$. Show that, as $n \rightarrow \infty$ with $\lambda$ fixed,

$$
\frac{1}{n} \sum_{t=1}^{r} \mathbf{x}_{t} \mathbf{x}_{t}^{\prime} \underset{\mathbf{p}}{\rightarrow} \lambda \boldsymbol{\Sigma}_{\mathbf{x x}}, \frac{1}{n} \sum_{t=r+1}^{n} \mathbf{x}_{t} \mathbf{x}_{t}^{\prime} \underset{\mathbf{p}}{\rightarrow}(1-\lambda) \boldsymbol{\Sigma}_{\mathbf{x x}}
$$

where $\boldsymbol{\Sigma}_{\mathbf{x x}} \equiv \mathrm{E}\left(\mathbf{x}_{t} \mathbf{x}_{t}^{\prime}\right)$.\\
(c) Show that, as $n \rightarrow \infty$ with $\lambda$ fixed,

$$
\frac{1}{\sqrt{n}} \sum_{t=1}^{r} \mathbf{x}_{t} \cdot \varepsilon_{t} \rightarrow N\left(\mathbf{0}, \lambda \sigma^{2} \boldsymbol{\Sigma}_{\mathbf{x x}}\right), \frac{1}{\sqrt{n}} \sum_{t=r+1}^{n} \mathbf{x}_{t} \cdot \varepsilon_{t} \rightarrow N\left(\mathbf{0},(1-\lambda) \sigma^{2} \boldsymbol{\Sigma}_{\mathbf{x x}}\right) .
$$

It can be shown (see, e.g., Stock, 1994, Section 5) that $\frac{1}{\sqrt{n}} \sum_{t=1}^{r} \mathbf{x}_{t} \varepsilon_{t}$ and $\frac{1}{\sqrt{n}} \sum_{t=r+1}^{n} \mathbf{x}_{t} \varepsilon_{t}$ are asymptotically uncorrelated. So

$$
\left[\begin{array}{c}
\frac{1}{\sqrt{n}} \sum_{t=1}^{r} \mathbf{x}_{t} \cdot \varepsilon_{t} \\
\frac{1}{\sqrt{n}} \sum_{t=r+1}^{n} \mathbf{x}_{t} \cdot \varepsilon_{t}
\end{array}\right] \rightarrow N\left(\mathbf{0},\left[\begin{array}{cc}
\lambda \sigma^{2} \boldsymbol{\Sigma}_{\mathbf{x} \mathbf{x}} & \mathbf{0} \\
\mathbf{0} & (1-\lambda) \sigma^{2} \mathbf{\Sigma}_{\mathbf{x x}}
\end{array}\right]\right)
$$

(d) Show that $\operatorname{Avar}\left(\mathbf{b}_{1}-\mathbf{b}_{2}\right)=\lambda^{-1} \sigma^{2} \boldsymbol{\Sigma}_{\mathbf{x x}}^{-1}+(1-\lambda)^{-1} \sigma^{2} \boldsymbol{\Sigma}_{\mathbf{x x}}^{-1}$. Hint: Let

$$
\mathbf{b} \equiv\left[\begin{array}{l}
\mathbf{b}_{1} \\
\mathbf{b}_{2}
\end{array}\right]
$$

Show that

$$
\operatorname{Avar}(\mathbf{b})=\left[\begin{array}{cc}
\lambda^{-1} \sigma^{2} \boldsymbol{\Sigma}_{\mathbf{x x}}^{-1} & \mathbf{0} \\
\mathbf{0} & (1-\lambda)^{-1} \sigma^{2} \boldsymbol{\Sigma}_{\mathbf{x x}}^{-1}
\end{array}\right]
$$

$\mathbf{b}_{1}-\mathbf{b}_{2}=\left[\mathbf{I}_{K} \vdots-\mathbf{I}_{K}\right] \mathbf{b}$. Use Lemma 2.4(c).\\
(e) Finally, show that $K F \rightarrow_{\mathrm{d}} \chi^{2}(K)$ as $n \rightarrow \infty$ with $\lambda=r / n$ fixed.

\section*{EMPIRICAL EXERCISES}
\begin{enumerate}
  \item Read the introduction and Sections I-IV of Fama (1975), before doing this exercise. In the data file MISHKIN.ASC, monthly data are provided on:
\end{enumerate}

Column 1: year\\
Column 2: month

Column 3: one-month inflation rate (in percent, annual rate; call this PAII)\\
Column 4: three-month inflation rate (in percent, annual rate; call this PAI3)\\
Column 5: one-month T-bill rate (in percent, annual rate; call this TB1)\\
Column 6: three-month T-bill rate (in percent, annual rate; call this TB3)\\
Column 7: CPI for urban consumers, all items (the 1982-1984 average is set to 100 ; call this $C P I$ ).

The sample period is February 1950 to December 1990 ( 491 observations). The data on PAII, PAI3, TB1, and TB3 are the same data used in Mishkin (1992) and were made available to us by him. The T-bill data were obtained from the Center for Research in Security Prices (CRSP) at the University of Chicago. The T-bill rates for the month are as of the last business day of the previous month (and so can be taken for the interest rates at the beginning of the month). The construction of PAII and PAI3 will be described toward the end of this exercise; for the time being, we will use the inflation derived from $C P I$.\\
(a) (Library/internet work) To check the accuracy of the data in MISHKIN.ASC, find relevant tables from back issues of Treasury Bulletin or Federal Reserve Bulletin for hard-copy data on T-bill rates. Or visit the web sites of the Board of Governors (\href{http://www.bog.frb.fed.us}{www.bog.frb.fed.us}) or the Treasury Department (\href{http://www.ustreas.gov}{www.ustreas.gov}) to accomplish the same. Can you find one-month T-bill rates? [Answer: Probably not.] Are the rates in MISHKIN.ASC close to those in the relevant tables? Do they appear to be at the beginning of the month?\\
(b) (Library/internet work) Find relevant tables from back issues of Monthly Labor Review (or visit the web site of the Bureau of Labor Statistics, \href{http://www.bls.gov}{www.bls.gov}) to verify that the CPI figures in MISHKIN.ASC are correct. Verify that the timing of the variable is such that a January CPI observation is the CPI for the month. Regarding the definition of the CPI, verify the following. (1) The CPI is for urban consumers, for all items including food and housing, and is not seasonally adjusted. (2) Prices of the component items of the index are sampled throughout the month. When is the CPI for the month announced?\\
(c) Is the CPI a fixed-weight index or a variable-weight index? Hint: Dig up your old intermediate macro textbook; graduate macro textbooks won't do.

The one-month T-bill rate for month $t$ in the data set is for the period from the beginning of month $t$ to the end of the month (as you just verified). Ideally, if we had data on the price level at the beginning of each period, we would\\
calculate the inflation rate for the same period as $\left(P_{t+1}-P_{t}\right) / P_{t}$, where $P_{t}$ is the beginning-of-the-period price level. We use CPI for month $t-1$ for $P_{t}$ (i.e., set $P_{t}=C P I_{t-1}$ ). Since the CPI component items are collected at different times during the month, there arises the inevitable misalignment of the inflation measure and the interest-rate data. Another problem is the timing of the release of the CPI figure. The efficient market theory assumes that $P_{t}$ is known to the market at the beginning of month $t$ when the T-bill rates for the month are set. However, the CPI for month $t-1$, which we take to be $P_{t}$, is not announced until sometime in the following month (month $t$ ). Thus we are assuming that people know the CPI for month $t-1$ at the beginning of month $t$.

TSP Tip: When reading in the data and calculating the inflation rate, you should exploit TSP's ability to handle calendar dates. The initial part of your TSP program might look like:

\begin{verbatim}
? The data are monthly, 1950:2 thru 1990:12
freq m;smpl 50:2 90:12;
? Read in the ASCII data
read(file='mishkin.asc') year month pail pai3
    tb1 tb3 cpi;
? Calculate inflation rate and the real rate
smpl 50:3 90:12;
pai=((cpi/cpi(-1))**12-1)*100;
r = tb1-pai;
\end{verbatim}

RATS Tip: Similarly, the initial part of your RATS program might look like:

\begin{verbatim}
* The data are monthly, 1950:2 thru 1990:12
cal 50 2 12;all 0 90:12
* Read in the ASCII data
open data mishkin.asc
data(org=obs) / year month pai1 pai3 tb1 tb3 cpi
* Calculate inflation rate and real rate
set pai = ((cpi(t)/cpi(t-1))**12-1)*100
set r = tbl(t)-pai(t)
\end{verbatim}

(d) Reproduce the results in Table 2.1. Because the T-bill rate is in percent and at an annual rate, the inflation rate must be measured in the same unit. Calculate $\pi_{t+1}$, which is to be matched with $T B 1_{t}$ (the one-month T-bill rate for month $t$ ), as the continuously compounded rate in percent:

$$
\left[\left(\frac{P_{t+1}}{P_{t}}\right)^{12}-1\right] \times 100 .
$$

Gauss Tip: One way to compute the sample autocovariances by formula (2.10.1) and the sample autocorrelations by (2.10.2) is the following. Let $z$ be the $n$-dimensional vector whose $i$-th component is $z_{i}$ and nobs be the sample size, $n$.

\begin{verbatim}
z=z-meanc(z); /* de-mean the series */
nlag=12; /* lag length */
rho=zeros(nlag,1);
j=1;do until j>nlag;
    rho[j]=z'*shiftr(z',j,0)'/nobs;
    /* sample autocovariances */
j=j+1;endo;
rho=rho./(z'*z/nobs);
    /* sample autocorrelations */
\end{verbatim}

TSP Tip: Use the BJIDENT command. It does not calculate $p$-values, so (unless you are willing to learn how to use TSP's matrix commands) give up the idea of producing $p$-values.

RATS Tip: Use CORRELATE or BOXJENK.\\
(e) Can you reproduce (2.11.9), which has robust standard errors? What is the interpretation of the intercept?

Gauss Tip: For the heteroskedasticity-robust standard errors, you have to calculate $\widehat{\mathbf{S}}$ by formula (2.5.1) with $\hat{\varepsilon}_{i}=e_{i}$ (OLS residual). An example of the Gauss codes is the following. Let x be the $n \times K$ data matrix and ehat the $n$-dimensional vector of OLS residuals. If $g$ is the $n \times K$ matrix whose $i$-th row is $\mathbf{x}_{i} \cdot e_{i}$, the Gauss code for generating g is $\mathrm{g}=\mathrm{x}$.*ehat. The Gauss code for calculating $\widehat{\mathbf{S}}$ is $\mathrm{g}^{\prime} \mathrm{g} /$ nobs, where nobs is the sample size. This, however, does not exploit the fact that $\widehat{\mathbf{S}}$ is symmetric. The more computationally efficient command is moment ( $g, 0$ ) /nobs.

TSP Tip: Use the option $\mathrm{HCTYPE}=0$ of OLSQ .\\
RATS Tip: Use the ROBUSTERRORS option of LINREG.\\
(f) (Davidson-MacKinnon correction of White) The finite-sample properties of the robust $t$-ratio might be improved by the Davidson-MacKinnon adjust-\\
ment discussed in the text. Calculate robust standard errors using the three formulas discussed in the text. The first is the degrees of freedom correction of multiplying $\widehat{\mathbf{S}}$ by $n /(n-K)$. The second is to calculate $\widehat{\mathbf{S}}$ by formula (2.5.5) with $d=1$. The third is (2.5.5) with $d=2$. These correspond to TSP's OLSQ options HCTYPE $=1,2,3$.

Gauss Tip: Let p be the $n$-dimensional vector to store the $p_{i}$ 's of (2.5.5), sxxinv $=\mathbf{S}_{\mathbf{x x}}^{-1}, x=$ the data matrix $\mathbf{X}$. One way to compute $p$ is

\begin{verbatim}
p=zeros(nobs,1);
i=1;do until i>nobs;
    p[i]=x[i,.]*sxxinv*x[i,.]'/nobs;
i=i+1;endo;
\end{verbatim}

(g) (Estimation under conditional homoskedasticity) Test market efficiency by regressing $\pi_{t+1}$ on a constant and TB1, under conditional homoskedasticity (2.6.1). Compare your results with those in (e) and (f). Which part is different?\\
(h) (Breusch-Godfrey test) For the specification in (g), conduct the BreuschGodfrey test for serial correlation with $p=12$. (The $n R^{2}$ statistic should be about 27.0.) Let $e_{t}(t=1,2, \ldots, n)$ be the OLS residual from the regression in (g). To perform the Breusch-Godfrey test as described in the text, we need to set $e_{t}(t=0,1, \ldots,-11)$ to zero in running the auxiliary regression for $t=1,2, \ldots, n$.

TSP Tip: The TSP codes for calculating the $n R^{2}$ statistic are the following. resid is the OLS residual from the original regression.

\begin{verbatim}
? Part (h): Breusch-Godfrey
resid=@res;
smpl 52:1 52:12;
resid=0;
smpl 53:1 71:7;
olsq resid c tb1 resid(-1) resid(-2) resid(-3)
    resid(-4) resid(-5) resid(-6)
    resid(-7) resid(-8) resid(-9)
    resid(-10) resid(-11)
    resid(-12);
set w=@nob*@rsq;cdf(chisq,df=12) w;
\end{verbatim}

Here, $w$ is the $n R^{2}$ statistic.

RATS Tip: The RATS codes for calculating the $n R^{2}$ statistic are the following. resid is the OLS residual from the original regression.

\begin{verbatim}
* Part (h): Breusch-Godfrey
set resid 53:1 71:7 = pai-%beta(1)-%beta(2)*tb1
set resid 52:1 52:12 = 0
linreg resid 53:1 71:7;
    # constant tb1 resid{1 to 12}
compute w = %nobs*%rsquared
cdf chisquared w 12
\end{verbatim}

Here, w is the $n R^{2}$ statistic.\\
Gauss Tip: As in the Gauss codes for (d), the shiftr command is useful here. Let $y(223 \times 1)$ and $\times(223 \times 2)$ be the vector of the dependent variable and the matrix of regressors, respectively, from the previous regression. Let $\mathrm{b}(2 \times 1)$ be the vector of estimated coefficients from the previous regression. Your Gauss program for creating the y and x for the auxiliary regression might look like:

\begin{verbatim}
"@ Part (h): Breusch-Godfrey @";
y=y-x*b; /* residual vector from the previous
    regression */
j=1;do until j>12;
    x=x~shiftr(y',j,0)';
j=j+1;endo;
\end{verbatim}

(i) (Reconstructing Fama) What are Fama's own point estimate and standard error of the nominal interest rate coefficient? Are they identical to your results? (Fama uses different notation, so you need to translate his results in our notation.) Why is his estimate of the intercept different from yours? (One obvious reason is that our data are different, but there is another reason.) Optional: What are the differences between his data and our data?\\
(j) (Seasonal dummies) The CPI used to calculate the inflation rate is not seasonally adjusted. To take account of seasonal factors while still using seasonally unadjusted data in the regression, define twelve monthly dummies, $M 1, M 2, \ldots, M 12$, and use them in place of a constant in the regression in (g). Does this make any difference to your results? What happens if you include a constant along with the twelve monthly dummies in the regression? Optional: Is there any reason for preferring seasonally unadjusted data to seasonally adjusted?

TSP Tip: Use the DUMMY command to create monthly dummies.\\
RATS Tip: Use the SEASONAL command to create monthly dummies.\\
A well-known problem with the CPI series is that its residential housing component is home mortgage interest costs for periods before 1983 and a more appropriate "rental equivalence" measure since 1983. The inflation variables PAII and PAI3 in MISHKIN.ASC are calculated from a price index that uses the rental-equivalence measure for all periods. For more details on this, see Section II of Huizinga and Mishkin (1984). The timing of the variables is such that a January observation for PAII is calculated from the December and January data on the price index and $P A I 3$ from the December and March data on the price index. So, under the assumption that the price index is for the end of the month, a January observation for PAII is the inflation rate during the month, and a January observation for PAI3 is the inflation rate from the beginning of January to the end of March.\\
(k) Estimate the Fama regression for 1/53-7/71 using this better measure of the one-month inflation rate and compare the results to those you obtained in (e).\\
(I) Estimate the Fama regression for the post-October 1979 period. Is the nominal interest rate coefficient much lower? (The coefficient should drop to 0.564.)\\
2. (Continuation of the Nerlove exercise of Chapter 1, p. 76)\\
(i) For Model 4, carry out White's $n R^{2}$ test for conditional heteroskedasticity. ( $n R^{2}$ should be 66.546).\\
(j) (optional) Estimate Model 4 by OLS. This is Step 1 of the procedure of Section 2.8. Carry out Step 2 by estimating the following auxiliary regression: regress $e_{i}^{2}$ on a constant and $1 / Q_{i}$. Verify that Step 3 is what you did in part (h).\\
(k) Calculate White standard errors for Model 4.

\section*{MONTE CARLO EXERCISES}
\begin{enumerate}
  \item (Degrees of freedom correction) In the model of the Monte Carlo exercise of Chapter 1, we assumed that the error term was normal. Assume instead that $\varepsilon_{i}$ is uniformly distributed between -0.5 and 0.5 .\\
(a) Verify that the DGP of the second simulation (where $\mathbf{X}$ differs across simulated samples) satisfies Assumptions 2.1-2.5 and 2.7. (It can be shown that
\end{enumerate}

Assumption 2.2 [ergodic stationarity] is also satisfied. This is an implication of Proposition 6.1(d).)\\
(b) Run the second simulation of the Chapter 1 Monte Carlo exercise with the uniformly distributed error term. In each replication, calculate the usual $t$-ratio, as before, and compare it with two critical values. The first is the same critical value (2.042) implied by the $t(30)$ distribution, and the second is the critical value (1.96) from $N(0,1)$. Compute the rejection frequency for each critical value. Which one is closer to the nominal size of $5 \%$ ?\\
2. (Box-Pierce vs. Ljung-Box) We wish to verify the claim that the small sample properties of the Box-Ljung $Q$ statistic are superior to those of the Box-Pierce $Q$. Generate a string of 50 i.i.d. random numbers with mean 0 . (Choose your favorite distribution.) Taking this string as data, do the following.\\
(1) Calculate the Box-Pierce $Q$ and the Ljung-Box $Q$ statistics (see (2.10.4) and (2.10.5)) with $p=4$. (The ensemble mean is 0 by construction. Nevertheless take the sample mean from the series when you compute the autocorrelations.)\\
(2) For each statistic, accept or reject the null of no serial correlation at a significance level of $5 \%$, assuming that the statistic is distributed as $\chi^{2}(4)$.

Do this a large number of times, each time using a different string of 50 i.i.d. random numbers generated afresh. For each $Q$ statistic, calculate the frequency of rejecting the null. If the finite-sample distribution of the statistic is well approximated by $\chi^{2}(4)$, the frequency should be close to $5 \%$. Which statistic gives you the frequency closer to the nominal size of $5 \%$ ? Do we fail to reject the null too often if we use the Box-Pierce $Q$ ?

\section*{ANSWERS TO SELECTED QUESTIONS}
\section*{ANALYTICAL EXERCISES}
\begin{enumerate}
  \setcounter{enumi}{2}
  \item From (2.3.7) of the text, $\mathbf{b}-\boldsymbol{\beta}=\mathbf{S}_{\mathbf{x x}}^{-1} \overline{\mathbf{g}}$. By Kolmogorov, $\mathbf{S}_{\mathbf{x x}} \rightarrow_{\mathrm{p}} \mathrm{E}\left(\mathbf{x}_{i} \mathbf{x}_{i}^{\prime}\right)$ and $\overline{\mathbf{g}} \rightarrow_{\mathrm{p}} \mathbf{0}$. So $\mathbf{b}$ is consistent. By the Lindeberg-Levy CLT, $\sqrt{n} \overline{\mathbf{g}} \rightarrow_{\mathrm{d}} N(\mathbf{0}, \mathbf{S})$. The rest of the proof should now be routine.
  \item What remains to be shown is that the LHS of (**) equals Avar(b) $=\boldsymbol{\Sigma}_{\mathbf{x x}}^{-1} \mathbf{S} \boldsymbol{\Sigma}_{\mathbf{x x}}^{-1}$.
\end{enumerate}

$$
\boldsymbol{\Sigma}_{\mathbf{x x}}=\operatorname{plim} \frac{1}{n} \mathbf{X}^{\prime} \mathbf{X}
$$

and

$$
\begin{aligned}
\mathbf{S} & =\mathrm{E}\left(\varepsilon_{i}^{2} \mathbf{x}_{i} \mathbf{x}_{i}^{\prime}\right) \\
& =\mathrm{E}\left[\mathrm{E}\left(\varepsilon_{i}^{2} \mid \mathbf{x}_{i}\right) \mathbf{x}_{i} \mathbf{x}_{i}^{\prime}\right] \\
& =\mathrm{E}\left[\left(\mathbf{z}_{i}^{\prime} \boldsymbol{\alpha}\right) \mathbf{x}_{i} \mathbf{x}_{i}^{\prime}\right] \quad \text { (by (2.8.1)) } \\
& =\operatorname{plim} \frac{1}{n} \sum_{i=1}^{n}\left(\mathbf{z}_{i}^{\prime} \boldsymbol{\alpha}\right) \mathbf{x}_{i} \mathbf{x}_{i}^{\prime} \quad \text { (by ergodic stationarity) } \\
& \left.=\operatorname{plim} \frac{1}{n} \mathbf{X}^{\prime} \mathbf{V} \mathbf{X} \quad \text { (by definition (2.8.4) of } \mathbf{V}\right) .
\end{aligned}
$$

\section*{EMPIRICAL EXERCISES}
1c. The CPI is probably the most widely used fixed-weight price index.\\
1e. The negative of the intercept $=r$ (the constant ex-ante real interest rate).\\
19. Point estimates are the same. Only standard errors are different.

1i. Looking at Fama's Table 3, his $R_{t}$ coefficient when the sample period is $1 / 53$ to $7 / 71$ is 0.98 with a s.e. of 0.10 . Very close, but not exactly the same. There are two possible explanations for the difference between our estimates and Fama's. First, in calculating the inflation rate for month $t, \pi_{t+1}=\left(P_{t+1}-P_{t}\right) / P_{t}$, it is not clear from Fama that $C P I_{t-1}$ rather than $C P I_{t}$ was used for $P_{t}$. Second, the weight for the CPI at the time of his writing may be for 1958. The weight for the CPI in our data is for 1982-1984. Our estimate of the intercept ( -0.868 ) differs from Fama's ( 0.00070 ) because, first, his dependent variable is the negative of the inflation rate and, second, the inflation rate and the T-bill rate are monthly rates in Fama.

1j. The seasonally adjusted series manufactured by the BLS is a sort of two-sided moving average. Thus, for example, seasonally adjusted $C P I_{t}$ depends on seasonally unadjusted values of future $C P I$, which is not in $I_{1}$. Thus if there is no reason to take account of seasonal factors in the relationship between the inflation rate and the nominal interest rate, one should use seasonally unadjusted data. If there is a need to take account of seasonal factors, it is better to include seasonal dummies in the regression rather than use seasonally adjusted data. Inclusion of seasonal dummies does not change results in any important way.

1k. The $R_{t}$ coefficient drops to 0.807 from 1.014 .

\section*{References}
Amemiya, T., 1977, "A Note on a Heteroscedastic Model," Journal of Econometrics, 6, 365-370.\\
$\_\_\_\_$ , 1985, Advanced Econometrics. Cambridge: Harvard University Press.\\
Anderson, T. W., 1971, The Statistical Analysis of Time Series, New York: Wiley.\\
Billingsley, P., 1961, "The Lindeberg-Levy Theorem for Martingales," in Proceedings of the American Mathematical Society, Volume 12, 788-792.\\
Box, G., and D. Pierce, 1970, "Distribution of Residual Autocorrelations in AutoregressiveIntegrated Moving Average Time Series Models," Journal of the American Statistical Association, 65, 1509-1526.\\
Breusch, T., 1978. "Testing for Autocorrelation in Dynamic Linear Models," Australian Economic Papers, 17, 334-355.\\
Davidson, R., and J. MacKinnon, 1993, Estimation and Inference in Econometrics, Oxford: Oxford University Press.\\
Durbin, J., 1970, "Testing for Serial Correlation in Least-Squares Regression When Some of the Regressors are Lagged Dependent Variables," Econometrica, 38, 410-421.\\
Engle, R., 1982, "Autoregressive Conditional Heteroscedasticity with Estimates of the Variance of United Kingdom Inflation," Econometrica, 50, 987-1007.\\
Fama, E., 1975, "Short-Term Interest Rates as Predictors of Inflation," American Economic Review, 65, 269-282.\\
$\_\_\_\_$ , 1976, Foundations of Finance, Basic Books.\\
Godfrey, L., 1978, "Testing against General Autoregressive and Moving Average Error Models When the Regressors Include Lagged Dependent Variables," Econometrica, 46, 1293-1301.\\
Greene, W., 1997, Econometric Analysis (3d ed.), New York: Prentice Hall.\\
Gregory, A., and M. Veall, 1985, "Formulating Wald Tests of Nonlinear Restrictions," Econometrica, 53, 1465-1468.\\
Hall, P., and C. Heyde, 1980, Martingale Limit Theory and Its Application, New York: Academic.\\
Hall, R., 1978, "Stochastic Implications of the Life Cycle-Permanent Income Hypothesis: Theory and Evidence," Journal of Political Economy, 86, 971-987.\\
Hamilton, J., 1994, Time Series Analysis, Princeton: Princeton University Press.\\
Huizinga, J., and F. Mishkin, 1984, "Inflation and Real Interest Rates on Assets with Different Risk Characteristics," Journal of Finance, 39, 699-714.\\
Judge, G., W. Griffiths, R. Hill, H. Lütkepohl, and T. Lee, 1985, The Theory and Practice of Econometrics (2nd ed.), New York: Wiley.\\
Karlin, S., and H. M. Taylor, 1975, A First Course in Stochastic Processes (2d ed.), New York: Academic.\\
Koopmans, T., and W. Hood, 1953, "The Estimation of Simultaneous Linear Economic Relationships," in W. Hood and T. Koopmans (eds.), Studies in Econometric Method, New Haven: Yale University Press.\\
Mishkin, F., 1992, "Is the Fisher Effect for Real?" Journal of Monetary Economics, 30, 195-215.\\
Newey, W., 1985. "Generalized Method of Moments Specification Testing," Journal of Econometrics, 29, 229-256.\\
Rao, C. R., 1973, Linear Satistical Inference and Its Applications (2nd ed.), New York: Wiley.\\
Stock, J., 1994, "Unit Roots, Structural Breaks and Trends," in R. Engle and D. McFadden (eds.), Handbook of Econometrics, Volume IV, New York: North Holland.\\
White, H., 1980, "A Heteroskedasticity-Consistent Covariance Matrix Estimator and a Direct Test for Heteroskedasticity," Econometrica, 48, 817-838.


\end{document}
